\documentclass[11pt,a4paper]{article}
\usepackage[T1]{fontenc}
\usepackage[utf8]{inputenc}
\usepackage{lmodern}
\usepackage[margin=25mm,headheight=15pt]{geometry}
\usepackage{amsmath,amssymb,amsthm,mathtools}
\usepackage{microtype,booktabs,longtable,array,enumitem}
\usepackage{xcolor,xurl,hyperref,fancyhdr}
\definecolor{ink}{HTML}{16324F}
\definecolor{accent}{HTML}{1A626B}
\hypersetup{colorlinks=true,linkcolor=ink,citecolor=accent,urlcolor=accent,
 pdfdisplaydoctitle=true,
 pdftitle={Positive Certificates and Primitive Arithmetic for Euler's Constant},
 pdfauthor={Keio--Shaitetsu research notes}}
\newtheorem{theorem}{Theorem}[section]
\newtheorem{lemma}[theorem]{Lemma}
\newtheorem{proposition}[theorem]{Proposition}
\newtheorem{corollary}[theorem]{Corollary}
\theoremstyle{definition}
\newtheorem{definition}[theorem]{Definition}
\theoremstyle{remark}
\newtheorem{remark}[theorem]{Remark}
\setlist[enumerate]{topsep=4pt,itemsep=3pt}
\setlist[itemize]{topsep=4pt,itemsep=3pt}
\setlength{\parindent}{0pt}
\setlength{\parskip}{5pt plus 1pt minus 1pt}
\setlength{\emergencystretch}{2em}
\allowdisplaybreaks[2]
\pagestyle{fancy}\fancyhf{}
\fancyhead[L]{\small\color{ink}Positive certificates and primitive arithmetic}
\fancyhead[R]{\small Euler's constant}\fancyfoot[C]{\thepage}
\renewcommand{\headrulewidth}{0.3pt}
\setcounter{tocdepth}{2}
\begin{document}
\begin{center}
{\LARGE\bfseries\color{ink}Positive Certificates and\\[4pt]
Primitive Arithmetic for Euler's Constant}\par
\vspace{8pt}
{\large Target realization, exact optimization,\\
and arithmetic obstructions}\par
\vspace{12pt}
{Keio--Shaitetsu research notes}\par
{31 October 2026}\par
\vspace{6pt}
{\small A consolidated research manuscript with complete proofs}
\end{center}
\begin{abstract}
We develop a unified framework for positive, logarithm-free approximation
to Euler's constant in which every estimate is compared with the actual
reduced denominator. Canonical Binet kernels yield rational brackets and
a least-denominator selection with an exact rationality dichotomy.
Farey geometry reduces the primitive-error minimum on a rational interval
to logarithmically many candidates. Global positive target realization,
effective polynomial-square representations, optimal Laguerre
integerization, prime-residue tests, and support compression provide
complementary analytic and arithmetic tools. We prove divergence for
entire canonical, fixed-polynomial, and dyadic interpolation classes.
A localized positive family has the exact arithmetic threshold
\(q_M=o(M)\), and a one-step cutoff change removes a specified adverse
prime factor from its reduced denominator. Signed cancellation gives
an exponentially small integerized upper bound but does not guarantee
nonvanishing. Additional synthesis results determine the sharp
asymptotic degree of positive dyadic certificates, extend the selection
dichotomy to general shrinking rational intervals, exhibit an unbounded
gap between least-denominator and primitive-optimal selection, and prove
an exponential height barrier for independent signed outputs.
All statements retain their explicit hypotheses. The resulting tools
do not establish the irrationality of Euler's constant.
\end{abstract}
\vspace{4pt}
\noindent\textbf{Keywords.} Euler's constant; rational approximation;
positive kernels; Farey fractions; Laguerre polynomials; integer linear forms.

\clearpage
\tableofcontents
\clearpage

\section{The objective and the ten result groups}
Throughout,
\[
 \gamma=\lim_{n\to\infty}(H_n-\log n),\qquad
 H_n=\sum_{j=1}^n\frac1j,\qquad H_0=0.
\]
For a reduced rational \(p/q\), \(q>0\), its \emph{primitive error} is
\[
 \mathcal E(p/q)=q\left|\gamma-\frac pq\right|=|q\gamma-p|.
\]
The coefficients are integers; the value itself is not asserted to be
an integer. If \(\gamma=a/b\) in lowest terms, then
\(b(q\gamma-p)=aq-bp\in\mathbb Z\). Consequently a proved infinite
sequence satisfying
\[
 0<|q_j\gamma-p_j|\longrightarrow0,\qquad p_j,q_j\in\mathbb Z,
\]
would establish irrationality. The five requirements are exact logarithm
cancellation and integer coefficients, nonvanishing, decay after actual
denominator reduction, an infinite family, and unconditional proofs of
these assertions for the same family.

This manuscript consolidates the ten result groups selected from three
research notes \cite{source-a,source-b,source-c}. Their constructions and
proofs are reproduced and connected here, so no separate manuscript or
supplement is needed to follow the argument. Standard external theorems
are stated when used and attributed in the bibliography. The order below
is a dependency order; the first column records the earlier research
priority ranking, not a claim of historical priority.

\small
\begin{longtable}{@{}p{0.06\textwidth}p{0.51\textwidth}p{0.35\textwidth}@{}}
\toprule Rank & Result group & Principal statements\\\midrule
\endfirsthead
\toprule Rank & Result group & Principal statements\\\midrule
\endhead
1 & Explicit least-denominator selection and its rationality dichotomy &
Theorems~\ref{next-cutoff-selection}, \ref{stage3-nested-brackets}.\\
2 & Exact finite primitive-error optimization and certificates &
Theorems~\ref{stage3-farey-logarithmic-sail}, \ref{stage3-farey-certified-optimizer},
\ref{stage3-farey-three-fractions}.\\
3 & Constructive positive target realization and complete reduction &
Theorems~\ref{next-cutoff-interpolation}, \ref{sea-analytic-positive-zeros}.\\
4 & Complete polynomial-square target image and effective degree &
Theorems~\ref{sea-square-exponential-lift}, \ref{sea-square-effective-degree}.\\
5 & Optimal Laguerre integerization and surviving-prime tests &
Theorems~\ref{next-square-normalization}, \ref{next-square-prime-block}.\\
6 & Primitive-cost obstructions for entire coefficient classes &
Theorems~\ref{next-cutoff-cost}, \ref{next-square-polynomial-filter},
\ref{stage3-arithmetic-dyadic-obstruction}.\\
7 & Support compression with exact denominator and error transfer &
Theorem~\ref{sea-arith-cutoff-compression}.\\
8 & Positive rational localization and the exact threshold &
Theorems~\ref{stage3-localization-sandwich}, \ref{stage3-localization-asymptotic}.\\
9 & Removal of an adverse prime factor by changing the cutoff &
Theorem~\ref{stage3-localization-primes}.\\
10 & Integerized signed cancellation and its nonvanishing obstruction &
Theorem~\ref{stage3-counting-main}.\\
\bottomrule
\end{longtable}
\normalsize

\paragraph{What is additionally proved in this consolidation.}
The synthesis sections prove a sharp degree law for globally nonnegative
dyadic kernels, a general interval-selection dichotomy with a Fibonacci
counterexample to asymptotic optimality of least-denominator selection,
and a lower bound for the height of two independent outputs in the
actual signed-cancellation lattice. These conclusions are proved here
without assuming the irrationality of \(\gamma\). Their relation to
classical polynomial and Diophantine tools is made explicit; no literature
priority or external peer-review status is asserted.

\paragraph{Conventions.}
All logarithms are natural unless a base is shown. The notation
\(v_p\) is the usual rational \(p\)-adic valuation, and
\(\operatorname{odd}(n)\) is the largest odd divisor of a nonzero integer
\(n\), with its sign discarded. Empty products are one. A polynomial's
coefficient height always refers to the explicitly stated basis.
Different uses of \(A,B\) and of \(q\) are locally specified: the
polynomial-square forms \(A(f),B(f)\) are distinct from canonical
coefficient normalizers, and a polynomial variable is distinguished
from a reduced denominator by its argument.

\section{Canonical compressed families with exact positive kernels}
\label{next-cutoff-section}

This section analyzes a complete low-degree family inside the shared-cutoff
construction. No irrationality claim is made. The new construction gives
positive, strictly monotone rational approximants in two parameters; its
primitive cost can also be settled uniformly.

Write
\[
 d\mu(u)=\frac{du}{e^{2\pi\sqrt u}-1},\qquad
 \mu_j=\int_0^\infty u^j\,d\mu(u)
       =\frac{|B_{2j+2}|}{2j+2}.
\]
For an integer $n\ge1$, Binet's identity \cite{dlmf-binet} is
\[
 H_n-\frac1{2n}=\gamma+\log n-\Phi(n^2),
 \qquad \Phi(z)=\int_0^\infty\frac{d\mu(u)}{u+z}.
\]
An elementary moment estimate used below is
\begin{equation}\label{next-cutoff-moment-ratio}
 \frac{\mu_{j+1}}{\mu_j}
 =\frac{(2j+3)(2j+2)}{4\pi^2}
   \frac{\zeta(2j+4)}{\zeta(2j+2)}
 <\frac{(j+1)^2}{4}.
\end{equation}
Indeed the zeta ratio is less than $1$, $\pi>3$, and
$(2j+3)(2j+2)<9(j+1)^2$.

\begin{theorem}[An exact positive and doubly monotone canonical family]
\label{next-cutoff-canonical}
For integers $m,t\ge0$, put
\[
 P_m(x)=\prod_{j=1}^{m}(4^jx-1),\quad
 A_m=P_m(1),\quad s_m=\frac{P_m'(1)}{P_m(1)},\quad
 a=2^t,\quad b=t+s_m,\quad R=4^{m+1}.
\]
Empty products have value $1$, so $s_0=0$. Define
\[
 C_{m,t}(x)=\frac{P_m(x)}{A_m}\bigl[1-b(x-1)\bigr]
           =\sum_i c_i x^i,
 \qquad
 r_{m,t}=\sum_i c_i\left(H_{a2^i}-\frac1{2a2^i}\right).
\]
Then $r_{m,t}$ is rational and
\begin{equation}\label{next-cutoff-positive-kernel}
 E_{m,t}:=\gamma-r_{m,t}
 =\int_0^\infty
 \frac{u^m\{u+a^2[R+b(R-1)]\}}
      {\prod_{i=0}^{m+1}(u+a^24^i)}\,d\mu(u)>0.
\end{equation}
Both coordinate increases are strict:
\[
 r_{m+1,t}>r_{m,t},\qquad r_{m,t+1}>r_{m,t}.
\]
Moreover, uniformly in both parameters,
\begin{equation}\label{next-cutoff-two-sided}
 \frac{a^{-2m-2}4^{-m(m+1)/2}}{3^{15}}
 \le E_{m,t}
 \le \mu_m(t+m+3)a^{-2m-2}4^{-m(m+1)/2}.
\end{equation}
In particular $r_{m,t}\to\gamma$ along every sequence $m+t\to\infty$.
\end{theorem}

\begin{proof}
The defining polynomial satisfies
\[
 C_{m,t}(1)=1,\quad C_{m,t}'(1)=-t,\quad
 C_{m,t}(4^{-j})=0\quad(1\le j\le m).
\]
Binet's identity therefore cancels all logarithms exactly. To evaluate the
remaining kernel, first observe the rational identity
\begin{equation}\label{next-cutoff-base-kernel}
 \sum_{i=0}^m\frac{[x^i]P_m(x)}{u+a^24^i}
 =\frac{A_m u^m}{\prod_{i=0}^m(u+a^24^i)}.
\end{equation}
After taking a common denominator, the numerator has degree at most $m$.
Its first $m$ Taylor coefficients at zero vanish because
$P_m(4^{-j})=0$ for $1\le j\le m$; its leading coefficient is
$P_m(1)=A_m$. This proves the identity. Multiplication by $x$ shifts every
index by one. Applying this observation to $C_{m,t}$ gives
\eqref{next-cutoff-positive-kernel}, whose integrand is positive for $u>0$.

For monotonicity, introduce the direction
\[
 J_{m,t}:=\int_0^\infty
 \sum_i\frac{[x^i]\{P_m(x)(x-1)^2/A_m\}}{u+a^24^i}\,d\mu(u).
\]
Two shifts of \eqref{next-cutoff-base-kernel} give exactly
\begin{equation}\label{next-cutoff-direction}
 J_{m,t}
 =a^2(R-1)\int_0^\infty
 \frac{u^m\{4a^2(R-1)-3u\}}
      {\prod_{i=0}^{m+2}(u+a^24^i)}\,d\mu(u).
\end{equation}
This integral is strictly positive. Indeed, for the positive measure
$u^m\,d\mu$ and the strictly decreasing positive weight
$w(u)=\prod_{i=0}^{m+2}(u+a^24^i)^{-1}$, the weighted mean of $u$ is at
most $\mu_{m+1}/\mu_m$. This follows by integrating
$(u-v)(w(u)-w(v))\le0$ against the product measure. By
\eqref{next-cutoff-moment-ratio}, that mean is less than $(m+1)^2/4$,
whereas $4a^2(R-1)/3$ is larger. For instance $(m+1)^2\le4^m$ and
$4(R-1)/3\ge R$ suffice.

Direct polynomial subtraction now gives
\[
 C_{m+1,t}-C_{m,t}
 =-\frac{R[R+b(R-1)]}{(R-1)^2}
        \frac{P_m}{A_m}(x-1)^2,
\]
\[
 xC_{m,t+1}-C_{m,t}
 =-(b+1)\frac{P_m}{A_m}(x-1)^2.
\]
In the second identity the common external base is $a=2^t$; the leading
$x$ accounts for the new base $2a$. Both differences of gamma errors
are negative multiples of $J_{m,t}>0$, proving strict monotonicity.

For the upper bound in \eqref{next-cutoff-two-sided}, bound each
denominator factor below by $a^24^i$. This gives
\[
 E_{m,t}\le
 \frac{\mu_{m+1}+a^2[R+b(R-1)]\mu_m}
 {a^{2m+4}4^{(m+1)(m+2)/2}}.
\]
Since
$s_m=m+\sum_{j=1}^m(4^j-1)^{-1}<m+1$,
and \eqref{next-cutoff-moment-ratio} implies
$\mu_{m+1}/(a^2R\mu_m)<1$, this proves the stated upper bound.

For the lower bound, integrate only over $1\le u\le2$.
On that interval $d\mu(u)>3^{-12}du$, because
$2\pi\sqrt u<12$ and $e<3$. Also
\[
 \prod_{i=0}^{m+1}(u+a^24^i)
 \le a^{2m+4}4^{(m+1)(m+2)/2}
      \exp\!\left(\sum_{i\ge0}\frac2{4^i}\right)
 <27a^{2m+4}4^{(m+1)(m+2)/2}.
\]
The numerator is at least $a^2R$ there. Combining these inequalities
proves the lower bound.

Finally,
\[
 \frac{\mu_{m+1}4^{-(m+1)(m+2)/2}}
      {\mu_m4^{-m(m+1)/2}}
 <\frac{(m+1)^2}{4^{m+2}}\le\frac1{16}.
\]
Thus the upper bound is at most
\[
 \frac{t+m+3}{12}\,4^{-t(m+1)}16^{-m}
 \le\frac{t+m+3}{12}\,4^{-t-m},
\]
which tends to zero uniformly as $t+m\to\infty$.
\end{proof}

\begin{theorem}[Complete compressed classification and raw-error optimum]
\label{next-cutoff-classification}
Fix $m,t$ and use the preceding notation. Every rational polynomial $C$
of degree at most $m+2$ satisfying
\[
 C(1)=1,\qquad C'(1)=-t,\qquad
 C(4^{-j})=0\quad(1\le j\le m)
\]
has the unique form
\[
 C_\kappa(x)=\frac{P_m(x)}{A_m}
       \bigl[1-b(x-1)+\kappa(x-1)^2\bigr],
 \qquad \kappa\in\mathbb Q.
\]
This includes every remainder of degree less than $m+3$ obtained from
the shared-cutoff polynomial
$(x-1)^2(2x-1)P_m(x)$, whenever these normalized constraints hold.

Writing $d=R-1$, its exact gamma-error kernel is
\[
 F_\kappa(u)=
 \frac{u^m(u^2+a^2A_\kappa u+a^4B_\kappa)}
      {\prod_{i=0}^{m+2}(u+a^24^i)},
\]
where
\[
 A_\kappa=5R+d(b-3\kappa),\qquad
 B_\kappa=4\{R^2+bRd+\kappa d^2\}.
\]
Consequently $F_\kappa(u)\ge0$ for all $u\ge0$ if and only if
\begin{equation}\label{next-cutoff-quadratic-positivity}
 B_\kappa\ge0,\qquad
 A_\kappa\ge0\ \hbox{or}\ A_\kappa^2\le4B_\kappa.
\end{equation}
Among all these pointwise nonnegative kernels, the smallest integral is
attained uniquely at
\[
 \kappa_{\min}=-\frac{R(R+bd)}{d^2}.
\]
At this parameter $C_{\kappa_{\min}}=C_{m+1,t}$, so the minimum error is
exactly $E_{m+1,t}$. This is an optimum for the unnormalized integral;
it is not an optimum for the final reduced-denominator error.
\end{theorem}

\begin{proof}
Divisibility by $P_m$ and the two conditions at $1$ give the unique
quadratic factor. Three shifts of \eqref{next-cutoff-base-kernel} yield
the displayed numerator. A monic quadratic $v^2+Av+B$ is nonnegative
on $v\ge0$ exactly when $B\ge0$ and either $A\ge0$ or its real minimum
$B-A^2/4$ is nonnegative; this proves
\eqref{next-cutoff-quadratic-positivity}.

The condition $B_\kappa\ge0$ forces $\kappa\ge\kappa_{\min}$.
At $\kappa_{\min}$, $B_\kappa=0$ and $A_\kappa>0$, so this endpoint is
admissible. Its polynomial additionally vanishes at $4^{-m-1}$, hence
the already proved polynomial identity identifies it with $C_{m+1,t}$.
The error is the affine function
\[
 E(\kappa)=E_{m,t}+\kappa J_{m,t}.
\]
The preceding proof established $J_{m,t}>0$ for every $m,t$, so the
minimum over all admissible parameters is attained uniquely at the
least one, as claimed.
\end{proof}

\paragraph{A rational double-zero example.}
For $m=0$, $t=4$, and $\kappa=80/9$, one has
\[
 C_\kappa(x)=\frac{125-196x+80x^2}{9},\qquad a=16,
\]
and the exact positive kernel is
\[
 F_\kappa(u)=
 \frac{(u-6144)^2}
 {(u+256)(u+1024)(u+4096)}.
\]
Thus an interior double zero, exact logarithmic cancellation, and global
positivity coexist in a three-node formula. No quality claim about its
final rational approximation is inferred from that zero.

\begin{theorem}[An explicit positive certificate with a square-root logarithmic degree bound]
\label{next-cutoff-interpolation}
For every rational $p<\gamma$, an exact rational search constructs a
globally positive rational kernel of mass exactly $\gamma-p$.
When $1/2\le p<\gamma$, it constructs a dyadic polynomial $C$ with
\[
 C(1)=1,\qquad C'(1)=0,
\]
whose dyadic kernel is positive for all $u>0$ and whose corrected harmonic
target is exactly $p$. If a separation bound
$0<\varepsilon\le\gamma-p$ is available, any integer $M\ge32$ satisfying
\[
 2^{-M^2/2}\le\varepsilon
\]
gives the explicit degree bound $\deg C\le M+1$ and support bound
$n\le2^{M+1}$. In particular the degree bound is
$O(\sqrt{\log(1/\varepsilon)})$ as $\varepsilon\to0$.

For $p=P/Q$ in lowest terms, the complete primitive integrated form is
exactly $Q\gamma-P$, regardless of the intermediate interpolation
coefficient denominators. The construction does not select a sequence
for which these primitive forms tend to zero.
\end{theorem}

\begin{proof}
For $p<1/2$, use
\[
 \frac1{u+1}+12(1/2-p)K_1(u),\qquad
 K_1(u)=\frac1{u+1}-\frac1{u+4}
       +\frac1{u+9/4}-\frac1{u+9}.
\]
Each difference in $K_1$ is positive, and Binet's integral combined with
digamma duplication \cite{dlmf-binet,dlmf-duplication}
gives $\int K_1\,d\mu=1/12$, proving this assertion directly.

Write $r_m=r_{m,0}$. The preceding theorem gives $r_0=1/2$ and
$r_m\uparrow\gamma$, with every $r_m$ rational and exactly computable.
If $p=1/2$, take $C=1$. Otherwise enumerate these exact rationals until
the first $m\ge1$ with $r_m\ge p$. Define
\[
 \lambda=\frac{p-r_{m-1}}{r_m-r_{m-1}}\in(0,1],\qquad
 C=(1-\lambda)C_{m-1,0}+\lambda C_{m,0}.
\]
All coefficients are rational. The constraints and target are exact by
linearity, and the kernel is a convex combination of two strictly
positive kernels. The search terminates precisely for $p<\gamma$ in
the stipulated range $p\ge1/2$.

For the quantitative bound, the moment formula gives
$\mu_m<(2m+1)!$. Hence
\[
 E_{m,0}\le(m+3)(2m+1)!\,2^{-m(m+1)}.
\]
For $m\ge32$, the elementary inequalities
$2m+1\le2^{m/4}$ and $m+3\le2^{m/2}$ imply
\[
 (m+3)(2m+1)!\le
 2^{m/2}(2m+1)^{2m+1}
 \le2^{m^2/2+3m/4}.
\]
Consequently $E_{m,0}<2^{-m^2/2}$.
If $M$ is as stated, then $r_M>\gamma-\varepsilon\ge p$, so the search
terminates with $m\le M$. The degree and support bounds follow.

The interpolation also has an exact nonnegative integer form. If the
two bracketing targets have primitive pairs $(q_0,p_0),(q_1,p_1)$,
put $\Delta=q_0p_1-q_1p_0>0$. For $p=P/Q$ reduced inside their interval,
\[
 u=Qp_1-Pq_1\ge0,\qquad v=Pq_0-Qp_0\ge0,
\]
and
\[
 u(q_0,p_0)+v(q_1,p_1)=\Delta(Q,P).
\]
Thus a nonnegative integer combination of the two positive primitive
kernels gives exactly $\Delta(Q\gamma-P)$, whose full coefficient gcd
is $\Delta$. This explicitly tracks the cancellation of all factors of
the determinant, including primes already present in the base targets.

Finally choose an integer $D$ that clears the integrated coefficient
pair and is divisible by $Q$. The pair is $(D,DP/Q)$ and its gcd is
$D/Q$, so complete reduction leaves $(Q,P)$. This arithmetic identity
does not bound $Q(\gamma-p)$ for targets that have not yet been selected.
\end{proof}


% New synthesis result for the unified gamma paper.
% This fragment uses the ambient theorem environments and the already defined
% Euler constant gamma, harmonic numbers H_n, and Binet measure d\mu.
% It contains no cross-reference to an external manuscript.

\section{The sharp degree scale for positive dyadic certificates}
\label{unified-degree-section}

The exact-target constructions allow arbitrary rational coefficients.  Their
height may therefore grow without restriction.  Nevertheless, positivity
imposes a quantitative lower bound that depends only on the dyadic support.
Together with the canonical construction, this determines the leading
asymptotic constant in the smallest possible certificate degree.  This is
an analytic complexity statement; it does not give primitive smallness.

\begin{lemma}[A monic polynomial bound on the half-line]
\label{unified-degree-monic}
Let $P\in\mathbb R[u]$ be monic of degree $d\ge0$ and nonnegative on
$[0,\infty)$.  Then
\[
 \int_0^1 P(u)\,du\ \ge\
 \mathfrak c_d:=\frac{1}{(d+1)\binom{d}{\lfloor d/2\rfloor}^{\!2}}.
\]
The bound is attained among such polynomials.
\end{lemma}

\begin{proof}
First, every polynomial nonnegative on the half-line has a representation
\[
 P(u)=S(u)^2+uT(u)^2,\qquad S,T\in\mathbb R[u].
\]
Here is a direct factorization proof in the monic case.  Every positive
real root has even multiplicity, so its factor is a square.  A zero root
contributes powers of $u=0^2+u\cdot1^2$.  A negative root contributes
$u+a=(\sqrt a)^2+u\cdot1^2$, where $a>0$.  For a pair of nonreal
conjugate roots, put $r=\sqrt{b^2+c^2}$, with $c\ne0$, and use
\[
 (u-b)^2+c^2=(u-r)^2+u\bigl(\sqrt{2(r-b)}\bigr)^2.
\]
Products remain of the stated form because
\[
 (S^2+uT^2)(V^2+uW^2)
   =(SV-uTW)^2+u(SW+TV)^2.
\]
This proves the representation, including possible repeated roots.

Write $d=2k+\alpha$, where $\alpha\in\{0,1\}$.  If $\alpha=0$,
the degree and leading coefficient force $\deg S=k$, with leading
coefficient $1$ or $-1$, and $\deg T\le k-1$.  If $\alpha=1$,
they force $\deg T=k$, with leading coefficient $1$ or $-1$, and
$\deg S\le k$.  Thus it suffices to bound the squared norm of a monic
polynomial of degree $k$ for the weight $u^\alpha$ on $[0,1]$.

For completeness, the relevant shifted Jacobi polynomial is
\[
 J_{k,\alpha}(u)=
 \frac{1}{k!u^\alpha}\frac{d^k}{du^k}
       \bigl[u^{k+\alpha}(u-1)^k\bigr].
\]
The expression is a polynomial, with leading coefficient
$B=\binom{2k+\alpha}{k}$.  Integration by parts $k$ times shows
that it is orthogonal, for the weight $u^\alpha$, to every polynomial
of degree less than $k$.  All boundary terms vanish.  A second such
calculation, followed by the elementary beta integral, gives
\[
 \int_0^1 u^\alpha J_{k,\alpha}(u)^2\,du
 =B\int_0^1u^{k+\alpha}(1-u)^k\,du
 =\frac{1}{2k+\alpha+1}.
\]
Consequently, for every monic $V$ of degree $k$, orthogonality yields
\[
 \int_0^1u^\alpha V(u)^2\,du
 \ge \int_0^1u^\alpha\left(\frac{J_{k,\alpha}(u)}B\right)^2du
 =\frac{1}{(2k+\alpha+1)B^2}=\mathfrak c_d.
\]
Applying this to $S$ or $T$, after changing its sign if necessary,
proves the bound.  Equality is attained by
$P=(J_{k,0}/B)^2$ when $d$ is even, and by
$P=u(J_{k,1}/B)^2$ when $d$ is odd.  Both are nonnegative on the
entire half-line.
\end{proof}

\begin{theorem}[A uniform lower bound at a fixed dyadic cutoff]
\label{unified-degree-cutoff}
Let $a\ge1$, let $d\ge0$ be an integer, and let $c_0,\ldots,c_d$ be
real numbers satisfying $\sum_{i=0}^d c_i=1$.  Suppose that
\[
 \mathcal K(u)=\sum_{i=0}^d\frac{c_i}{u+a^24^i}\ge0
 \qquad(u\ge0).
\]
For the Binet measure $d\mu(u)=du/(e^{2\pi\sqrt u}-1)$, one has
\begin{align}
 \int_0^\infty\mathcal K(u)\,d\mu(u)
 &\ge
 \frac{\mathfrak c_d}{3^7\prod_{i=0}^d(1+a^24^i)}
 \label{unified-degree-product-bound}\\
 &\ge
 \frac{1}{3^8(d+1)a^{2d+2}2^{d^2+3d}}.
 \label{unified-degree-elementary-bound}
\end{align}
In particular, the bounds hold without any restriction on coefficient
height, coefficient signs, or rational denominators.
\end{theorem}

\begin{proof}
Put
\[
 D(u)=\prod_{i=0}^d(u+a^24^i),\qquad P(u)=D(u)\mathcal K(u).
\]
The numerator $P$ is a polynomial of degree exactly $d$, and its leading
coefficient is $\sum_i c_i=1$.  Since $D>0$ on the half-line, $P$ is
nonnegative there.  On $0<u\le1$, the elementary inequalities
$2\pi<7$ and $e<3$ give
\[
 \frac{1}{e^{2\pi\sqrt u}-1}>3^{-7},\qquad
 D(u)\le\prod_{i=0}^d(1+a^24^i).
\]
The preceding lemma, applied after restricting the integral to $[0,1]$,
proves \eqref{unified-degree-product-bound}.

Furthermore,
\[
 \prod_{i=0}^d(1+a^24^i)
 =a^{2d+2}2^{d(d+1)}
       \prod_{i=0}^d(1+a^{-2}4^{-i})
 <3a^{2d+2}2^{d(d+1)}.
\]
Indeed the last product is at most
$2\exp(\sum_{i\ge1}4^{-i})=2e^{1/3}<3$.
Finally $\binom d{\lfloor d/2\rfloor}\le2^d$, so
$\mathfrak c_d\ge1/((d+1)2^{2d})$.  These estimates give
\eqref{unified-degree-elementary-bound}.
\end{proof}

\begin{theorem}[Sharp asymptotic degree for exact positive dyadic targets]
\label{unified-degree-sharp}
For a rational number $p\in[1/2,\gamma)$, let $d_*(p)$ be the least
degree of a polynomial $C(x)=\sum_i c_i x^i\in\mathbb Q[x]$ satisfying
\[
 C(1)=1,\qquad C'(1)=0,\qquad
 \mathcal K_C(u):=\sum_i\frac{c_i}{u+4^i}\ge0\quad(u\ge0),
\]
and having the exact corrected harmonic target
\[
 p=\sum_i c_i\left(H_{2^i}-\frac1{2^{i+1}}\right).
\]
This minimum exists.  As rational $p$ approaches $\gamma$ from below,
\begin{equation}
 \frac{d_*(p)}{\sqrt{\log_2(1/(\gamma-p))}}\longrightarrow1.
 \label{unified-degree-sharp-limit}
\end{equation}
No assumption about the rationality or irrationality of $\gamma$ is
used.  If $p=P/Q$ is reduced, complete integer normalization still
leaves exactly $Q\gamma-P$; the degree conclusion is independent of
whether this primitive error is small.
\end{theorem}

\begin{proof}
Binet's identity gives
\[
 H_{2^i}-\frac1{2^{i+1}}
 =\gamma+i\log2-\int_0^\infty\frac{d\mu(u)}{u+4^i}.
\]
The two normalization conditions therefore imply
\[
 \gamma-p=\int_0^\infty\mathcal K_C(u)\,d\mu(u)>0.
\]
Put $\delta=\gamma-p$ and $L=\log_2(1/\delta)$.  The preceding
theorem with $a=1$ yields, for every admissible degree $d$,
\[
 L\le d^2+3d+\log_2(d+1)+8\log_2 3.
\]
In particular $d_*(p)\to\infty$ as $\delta\to0$, and this inequality
implies
\[
 \liminf_{p\uparrow\gamma}\frac{d_*(p)}{\sqrt L}\ge1.
\]

We recall the explicit canonical family to prove the matching upper bound.
Set
\[
 P_m(x)=\prod_{j=1}^m(4^jx-1),\quad
 A_m=P_m(1),\quad s_m=\frac{P_m'(1)}{P_m(1)},\quad
 C_m(x)=\frac{P_m(x)}{A_m}[1-s_m(x-1)].
\]
For $m=0$, take $P_0=A_0=C_0=1$ and $s_0=0$.  These polynomials
have $C_m(1)=1$, $C_m'(1)=0$, and degree at most $m+1$.
The partial-fraction identity
\[
 \sum_i\frac{[x^i]P_m(x)}{u+4^i}
   =\frac{A_m u^m}{\prod_{i=0}^m(u+4^i)}
\]
follows by comparing the numerator's leading coefficient and its first
$m$ Taylor coefficients at zero.  Multiplication by $x$ shifts the
indices by one; it follows that, with $R_m=4^{m+1}$,
\[
 \mathcal K_{C_m}(u)=
 \frac{u^m\{u+R_m+s_m(R_m-1)\}}
      {\prod_{i=0}^{m+1}(u+4^i)}>0\qquad(u>0).
\]
Let $r_m$ be its rational corrected target and put
$E_m=\gamma-r_m$.  In particular $r_0=1/2$.

The moment formula
\[
 \mu_j=\int_0^\infty u^j\,d\mu(u)
 =\frac{2(2j+1)!\,\zeta(2j+2)}{(2\pi)^{2j+2}}
\]
gives $\mu_j<(2j+1)!$ and
$\mu_{j+1}/\mu_j<(j+1)^2/4$.  The latter inequality follows from
$\zeta(2j+4)/\zeta(2j+2)<1$, $\pi>3$, and
$(2j+3)(2j+2)<9(j+1)^2$.
Also $s_m=m+\sum_{j=1}^m(4^j-1)^{-1}<m+1$.
Bounding each factor in the denominator of the positive kernel below
by $4^i$ now gives
\[
 E_m\le\mu_m(m+3)2^{-m(m+1)}
 <(m+3)(2m+1)!\,2^{-m(m+1)}.
\]
Consequently
\[
 \log_2 E_m\le -m^2+O(m\log(m+1)),\qquad E_m\longrightarrow0.
\]
The implied constant is absolute.

For $p>1/2$, enumerate the rational numbers $r_m$ until the first
index $j\ge1$ with $r_j\ge p$.  Such an index exists because
$E_m\to0$.  Minimality gives $r_{j-1}<p\le r_j$, and hence
\[
 \lambda=\frac{p-r_{j-1}}{r_j-r_{j-1}}\in(0,1],\qquad
 C=(1-\lambda)C_{j-1}+\lambda C_j
\]
is an admissible rational polynomial of degree at most $j+1$.
Its positivity and exact target follow from linearity and convexity.
For $p=1/2$, simply use $C=1$.  This proves existence of $d_*(p)$
by exact rational comparisons, without an equality oracle for $\gamma$.

Fix $\eta>0$ and take $m=\lceil(1+\eta)\sqrt L\rceil$.
As $L\to\infty$,
\[
 \log_2 E_m\le -(1+\eta)^2L+O(\sqrt L\log L)<-L.
\]
Thus $E_m<\delta$, so $r_m>p$, and the first crossing index above
satisfies $j\le m$.  Therefore
\[
 \limsup_{p\uparrow\gamma}\frac{d_*(p)}{\sqrt L}\le1+\eta.
\]
Letting $\eta\downarrow0$ proves \eqref{unified-degree-sharp-limit}.
Finally, an integer $D$ clearing the integrated rational pair gives
$(D,DP/Q)$ with gcd $D/Q$, leaving exactly $(Q,P)$.  This verifies
the last assertion without a bound on intermediate coefficient heights.
\end{proof}

\paragraph{A quantitative price for positive support compression.}
Suppose that an original normalized certificate has error $E>0$, and a
target-changing operation replaces it by a globally nonnegative kernel
supported at $a2^i$, $0\le i\le d$, with the same gamma coefficient $1$.
Assume the complete logarithm condition is preserved, so both masses are
log-free gamma errors.  If their targets are $p$ and $\widetilde p$, then
\[
 p-\widetilde p
 = (\gamma-\widetilde p)-(\gamma-p)
 \ge\frac{1}{3^8(d+1)a^{2d+2}2^{d^2+3d}}-E.
\]
In particular, if $E$ is at most half of this cutoff bound, any such
positive compression must lower the target by at least half of the
bound.  This applies to shared-cutoff compression whenever its normalized
output kernel remains globally nonnegative.  A certificate merely proving
that its integrated error is positive does not suffice for this conclusion.

The sharp degree law concerns a fixed base and a specified pole grid.
It does not bound arbitrary rational-node representations, support with a
varying external base, or signed kernels with only a positive integral.
Nor does it control the denominator $Q$ of a target or prove
$Q(\gamma-P/Q)\to0$.

\section{Global positive certificates with prescribed local data}
\label{sea-analytic}

The first three results join three requirements that were previously handled
separately: control on the whole half-line, exact local vanishing, and an exact
rational integral target. They do not provide a bound for the denominator of
a well-chosen target.

For $a>0$ and $u\geq0$, put
\begin{align}
h_a(u)&=\frac1{u+a^2},&
K_a(u)&=h_a(u)+h_{a+1/2}(u)-h_{2a}(u)-h_{2a+1}(u),\\
k_a(u)&=(1+u)K_a(u),&
d\mu(u)&=\frac{du}{e^{2\pi\sqrt{u}}-1},\qquad
\ell(f)=\int_0^\infty\frac{f(u)}{1+u}\,d\mu(u).
\end{align}
Binet's formula and digamma duplication \cite{dlmf-binet,dlmf-duplication}
give
\begin{equation}\label{sea-analytic-masses}
\int h_1\,d\mu=\gamma-\frac12,\qquad
\ell(k_a)=\int K_a\,d\mu=\frac1{4a(2a+1)},\qquad
0<\ell(1)<\frac1{12}.
\end{equation}
For completeness, Binet gives $I(a):=\int h_a\,d\mu
=\log a-\psi(a)-1/(2a)$. Substituting the duplication identity into
$I(a)+I(a+1/2)-I(2a)-I(2a+1)$ yields the rational mass in
\eqref{sea-analytic-masses}. The total mass of $\mu$ is $1/12$,
so $0<\int h_1\,d\mu<1/12$ follows directly. In particular all the
displayed kernel masses are rational when $a$ is rational.

\begin{theorem}[Density with global and derivative control]
\label{sea-analytic-density}
Let $M\geq0$ be an integer and $L>0$. The rational linear span of
$\{k_a:a\in\mathbb Q_{>0}\}$ is dense in the Banach space
\[
X_{M,L}=\{f\in C_0([0,\infty)):f|_{[0,L]}\in C^M([0,L])\},
\quad
\|f\|_X=\|f\|_\infty+\sum_{j=1}^M\|f^{(j)}\|_{[0,L]}.
\]
No derivative control at infinity is asserted.
\end{theorem}

\begin{proof}
Completeness follows from uniform convergence of the functions and their
local derivatives. Suppose a continuous real functional $\Lambda$ annihilates
all rational-node kernels. Hahn--Banach, applied to the embedding of $X$
in $C_0([0,\infty))\oplus C([0,L])^M$, and the measure representation of
the dual give finite signed measures $\nu_0,\ldots,\nu_M$ such that
\[
\Lambda(f)=\int f\,d\nu_0+
\sum_{j=1}^M\int_{[0,L]}f^{(j)}\,d\nu_j.
\]
Apply this integral expression to $g_a(u)=(1+u)/(u+a^2)$ and denote the
result by $F(a)$. The expression is meaningful although $g_a\notin C_0$.
For $a\geq1$, $|g_a|\leq1$ and $g_a(u)\to0$ at each finite $u$ as
$a\to\infty$. Moreover
\[
g_a(u)=1+\frac{1-a^2}{u+a^2}
\]
shows that all its positive-order derivatives tend uniformly to zero on
$[0,L]$. Thus $F(a)\to0$ by dominated convergence.

Continuity in $a$ extends the annihilation identity from rational nodes to
all $a>0$:
\[
F(a)+F(a+1/2)-F(2a)-F(2a+1)=0.
\]
Set $D(a)=F(a)-F(2a)$. Then $D(a+1/2)=-D(a)$, whereas
$D(a+j/2)\to0$. Hence $D=0$. Iterating $F(a)=F(2a)$ and using the
same limit proves $F=0$.

Write $x=(1+u)^{-1}$. Differentiation with respect to $z=a^2$ at $z=1$
now gives $\Lambda(x^n)=0$ for every $n\geq1$. Differentiation under
each measure is justified by uniform bounds in a fixed neighborhood of
$z=1$, including the required local derivatives.

Polynomials in $x$ with zero constant term are dense in $X$. Here are
details for this last approximation step. The change of variable identifies
$X$ with functions continuous on $[0,1]$, zero at $0$, and $C^M$ on
$[b,1]$, where $b=(1+L)^{-1}$; the local derivative norms are equivalent.
Extend the Taylor polynomial at $b$ to its left and blend it, in a small
interval strictly to the left of $b$, into a smooth uniform approximation
of the remaining continuous function. This gives a global $C^M$ function
that is unchanged on $[b,1]$ and arbitrarily close uniformly. A small smooth
bump below $b$ restores its value $0$ at $0$. Approximate its $M$th
derivative uniformly by a polynomial and integrate $M$ times with the
initial derivatives fixed. The resulting polynomial is close in $C^M$ and
is still zero at $0$. For $M=0$, ordinary polynomial approximation followed
by subtraction of the value at $0$ suffices. Therefore $\Lambda=0$.
This proves real density. Approximation of the finitely many coefficients
in each combination by rational numbers proves rational-span density.
\end{proof}

\begin{theorem}[Exact rational jets and mass with a bounded correction]
\label{sea-analytic-projection}
Let $T_1,\ldots,T_s$ consist of distinct derivative evaluations
$f\mapsto f^{(r)}(x)$, with $x\in[0,L]\cap\mathbb Q$ and $0\leq r\leq M$,
together with the single functional $\ell$. There are positive rational
nodes $a_1,\ldots,a_s$ for which $C_{ji}=T_j(k_{a_i})$ is an invertible
rational matrix. The nodes can be found by exact rational enumeration.
For $b\in\mathbb Q^s$, define
\[
\mathcal P_b(v)=v+\sum_{i=1}^s
\bigl[C^{-1}(b-T(v))\bigr]_i k_{a_i}.
\]
This map satisfies $T(\mathcal P_b(v))=b$ and preserves the rational
kernel span. If $T(F)=b$, then
\[
\|\mathcal P_b(v)-F\|_X\leq C_T\|v-F\|_X,
\qquad
C_T=1+\sum_{i,j}U_i|(C^{-1})_{ij}|V_j,
\]
where $U_i$ is any rational upper bound for $\|k_{a_i}\|_X$, and $V_j$
can be taken as $1$ for an evaluation and $1/12$ for the mass functional.
The upper bound $C_T$ is therefore computable in rational arithmetic.
\end{theorem}

\begin{proof}
The functionals are independent. A smooth positive bump away from the
finitely many evaluation points first eliminates the coefficient of $\ell$
in any relation. Localized polynomial bumps then separate the derivative
evaluations. Theorem~\ref{sea-analytic-density} implies full rank on the
kernel span. All matrix entries at rational nodes are rational, so
enumerating nodes and retaining independent columns terminates.
The exact constraint identity follows by multiplication by $C$. The error
bound follows from $b-T(v)=T(F-v)$ and the triangle inequality. Explicit
rational norm bounds follow by summing, over the four terms defining $k_a$,
\[
\|g_a\|_\infty\leq\max(1,a^{-2}),\qquad
\|g_a^{(r)}\|_{[0,L]}\leq
\frac{r!|1-a^2|}{a^{2r+2}}\quad(r\geq1).
\]
Finally $\|\ell\|\leq1/12$ follows from
\eqref{sea-analytic-masses}.
\end{proof}

\begin{theorem}[Prescribed zeros and an exact rational target]
\label{sea-analytic-positive-zeros}
Fix finitely many distinct nonnegative rational points $x_j$ and positive
integer multiplicities $m_j$, with $m_j$ even whenever $x_j>0$.
Let $d=\sum_jm_j$ and
\[
W(u)=\frac{\prod_j(u-x_j)^{m_j}}{(1+u)^d};
\]
use $W=1$ for an empty set. For rational $p$, the following are equivalent:
\begin{enumerate}[label=(\roman*)]
\item $p<\gamma$.
\item There exist rational $c_i$, positive rational $a_i$, and rational
$\eta>0$ such that
\[
R(u)=h_1(u)-\sum_i c_iK_{a_i}(u)
\geq\frac{\eta W(u)}{1+u}\qquad(u\geq0),
\quad
\int R\,d\mu=\gamma-p,
\]
and the zeros of $R$ on $[0,\infty)$ are exactly the specified points,
with exactly the specified multiplicities.
\end{enumerate}
Moreover $R(u)\sim1/u$ at infinity. If $p=P/Q$ is reduced, complete
primitive normalization of its integrated identity gives exactly
$Q\gamma-P$, independently of the intermediate coefficient denominators.
\end{theorem}

\begin{proof}
Suppose $\epsilon=\gamma-p>0$. There exists a smooth function $G>0$,
equal to $1$ sufficiently far out, such that $\ell(WG)=\epsilon$.
Indeed, let $G_{\rm low}$ be a small positive constant on a long initial
interval and $1$ in the far tail. Boundedness of $W$ and finiteness of
$h_1d\mu$ show that these masses can tend to $0$. On the other hand,
$G_{\rm high}=1+C\varphi$, for a nonnegative smooth bump $\varphi$
away from the zero set, gives arbitrarily large masses. A suitable convex
combination realizes $\epsilon$ exactly while remaining positive and
equal to $1$ in the tail.

Put $F_*=1-WG$. It belongs to $X_{M,L}$ for
$M=\max_jm_j$ and $L$ beyond all the points. At each $x_j$ its value is
$1$ and its derivatives of orders $1,\ldots,m_j-1$ are zero. Also
\[
\ell(F_*)=(\gamma-1/2)-(\gamma-p)=p-1/2\in\mathbb Q.
\]
For an empty zero set take $M=0$. Apply
Theorems~\ref{sea-analytic-density} and \ref{sea-analytic-projection}
to approximate $F_*$ by a rational combination $A$ of the weighted kernels
$k_a=(1+u)K_a$, with all these
jets and the mass exact. Near $x_j$, Taylor's theorem bounds the error by
\[
|A(u)-F_*(u)|\leq
\frac{\|(A-F_*)^{(m_j)}\|}{m_j!}|u-x_j|^{m_j}.
\]
The leading coefficient of $WG$ at that point is strictly positive;
the parity restriction ensures positivity on both sides of an interior
zero. Sufficiently close approximation makes $1-A$ positive near each
point, with precisely that order of vanishing. Away from small disjoint
neighborhoods, $WG$ has a positive lower bound, since its limit is $1$.
Uniform closeness proves positivity there too.

Thus $R=h_1(1-A)$ has the required zero set and signs. The ratio
$(1-A)/W$ extends positively and continuously across each zero and
infinity. Its positive minimum admits a smaller positive rational
$\eta$, proving the stated global bound. The mass is exact by construction;
the tail follows from $K_a(u)=O(u^{-2})$. Conversely, the global bound
has positive integral, so it implies $\gamma-p>0$.

For the arithmetic claim, take a positive integer $D$ clearing the
coefficients and both $D/2$ and $D(p-1/2)$. Then $Q\mid D$, and the
integer coefficient pair of the integrated identity is
$(D,DP/Q)$. Its gcd is $D/Q$, giving exactly $Q\gamma-P$.
\end{proof}

\paragraph{Finite verification and effective search.}
For fixed rational data, enumerate rational kernel combinations and apply
the exact projection. The denominators of $R-\eta W/(1+u)$ are positive
on the half-line. Clearing them reduces verification to an exact polynomial
sign test. Enumerating positive rational $\eta$ therefore gives a search
that terminates exactly when $p<\gamma$. No bound for its degree, height,
or running time is asserted. Most importantly, it does not select a
sequence for which $Q\gamma-P\to0$.

Two small certificates can be checked without a search:
\begin{align}
R_1&=h_1-\frac{12}{13}K_1
=\frac{u(52u^2+325u+153)}
{13(u+1)(u+4)(u+9)(4u+9)},\\
\int R_1\,d\mu&=\gamma-\frac{15}{26};\nonumber\\
R_2&=h_1+\frac{17}{5640}K_{1/2}-\frac{351}{376}K_1
=\frac{u^2(752u^2+4820u+2943)}
{47(u+1)(u+4)(u+9)(4u+1)(4u+9)},\\
\int R_2\,d\mu&=\gamma-\frac{6509}{11280}.\nonumber
\end{align}
Positive numerator coefficients prove the sign globally and give exact
orders $1$ and $2$ at the origin. These are worked finite examples,
not evidence for an unproved infinite pattern.


\section{Nested rational brackets from the canonical kernels}
Define the positive measure and its moments by
\[
 d\mu(u)=\frac{du}{e^{2\pi\sqrt u}-1},\qquad
 \mu_j=\int_0^\infty u^j\,d\mu(u)
       =\frac{|B_{2j+2}|}{2j+2}.
\]
Here $B_k$ denotes a Bernoulli number; the polynomial quantity $B_m$
used later is explicitly defined there and has a different role.
Binet's identity \cite{dlmf-binet} gives, for integers $n\ge1$,
\begin{equation}\label{stage3-binet}
 H_n-\frac1{2n}=\gamma+\log n-
       \int_0^\infty\frac{d\mu(u)}{u+n^2}.
\end{equation}
For $m\ge0$, set
\[
 P_m(x)=\prod_{j=1}^m(4^jx-1),\quad A_m=P_m(1),\quad
 s_m=\frac{P_m'(1)}{A_m},\quad
 C_m(x)=\frac{P_m(x)}{A_m}[1-s_m(x-1)].
\]
Empty products are $1$. Put
\[
 r_m=\sum_i[x^i]C_m(x)
       \left(H_{2^i}-\frac1{2^{i+1}}\right),\qquad
 E_m=\gamma-r_m.
\]
The first values are $r_0=1/2$, $r_1=31/54$, and
$r_2=27269/47250$.

\begin{lemma}[Positive kernels and a sharpened contraction]
\label{stage3-canonical-baseline}
For every $m\ge0$, with $R=4^{m+1}$ and
$\lambda_m=R+s_m(R-1)$,
\begin{equation}\label{stage3-canonical-integral}
 E_m=\int_0^\infty
 \frac{u^m(u+\lambda_m)}{\prod_{i=0}^{m+1}(u+4^i)}\,d\mu(u)>0.
\end{equation}
Define
\[
 \rho_m=\frac{9(m+1)^2}{16\cdot4^{m+1}}\le\frac9{64}.
\]
Then
\begin{equation}\label{stage3-sharpened-ratio}
 0<\frac{E_{m+1}}{E_m}<\frac23\rho_m.
\end{equation}
In particular $r_m$ increases strictly to $\gamma$ and
$E_{m+1}/E_m\to0$.
\end{lemma}

\begin{proof}
The polynomial has $C_m(1)=1$ and $C'_m(1)=0$, so
\eqref{stage3-binet} cancels every logarithm exactly. The elementary
rational identity
\[
 \sum_{i=0}^m\frac{[x^i]P_m(x)}{u+4^i}
 =\frac{A_mu^m}{\prod_{i=0}^m(u+4^i)}
\]
follows by taking a common denominator: its numerator has degree at
most $m$, its first $m$ Taylor coefficients vanish because
$P_m(4^{-j})=0$ for $1\le j\le m$, and its leading coefficient is
$P_m(1)=A_m$. Shifting the indices once evaluates the contribution
from $xP_m$ and gives \eqref{stage3-canonical-integral}.

The exact moment formula also gives the stronger estimate
\[
 \frac{\mu_{m+1}}{\mu_m}
 =\frac{(2m+3)(2m+2)}{4\pi^2}
       \frac{\zeta(2m+4)}{\zeta(2m+2)}
 <\frac{(m+1)^2}{6}.
\]
Indeed $(2m+3)(2m+2)\le6(m+1)^2$, $\pi>3$, and the zeta ratio is
strictly less than $1$. The kernels $F_m$ in
\eqref{stage3-canonical-integral} satisfy
\[
 \frac{F_{m+1}(u)}{F_m(u)}
 =\frac{u(u+\lambda_{m+1})}{(u+4R)(u+\lambda_m)}
 \le\frac{9u}{4R}.
\]
Here $s_{m+1}=s_m+R/(R-1)$ and
\[
 9\lambda_m-\lambda_{m+1}
 =\frac{R(R-4)}{R-1}+s_m(5R-8)\ge0.
\]
The function
$W_m(u)=(u+\lambda_m)/\prod_{i=0}^{m+1}(u+4^i)$
is strictly decreasing. Its logarithmic derivative is negative because
$\lambda_m\ge R$. For the measure $u^m d\mu$, multiplication by a
decreasing positive function cannot increase the mean of $u$.
This follows by integrating
$(u-v)(W_m(u)-W_m(v))\le0$ against its product measure.
Consequently
\[
 \frac{E_{m+1}}{E_m}
 \le\frac9{4R}\frac{\mu_{m+1}}{\mu_m}
 <\frac{3(m+1)^2}{8R}=\frac23\rho_m.
\]
The ratio is uniformly at most $3/32<1$, which proves convergence of
$E_m$ to zero. Since $\rho_m\to0$, the remaining assertions follow.
\end{proof}

\begin{theorem}[Strictly nested rational upper brackets]
\label{stage3-nested-brackets}
Put $w_m=r_{m+1}-r_m>0$ and
\[
 h_m=r_m+\frac{w_m}{1-\rho_m},\qquad
 \delta_m=h_m-\gamma.
\]
These quantities satisfy
\[
 r_m<r_{m+1}<\gamma<h_{m+1}<h_m,
 \qquad r_m\uparrow\gamma,\quad h_m\downarrow\gamma.
\]
More precisely,
\begin{equation}\label{stage3-upper-gap-sandwich}
 \frac{\rho_m E_m}{3(1-\rho_m)}<\delta_m
 <\frac{\rho_m E_m}{1-\rho_m},
\end{equation}
and
\[
 0<\frac{\delta_{m+1}}{\delta_m}
 <2\rho_{m+1}\frac{1-\rho_m}{1-\rho_{m+1}}
 \le\frac9{32}.
\]
Every endpoint $r_m,h_m$ is obtained using rational arithmetic only.
\end{theorem}

\begin{proof}
Since $w_m=E_m-E_{m+1}$,
\[
 \delta_m=\frac{\rho_m E_m-E_{m+1}}{1-\rho_m}.
\]
The sharpened ratio gives the lower bound in
\eqref{stage3-upper-gap-sandwich}; positivity of $E_{m+1}$ gives its
upper bound. Dividing these bounds at successive indices and using
\eqref{stage3-sharpened-ratio} gives the displayed contraction.
The sequence $\rho_m$ is nonincreasing, since
$\rho_{m+1}/\rho_m=(m+2)^2/(4(m+1)^2)\le1$.
This proves strict decrease, convergence, and the uniform constant.
\end{proof}

For reference, let $\widehat p_m/\widehat q_m$ be the first mediant
accepted by the Stern--Brocot search in $[r_m,r_{m+1}]$. Unimodularity
of the two search parents proves that this is a least-denominator
fraction in the interval. The search terminates: any fixed rational
in the target interval bounds the increasing parent denominator sum
until acceptance. All comparisons use the rational endpoints.
Thus
\begin{equation}\label{stage3-selected-form}
 0<\widehat q_m\gamma-\widehat p_m
 \le\frac{\widehat q_m w_m}{1-\rho_m}.
\end{equation}
Also $\widehat q_m\to\infty$, because finitely many rational numbers
have bounded denominator in $[r_0,1]$, and each fixed one below
$\gamma$ is eventually below $r_m$. This proves infinitely many
distinct targets, without assuming irrationality. The right side of
\eqref{stage3-selected-form} is an explicit rational upper certificate;
its convergence to zero is the unresolved condition.


\begin{theorem}[Least-denominator selection and an exact rationality dichotomy]
\label{next-cutoff-selection}
Write $r_m=r_{m,0}$ and define the positive rational numbers
\[
 w_m=r_{m+1}-r_m,\qquad
 \rho_m=\frac{9(m+1)^2}{16\cdot4^{m+1}}\le\frac9{64}.
\]
Choose a reduced fraction $\widehat p_m/\widehat q_m$ of least positive
denominator in the closed rational interval $[r_m,r_{m+1}]$.
It can be found by the Stern--Brocot interval algorithm using rational
comparisons only; take its first accepted mediant as the definite selection.
The denominators satisfy $\widehat q_m\to\infty$,
without any rationality assumption on $\gamma$. Set
\[
 \mathcal U_m=\frac{\widehat q_m w_m}{1-\rho_m}.
\]
The associated forms have integer coefficients and satisfy the positive
rational certificate
\begin{equation}\label{next-cutoff-selected-certificate}
 0<\widehat q_m\gamma-\widehat p_m\le\mathcal U_m.
\end{equation}
For every $j\ge m+1$, the sharper fully rational bound
\begin{equation}\label{next-cutoff-refined-certificate}
 \widehat q_m\gamma-\widehat p_m
 \le
 \widehat q_m\left(r_j-\frac{\widehat p_m}{\widehat q_m}
                   +\frac{w_j}{1-\rho_j}\right)
\end{equation}
also holds.

The limit of the completely rational sequence $\mathcal U_m$ exists,
with the exact dichotomy
\begin{equation}\label{next-cutoff-dichotomy}
 \lim_{m\to\infty}\mathcal U_m=
 \begin{cases}
  0,&\gamma\notin\mathbb Q,\\
  1/b,&\gamma=a/b\text{ in lowest terms}.
 \end{cases}
\end{equation}
In the rational case the primitive errors
$\widehat q_m\gamma-\widehat p_m$ are eventually exactly $1/b$.
Consequently proving $\mathcal U_m\to0$ would prove gamma irrationality;
the theorem does not establish which branch in
\eqref{next-cutoff-dichotomy} holds.
\end{theorem}

\begin{proof}
Let $F_m$ denote the kernel in \eqref{next-cutoff-positive-kernel} at
$t=0$. Put $E_m=\gamma-r_m>0$ and
$\lambda_m=4^{m+1}+s_m(4^{m+1}-1)$. The exact kernels give
\[
 \frac{F_{m+1}(u)}{F_m(u)}
 =\frac{u(u+\lambda_{m+1})}
 {(u+4^{m+2})(u+\lambda_m)}.
\]
Writing $R=4^{m+1}$ and $d=R-1$, the identity
$s_{m+1}=s_m+R/d$ shows
\[
 \lambda_{m+1}
 =4R+s_m(4R-1)+\frac{R(4R-1)}d
 \le9\lambda_m.
\]
Therefore $F_{m+1}/F_m\le9u/(4R)$.
The function
\[
 \frac{u+\lambda_m}{\prod_{i=0}^{m+1}(u+4^i)}
\]
is strictly decreasing, since $\lambda_m\ge4^{m+1}$ and its logarithmic
derivative is negative. The same weighted-mean argument as before thus
gives
\[
 \frac{E_{m+1}}{E_m}
 \le\frac9{4R}\frac{\mu_{m+1}}{\mu_m}<\rho_m.
\]
In particular $\rho_m\to0$, and
\[
 w_m<E_m\le\frac{w_m}{1-\rho_m}.
\]
These bounds prove \eqref{next-cutoff-selected-certificate} and
\eqref{next-cutoff-refined-certificate}, because
$\widehat p_m/\widehat q_m\le r_{m+1}<\gamma$.

For completeness, start the Stern--Brocot search with Farey neighbours
$0/1$ and $1/0$. Replace the lower neighbour by the mediant when the
mediant is below $r_m$, and replace the upper neighbour when it is above
$r_{m+1}$. Stop when the mediant lies in the target interval.
The elementary integral comparison gives $\gamma\le3/2-\log2<1$,
so the first mediant $1/1$ replaces the upper infinite endpoint.
Thereafter both neighbour denominators are positive.
The neighbours $a/b<c/d$ always satisfy $bc-ad=1$.
Every reduced fraction strictly between them has denominator at least
$b+d$, because its coefficient representation in these two unimodular
columns consists of positive integers. The first accepted mediant
therefore has least denominator. Termination follows by keeping any one
rational in the nonempty target interval: until acceptance the sum of
the neighbour denominators increases and is bounded by that rational's
denominator. This also proves that the selection uses no gamma oracle.

For any fixed denominator bound, there are only finitely many reduced
fractions in the bounded interval $[r_0,\gamma)$ satisfying that bound.
Each is eventually smaller than $r_m$, because $r_m\uparrow\gamma$.
Consequently none can be selected for all sufficiently large $m$, which
proves $\widehat q_m\to\infty$ and supplies infinitely many distinct targets.

Suppose first that $\gamma$ is irrational. Fix $0<\eta<1/2$.
An irrational rotation visits every fixed open arc with bounded gaps.
Here is the elementary compactness argument: the pigeonhole principle
makes the subgroup generated by $\gamma$ dense modulo $1$; finitely many
translates of a fixed open arc cover the circle. A common shift makes
their finitely many integer indices nonnegative. Their maximum then
bounds a return time from every starting point.

Apply this to the arc $(\eta/4,\eta/2)$. For all large $m$, there is an
integer
\[
 \frac{\eta}{2E_m}\le q\le\frac{\eta}{E_m}
 \quad\hbox{with}\quad
 \frac\eta4<\{q\gamma\}<\frac\eta2.
\]
For $p=\lfloor q\gamma\rfloor$ this gives
\[
 E_m/4<\gamma-p/q<E_m.
\]
Since $E_{m+1}/E_m\to0$, the fraction lies in $[r_m,r_{m+1}]$ for all
large $m$. Reduction can only lower its denominator, so the minimality
of $\widehat q_m$ gives $\widehat q_m E_m\le\eta$.
As $\eta$ is arbitrary, $\widehat q_m E_m\to0$ and hence
$\mathcal U_m\to0$.

Now suppose $\gamma=a/b$ in lowest terms. Every fraction $p/q<\gamma$
has gap at least $1/(bq)$, so
$\widehat q_m E_m\ge1/b$.
Choose the smallest integer $q\ge1/(bE_m)$ in the residue class
$aq\equiv1\pmod b$. Then
\[
 \frac1{bE_m}\le q<\frac1{bE_m}+b,\qquad
 p=\frac{aq-1}{b}\in\mathbb Z,\qquad \gcd(p,q)=1.
\]
Its gap satisfies
\[
 \frac{E_m}{1+b^2E_m}<\gamma-p/q\le E_m.
\]
For large $m$ this gap exceeds $E_{m+1}$, so the fraction is in the target
interval. Therefore
\[
 \frac1b\le\widehat q_m E_m\le\frac1b+bE_m,
\]
and $\widehat q_m E_m\to1/b$. Since
$w_m/E_m=1-E_{m+1}/E_m\to1$ and $\rho_m\to0$,
this proves the rational branch of \eqref{next-cutoff-dichotomy}.
Finally each
$\widehat q_m\gamma-\widehat p_m$ is a positive integral multiple of
$1/b$ and is at most $\widehat q_m E_m\to1/b$.
It must therefore equal $1/b$ for all sufficiently large $m$.
\end{proof}




\section{Logarithmic Farey reduction of the primitive optimization}
\label{stage3-farey-section}

Let $\alpha=a/b<\beta=c/d<\xi$, where the endpoint fractions are reduced
and $b,d>0$. Set $v_0=(b,a)$, $v_1=(d,c)$ and
$\Delta=\det(v_0,v_1)=bc-ad>0$. The linear functional
\[
 \mathcal L_\xi(q,p)=q\xi-p
\]
is strictly positive on every nonzero integer vector in the cone spanned
by $v_0,v_1$. The following reduction replaces enumeration of a fundamental
parallelogram of size $\Delta$ by a logarithmic number of candidates.

\begin{theorem}[A logarithmic candidate set for the exact primitive minimum]
\label{stage3-farey-logarithmic-sail}
One can construct, using integer arithmetic only, a set $G$ of primitive
vectors $(q,p)$ with $q>0$ and $\alpha\le p/q\le\beta$ such that
\[
 |G|\le\lfloor\log_2\Delta\rfloor+2
\]
and, simultaneously for every real $\xi>\beta$,
\begin{equation}\label{stage3-farey-sail-minimum}
 \min_{\substack{\alpha\le p/q\le\beta\\q>0,\ \gcd(p,q)=1}}
       (q\xi-p)
 =
 \min_{(q,p)\in G}(q\xi-p).
\end{equation}
More generally the same equality holds for every linear functional
strictly positive on the nonzero cone.
After one extended Euclidean computation for $a,b$, the construction uses
at most $2\lfloor\log_2\Delta\rfloor+1$ ceiling-Euclidean updates when
$\Delta>1$. Each collinear run is handled by a single integer division.
For $\Delta=1$, the two endpoint vectors suffice.

The resulting lower envelope is a finite rational piecewise-linear
function of $\xi$. No test for whether $\xi$ equals one of its rational
breakpoints is part of the construction.
\end{theorem}

\begin{proof}
The case $\Delta=1$ follows from the unimodular endpoint basis. Suppose
$\Delta>1$. Choose integers $u,v$ with $bu+av=1$ and put
$r_{\rm raw}=ud+vc$. A unimodular shear after
\[
 U=\begin{pmatrix}u&v\\-a&b\end{pmatrix}
\]
sends the two endpoint vectors to $(1,0)$ and $(r,\Delta)$, where
$r\equiv r_{\rm raw}\pmod\Delta$ and $1\le r<\Delta$.
The last assertion follows because the upper endpoint is primitive,
so $\gcd(r,\Delta)=1$.

In these transformed coordinates start with
\[
 z_0=(1,0),\quad z_1=(1,1),\qquad
 D_0=\Delta,\quad D_1=\Delta-r.
\]
While $D_i>0$, use
\[
 h_i=\left\lceil\frac{D_{i-1}}{D_i}\right\rceil\ge2,\qquad
 z_{i+1}=h_i z_i-z_{i-1},\qquad
 D_{i+1}=h_iD_i-D_{i-1}.
\]
Then
\[
 0\le D_{i+1}<D_i,\qquad
 \det(z_i,z_{i+1})=1,\qquad
 D_i=\det(z_i,(r,\Delta)).
\]
The coordinates are nondecreasing, with increasing second coordinate,
so every ray remains in the cone and the terminal vector is exactly
$(r,\Delta)$: it is a positive parallel primitive vector.
The consecutive rays form a unimodular fan. Every integer vector in
the cone is therefore a nonnegative integer combination of two adjacent
fan vectors. A positive linear functional on such a nonzero combination
is at least its value on one of those generators. Thus all fan vectors
form an exact candidate set.

When $h_i=2$, the vector $z_i$ lies on the line segment between
$z_{i-1}$ and $z_{i+1}$. Throughout a run of such updates the vectors
form an arithmetic progression, and a linear functional is minimized
at one of the run endpoints. Consequently it is enough to retain
$z_0$, the terminal vector, and precisely the vectors $z_i$ for which
$h_i\ge3$.

Every retained interior vertex satisfies
$D_i<D_{i-1}/2$. Since the positive determinant remainders decrease,
there are at most $\lfloor\log_2\Delta\rfloor$ such events.
This proves the candidate count after applying the inverse unimodular
map. All candidate vectors remain primitive, and their first coordinates
are positive because they lie in the original cone.

To skip a run with $h_i=2$, put $g=D_{i-1}-D_i>0$ and
$k=\lfloor D_i/g\rfloor\ge1$. The next $k$ updates all have quotient $2$,
and can be replaced by
\[
 z_{i+k}=z_i+k(z_i-z_{i-1}),\qquad D_{i+k}=D_i-kg.
\]
The predecessor is obtained by replacing $k$ by $k-1$. Either the
remainder becomes zero or the next quotient is at least $3$.
There is at most one such run between retained interior vertices,
giving the stated update bound. Equality
\eqref{stage3-farey-sail-minimum} follows from generation and removal of
collinear candidates, and rationality of the breakpoints follows by
equating two integer affine functions.
\end{proof}

\begin{theorem}[A terminating relative optimizer without an equality oracle]
\label{stage3-farey-certified-optimizer}
Let $G$ be the preceding finite set, and suppose rational brackets
$L_j\le\xi\le U_j$ are available with $L_j>\beta$ and
$U_j-L_j\to0$. Define
\[
 A_j=\min_{(q,p)\in G}(qL_j-p),\qquad
 B_j=\min_{(q,p)\in G}(qU_j-p).
\]
Writing $M_\xi$ for the primitive minimum in
\eqref{stage3-farey-sail-minimum}, one has
\[
 0<A_j\le M_\xi\le B_j,\qquad
 B_j-A_j\le Q_{\max}(U_j-L_j),
 \quad Q_{\max}=\max_{(q,p)\in G}q.
\]
For every rational $\tau>0$, refinement until
$B_j\le(1+\tau)A_j$ terminates. If $\widehat v_j$ is any candidate
attaining $B_j$, then
\[
 M_\xi\le\mathcal L_\xi(\widehat v_j)
 \le B_j\le(1+\tau)M_\xi.
\]
At any step all candidates with $qL_j-p>B_j$ may be discarded without
losing a true optimum.

For the canonical gamma intervals $[r_m,r_{m+1}]$, the rational
brackets
\[
 L_j=r_j,\qquad
 U_j=r_j+\frac{r_{j+1}-r_j}{1-\rho_j},
 \qquad j\ge m+2,
\]
supply such an optimizer. This gives an explicit infinite family of
positive primitive forms within any prescribed fixed relative factor
of the true interval optimum. It does not prove that those optima
tend to zero.
\end{theorem}

\begin{proof}
For each candidate, replacing $\xi$ by $L_j$ or $U_j$ gives the stated
lower and upper bounds. If a vector minimizes the lower expression, its
upper expression is at most
$A_j+Q_{\max}(U_j-L_j)$; minimizing the upper expression proves the
difference estimate. The finite set has a strictly positive minimum
$M_\xi$, so both $A_j$ and $B_j$ tend to $M_\xi$. Therefore the relative
stopping condition holds after finitely many refinements. This argument
does not require identifying a unique optimum or deciding whether
$\xi$ equals a rational breakpoint. The bound for an upper-minimizing
candidate follows directly. A discarded vector has actual value greater
than $B_j$, whereas an undiscarded upper-minimizing vector has actual
value at most $B_j$, so discarding is safe. The previously proved
canonical gamma bounds verify all hypotheses of the application.
\end{proof}

\begin{theorem}[A three-fraction certificate for an exact primitive optimum]
\label{stage3-farey-three-fractions}
Suppose $a/b<c/d$ are Farey neighbours, so $bc-ad=1$, and
\[
 M=\frac{a+c}{b+d}<\xi<\frac cd.
\]
Then for every reduced fraction $p/q<M$ with $q>0$,
\[
 q\xi-p\ge b\xi-a>0,
\]
with equality only for $p/q=a/b$. Consequently, whenever
\[
 \alpha\le\frac ab\le\beta<M<\xi<\frac cd,
\]
the fraction $a/b$ is the unique primitive-error minimizer on
$[\alpha,\beta]$, and
\begin{equation}\label{stage3-farey-three-error}
 0<b\xi-a
 =\frac1d-b\left(\frac cd-\xi\right)<\frac1d.
\end{equation}
For gamma, both strict comparisons must be certified: choose rational
brackets with $M<L_k\le\gamma\le U_j<c/d$.
Thus three fractions and rational comparisons certify an exact interval
optimum without an equality oracle. In the canonical application below,
the next interval's selected mediant satisfies $M<r_{m+2}$, while a
computed upper bracket lies below $c/d$.
\end{theorem}

\begin{proof}
First suppose $p/q\le a/b$. If $q\ge b$, then
$q(\xi-p/q)\ge b(\xi-a/b)$, with equality only at the same primitive
fraction. If $q<b$, the strict upper bound $\xi<c/d$ gives
\[
 (q\xi-p)-(b\xi-a)
 >
 \frac{cq-dp-1}{d}\ge0,
\]
because $cq-dp$ is a positive integer.

If $a/b<p/q<M$, unimodularity gives the positive integers
\[
 s=cq-dp,\qquad t=bp-aq,\qquad
 (q,p)=s(b,a)+t(d,c).
\]
The inequality $p/q<M$ is equivalent to $s>t$. Hence
\[
 (q,p)=(s-t)(b,a)+t(b+d,a+c).
\]
Both associated linear forms are strictly positive because
$a/b<M<\xi$. Their positive integer combination has value strictly
greater than $b\xi-a$. This proves the minimum statement.
Finally $bc-ad=1$ gives
$c/d-a/b=1/(bd)$ and therefore
\eqref{stage3-farey-three-error}.
\end{proof}

\begin{proposition}[A quantitative parent-denominator target]
\label{stage3-farey-parent-target}
For the least-denominator fraction of $[r_m,r_{m+1}]$, let its final
Stern--Brocot parents be $a_m/b_m<c_m/d_m$. Then
\[
 \widehat q_m=b_m+d_m,\qquad
 \widehat q_m w_m<
 \frac1{b_m}+\frac1{d_m},
\]
and consequently
\begin{equation}\label{stage3-farey-parent-upper}
 \mathcal U_m<
 \frac{b_m^{-1}+d_m^{-1}}{1-\rho_m}.
\end{equation}
The lower-parent denominators satisfy $b_m\to\infty$ without assuming
gamma irrationality. The upper-parent denominators satisfy the exact
alternative:
\begin{itemize}
\item if gamma is irrational, then $d_m\to\infty$;
\item if gamma is rational, then $c_m/d_m$ is eventually exactly gamma,
so $d_m$ is eventually its reduced denominator.
\end{itemize}
In particular, an independent proof that $d_m$ is unbounded would
suffice for condition 3. This proposition does not supply that proof.
\end{proposition}

\begin{proof}
The selected fraction is the mediant of its final parents, so its
denominator is $b_m+d_m$. The target interval lies strictly inside
their Farey interval of length $1/(b_md_m)$, proving the quantitative
inequalities.

The selected denominators tend to infinity: otherwise a fixed rational
would lie in infinitely many intervals converging strictly from below
to gamma, forcing that rational to equal gamma, a contradiction.
If the lower denominators had a bounded subsequence, finitely many
possible lower parents would give a constant subsubsequence. The upper
denominators would then tend to infinity, and the Farey gap would tend
to zero, forcing that fixed lower parent to equal gamma. This contradicts
its being strictly below every target interval. Thus $b_m\to\infty$.

The same compactness argument shows that a bounded subsequence of
$d_m$ forces gamma to be rational. Conversely, if gamma is a reduced
rational, the finite Stern--Brocot path reaches it. For all sufficiently
late intervals lying strictly below and sufficiently close to it, this
rational becomes the right parent during that common initial path.
Thereafter mediants with that right parent increase towards gamma.
If a mediant overshot the target interval but still lay below gamma,
it would replace the right parent. For late intervals this cannot occur
before a mediant enters the interval: consecutive such mediants have
gaps from gamma of the form $1/[d(b+kd)]$, with successive gap ratio
tending to $1$, whereas $E_{m+1}/E_m\to0$.
More explicitly, choose the first $k$ with gap at most $E_m$.
Its gap is at least $E_m/(1+d^2E_m)$, which exceeds $E_{m+1}$ for all
large $m$. It therefore lies inside the target interval and the search
stops with right parent gamma unchanged. This proves eventual constancy.
\end{proof}

\begin{proposition}[An exact integer test for the final upper Farey parent]
\label{stage3-farey-upper-parent-test}
Let $0<\alpha=p_0/q_0<\beta=p_1/q_1<1$ be reduced rational endpoints.
For a proposed reduced upper parent $c/d$ with
$\beta<c/d\le1$, put
\[
 k_i=cq_i-dp_i>0\quad(i=0,1).
\]
Choose $b_0\in\{0,\ldots,d-1\}$ with $b_0c\equiv1\pmod d$;
for $d=1$ use $b_0=0$. Define the integer
\[
 B_*=b_0+d\left\lfloor
       \frac{q_0-1-b_0k_0}{dk_0}
       \right\rfloor.
\]
Then $c/d$ is exactly the final upper parent in the Stern--Brocot
least-denominator search on $[\alpha,\beta]$ if and only if
\begin{equation}\label{stage3-farey-upper-parent-integer-test}
 B_*>0,\qquad (B_*+d)k_1\le q_1.
\end{equation}
When these inequalities hold, the lower parent is
\[
 \frac{A_*}{B_*},\qquad A_*=\frac{B_*c-1}{d},
\]
and the selected fraction is
$(A_*+c)/(B_*+d)$.

For the canonical endpoints, this translates bounded upper-parent
denominators into explicit floor and congruence inequalities involving
only their harmonic numerator and denominator pairs. In particular,
\[
 \gcd(q_i,k_i)=\gcd(q_i,d).
\]
The existence of many primes dividing $q_i$ and not dividing $d$ merely
makes $k_i$ a unit at those primes; it does not by itself contradict
\eqref{stage3-farey-upper-parent-integer-test}.
\end{proposition}

\begin{proof}
Set
\[
 X=\frac{q_0}{k_0}=\frac1{d(c/d-\alpha)},\qquad
 Y=\frac{q_1}{k_1}=\frac1{d(c/d-\beta)}.
\]
A Farey lower neighbour $A/B$ of $c/d$ must satisfy
$Bc-Ad=1$, so $Bc\equiv1\pmod d$.
The conditions $A/B<\alpha$ and
$(A+c)/(B+d)\ge\alpha$ are respectively
\[
 B<X,\qquad B+d\ge X.
\]
There is exactly one integer in the residue class $b_0\bmod d$
satisfying these two inequalities: the largest such integer strictly
below $X$. Since $Bk_0<q_0$ is equivalent to $Bk_0\le q_0-1$,
that integer is exactly $B_*$ in the statement.
It must be positive. The remaining condition that the mediant be at
most $\beta$ is $B_*+d\le Y$, which is precisely the second inequality
in \eqref{stage3-farey-upper-parent-integer-test}.

Conversely these inequalities give a Farey pair whose open interval
contains $[\alpha,\beta]$ and whose mediant belongs to that target
interval. Every fraction in the open Farey interval has denominator at
least the sum of the parent denominators, so the mediant is the unique
least-denominator fraction in the target interval. A reduced fraction
in $(0,1)$ has a unique Farey parent pair in $[0,1]$ with positive parent
denominators: the right denominator is determined by the inverse
congruence modulo the mediant denominator. Thus this pair is exactly
the final pair from the Stern--Brocot search.
The gcd assertion follows from $\gcd(p_i,q_i)=1$.
\end{proof}

\paragraph{Why positivity and prime-rich denominators alone do not settle the target.}
The obstruction can be exhibited exactly for the rational limit $\xi=1/2$.
Choose integers $D_m$ divisible by $4$, with
$D_{m+1}\ge D_m^2$ and $D_m\to\infty$, and put
\[
 \ell_m=\frac12-\frac1{D_m}.
\]
The endpoint $\ell_m$ has reduced denominator exactly $D_m$, so the
$D_m$ can be chosen to include any prescribed increasing lists of primes.
The errors have ratio $D_m/D_{m+1}\to0$. They even have positive integral
representations against the same Binet measure: since $\mu_0=1/12$,
\[
 \xi-\ell_m=\int\frac{12}{D_m}\,d\mu>0.
\]
Yet the least-denominator rational in
$[\ell_m,\ell_{m+1}]$ is exactly
\[
 \frac{D_m/4}{D_m/2+1},
\]
and its primitive error is exactly $1/2$ for every $m$.
Indeed odd denominators require $q\ge D_m/2$ and the smallest eligible
odd denominator is $D_m/2+1$; even denominators require $q\ge D_m$.
The displayed fraction belongs to the interval because
$D_{m+1}\ge D_m+2$.
Thus arbitrarily fast raw convergence, global positivity, and prescribed
prime factors of the endpoint denominators do not by themselves imply
primitive-error decay. This example does not imitate every
gamma-specific harmonic coefficient identity.



\subsection{Exact finite optima and a rational numerical certificate}
In each of the eleven canonical intervals $m=0,\ldots,10$, the
relative optimizer's surviving set reduces to one vector.
The same vector is independently certified by the three-fraction
criterion. The following selected rows show the actual optimum,
the number of logarithmic-sail candidates, and an upward rational
rounding of its positive primitive error.

\begin{center}
\small
\renewcommand{\arraystretch}{1.55}
\begin{tabular}{@{}rcrl@{}}
\toprule
$m$ & Exact optimal target & Candidates & Strict upper error bound\\
\midrule
0 & $\dfrac{4}{7}$ & 3 & $4.05096544\times10^{-2}$ \\[3pt]
1 & $\dfrac{15}{26}$ & 6 & $7.60728744\times10^{-3}$ \\[3pt]
2 & $\dfrac{157}{272}$ & 15 & $2.66085322\times10^{-3}$ \\[3pt]
4 & $\dfrac{33841}{58628}$ & 49 & $1.84706856\times10^{-6}$ \\[3pt]
6 & $\dfrac{9062634}{15700603}$ & 172 & $1.53583719\times10^{-9}$ \\[3pt]
8 & $\dfrac{8788062966751}{15224921118955}$ & 723 & $9.04971062\times10^{-14}$ \\[3pt]
10 & $\dfrac{2323006970406828281}{4024504377931443899}$ & 2728 & $3.81603071\times10^{-19}$ \\[3pt]
\bottomrule
\end{tabular}
\normalsize
\end{center}

For $m=10$, the certifying triple is
\[
 \frac{2323006970406828281}{4024504377931443899},\qquad
 \frac{3378631047479839813}{5853325287102525206},\qquad
 \frac{1055624077073011532}{1828820909171081307}.
\]
The first and third fractions are Farey neighbours; the middle is
their mediant. The first lies in $[r_{10},r_{11}]$, while the middle
lies strictly above $r_{11}$ and strictly below $r_{12}<\gamma$.
The rational upper bracket $h_{12}$ is below the third fraction.
Thus the first fraction is the unique primitive-error optimum over
the entire interval $[r_{10},r_{11}]$.

A sharper, independently constructed rational Binet bracket proves
\[
 \begin{gathered}
 \mathcal F_{10}
 =4024504377931443899\,\gamma-2323006970406828281,\\
 \frac{38160307067862345768}{10^{38}}
 <\mathcal F_{10}
 <\frac{38160307067862345769}{10^{38}}.
 \end{gathered}
\]
Here is the exact bracket construction used by the verifier.
Put
\[
 \ell_-=
 2\sum_{j=0}^{109}\frac1{(2j+1)3^{2j+1}},\qquad
 \ell_+=\ell_-+\frac{2}{221\cdot3^{221}}\frac98.
\]
The positive series for $2\operatorname{artanh}(1/3)$ gives
$\ell_-<\log2<\ell_+$. Define
\[
 S=\sum_{j=0}^{19}\frac{(-1)^j\mu_j}{64^{2j+2}},
 \qquad R=\frac{\mu_{20}}{64^{42}}.
\]
Finite division of $1/(64^2+u)$ with an even number of terms shows
\[
 S<\int_0^\infty\frac{d\mu(u)}{64^2+u}<S+R.
\]
Consequently the entirely rational endpoints
\[
 g_-=H_{64}-\frac1{128}-6\ell_++S,\qquad
 g_+=H_{64}-\frac1{128}-6\ell_-+S+R
\]
satisfy $g_-<\gamma<g_+$. Comparing the desired integer form at
these two endpoints proves the displayed strict inequalities by
integer cross multiplication.

These are finite optimum certificates. They neither prove that the
interval optima tend to zero nor assert a computational record for
excluding rational denominators of $\gamma$.




% Additional synthesis for the unified gamma manuscript, 31 October 2026.
% Body fragment; the host supplies theorem environments and amsmath/amssymb.
% References use labels preserved from the original constructive-selection
% and stage-three manuscripts. No assertion of historical priority is made.

\section{Selection and primitive optimization on shrinking intervals}
\label{unified-selection-section}

The next two results clarify the relationship between least-denominator
selection and optimization of the actual primitive error. The first is a
synthesis of the preceding constructions with standard irrational-rotation
and Farey arguments. Its interval hypothesis is weaker than the contraction
available for the canonical gamma kernels. The second shows that the two
selection rules can nevertheless have arbitrarily different finite errors.
Neither result determines whether Euler's constant is rational.

For a real number $\xi$, define
\[
 \kappa(\xi)=
 \begin{cases}
  0,&\xi\notin\mathbb Q,\\
  1/b,&\xi=a/b\text{ in lowest terms, }b>0.
 \end{cases}
\]
All rational targets below are reduced, with positive denominators. Their
primitive errors are $q\xi-p$; positivity follows from $p/q<\xi$.

\begin{theorem}[A common limiting dichotomy for three selection rules]
\label{unified-selection-window-dichotomy}
Let $\xi\in\mathbb R$ and let
\[
 I_m=[\alpha_m,\beta_m],\qquad
 \alpha_m,\beta_m\in\mathbb Q,
 \qquad \alpha_m<\beta_m<\xi.
\]
Put
\[
 E_m=\xi-\alpha_m,\qquad E_m^+=\xi-\beta_m.
\]
Assume
\begin{equation}\label{unified-selection-width-hypothesis}
 E_m\longrightarrow0,
 \qquad \limsup_{m\to\infty}\frac{E_m^+}{E_m}<1.
\end{equation}
No monotonicity or nesting of the intervals is required.

Choose $p_m/q_m\in I_m$ of least positive denominator and set
\[
 S_m=q_m\xi-p_m,\qquad
 M_m=\min_{\substack{p/q\in I_m\\q>0,\ \gcd(p,q)=1}}
                     (q\xi-p).
\]
The minimum defining $M_m$ is attained and is strictly positive. Then
\begin{equation}\label{unified-selection-three-limits}
 q_m\longrightarrow\infty,
 \qquad
 \lim q_mE_m=\lim S_m=\lim M_m=\kappa(\xi).
\end{equation}
If $\xi=a/b$ is rational in lowest terms, then
\begin{equation}\label{unified-selection-rational-stabilization}
 S_m=M_m=1/b
 \qquad\text{for all sufficiently large }m.
\end{equation}

More generally, suppose a reduced target $P_m/Q_m\in I_m$ satisfies
\begin{equation}\label{unified-selection-relative-rule}
 M_m\le T_m:=Q_m\xi-P_m\le(1+\tau_m)M_m,
 \qquad 0<\tau_m<1.
\end{equation}
Then $Q_m\to\infty$ and $T_m\to\kappa(\xi)$. In the rational case,
$T_m=1/b$ eventually. If additionally $\tau_m\to0$, then
\begin{equation}\label{unified-selection-asymptotic-optimality}
 \frac{T_m}{M_m}\longrightarrow1.
\end{equation}
When the endpoints and narrowing rational brackets for $\xi$ are
effectively available, the relative rule can be implemented by
Theorem~\ref{stage3-farey-certified-optimizer}, for any effectively given
rational sequence $0<\tau_m<1$. It does not require an equality oracle.

Finally, any positive numbers $C_m$ with $C_m/E_m\to1$ satisfy
\begin{equation}\label{unified-selection-upper-limit}
 q_m C_m\longrightarrow\kappa(\xi).
\end{equation}
In particular this applies to any asymptotically exact rational upper
certificate for $E_m$.
\end{theorem}

\begin{proof}
For a fixed interval, take any admissible primitive pair $(q_*,p_*)$.
Since $E_m^+>0$, every admissible pair satisfies
\[
 q\xi-p=q(\xi-p/q)\ge qE_m^+.
\]
Pairs with $q>(q_*\xi-p_*)/E_m^+$ cannot improve the chosen pair.
Only finitely many fractions in a bounded interval have a bounded
denominator, proving existence and positivity of $M_m$.
The same finiteness observation shows that every sequence of rational
targets chosen in $I_m$ has denominator tending to infinity: otherwise
a fixed rational would occur along a subsequence converging to $\xi$
strictly from below, forcing that rational to equal $\xi$.

Choose a fixed number $\theta\in(0,1)$ such that
\[
 E_m^+/E_m\le\theta
\]
for all sufficiently large $m$. We first suppose that $\xi$ is irrational.
We use only the bounded-return property of an irrational rotation, not
equidistribution. The elementary density of the forward orbit of an
irrational rotation implies that, for every nonempty open arc $J$ of
$\mathbb R/\mathbb Z$, the sets
$J-j\xi$, $j\ge0$, cover the circle. Compactness gives a finite subcover.
Consequently there is an integer $K_J$ such that, from every starting
point, some iterate with index between $0$ and $K_J$ lies in $J$.
The density assertion itself follows from the usual pigeonhole
approximation of an irrational rotation by arbitrarily small nonzero
rotations.

Fix $h>0$ with $\theta(1+h)<1$, and fix $0<\eta<1$. Apply the bounded
return property to
\[
 J=\bigl(\theta(1+h)\eta,\eta\bigr).
\]
For all sufficiently large $m$, the interval
$[\eta/E_m,(1+h)\eta/E_m]$ is long enough to contain such a return.
Thus there is an integer $q$ with
\[
 \frac{\eta}{E_m}\le q\le\frac{(1+h)\eta}{E_m},
 \qquad \theta(1+h)\eta<\{q\xi\}<\eta.
\]
Put $p=\lfloor q\xi\rfloor$. Then
\[
 E_m^+\le\theta E_m
   <\xi-p/q<E_m.
\]
After reduction the fraction still belongs to $I_m$ and its denominator
does not increase. Minimality of $q_m$ therefore gives
\[
 q_mE_m\le(1+h)\eta.
\]
Since $\eta>0$ is arbitrary, $q_mE_m\to0$. The inequalities
\[
 0<M_m\le S_m\le q_mE_m
\]
prove the other two limits in the irrational case.

Now let $\xi=a/b$ be rational in lowest terms. Every positive primitive
error is a positive integer multiple of $1/b$, so
\[
 q_mE_m\ge S_m\ge M_m\ge1/b.
\]
Choose the smallest integer $q\ge1/(bE_m)$ in the residue class
$aq\equiv1\pmod b$, and put $p=(aq-1)/b$. For large $m$, this is a
positive denominator and
\[
 \frac1{bE_m}\le q<\frac1{bE_m}+b,
 \qquad \gcd(p,q)=1,
\]
because $aq-bp=1$. Its gap satisfies
\[
 \frac{E_m}{1+b^2E_m}<\xi-p/q\le E_m.
\]
For sufficiently large $m$, the left bound exceeds $\theta E_m$ and
therefore exceeds $E_m^+$. This fraction belongs to $I_m$. Hence
\[
 \frac1b\le q_mE_m\le\frac1b+bE_m,
\]
which proves $q_mE_m\to1/b$. Since $S_m$ is a positive integer multiple
of $1/b$ and $S_m\le q_mE_m$, eventually $S_m=1/b$. The inequalities
$1/b\le M_m\le S_m$ then give $M_m=1/b$ as well.

For a relative optimizer, the irrational case follows from
$0<T_m<2M_m\to0$. In the rational case, eventually
\[
 1/b\le T_m\le(1+\tau_m)/b<2/b.
\]
Its membership in $(1/b)\mathbb Z$ forces $T_m=1/b$.
If $\tau_m\to0$, division of
\eqref{unified-selection-relative-rule} by $M_m>0$ proves
\eqref{unified-selection-asymptotic-optimality}. The effective assertion
uses the finite candidate set of
Theorem~\ref{stage3-farey-logarithmic-sail} and the terminating refinement
rule of Theorem~\ref{stage3-farey-certified-optimizer}; for each fixed
$m$, its strict gap $\xi-\beta_m>0$ lets the available lower rational
bracket eventually exceed $\beta_m$.
Finally \eqref{unified-selection-upper-limit} follows by multiplying
$q_mE_m\to\kappa(\xi)$ by $C_m/E_m\to1$.
\end{proof}

\paragraph{The exact Farey-parent alternative.}
Under the hypotheses of
Theorem~\ref{unified-selection-window-dichotomy}, the parent criterion
also holds for the general intervals, with a short proof that uses the
rational stabilization just established. For all sufficiently large $m$,
$q_m>1$. Let
\[
 \frac{a_m}{b_m}<\frac{p_m}{q_m}<\frac{c_m}{d_m},
 \qquad b_m+d_m=q_m,\quad a_m+c_m=p_m,\quad
 b_mc_m-a_md_m=1,
\]
be its canonical Farey parents, whose denominators lie between $1$ and
$q_m-1$. They are the final Stern--Brocot parents in the positive-interval
application. Since either parent has smaller denominator than $q_m$,
minimality implies
\[
 a_m/b_m<\alpha_m<\beta_m<c_m/d_m.
\]
Their gap is $1/(b_md_m)$, and consequently
\begin{equation}\label{unified-selection-parent-gap}
 q_m(\beta_m-\alpha_m)<\frac1{b_m}+\frac1{d_m}.
\end{equation}

One always has $b_m\to\infty$. Indeed, on a subsequence with bounded
$b_m$, one has $d_m=q_m-b_m\to\infty$. The parent gap then tends to
zero and forces the lower parents to converge to $\xi$. A bounded
denominator leaves only finitely many fractions near $\xi$, so some
fixed lower parent would equal $\xi$, contradicting its strict position
below $\alpha_m<\xi$.
The same argument applied to a bounded subsequence of $d_m$ forces
$\xi$ to be rational. Thus irrational $\xi$ implies $d_m\to\infty$.

Conversely, if $\xi=a/b$ is reduced, the theorem gives
$aq_m-bp_m=1$ for all sufficiently large $m$. Also $q_m>b$ eventually.
The right Farey-parent denominator is the unique integer
$d\in\{1,\ldots,q_m-1\}$ satisfying $p_md\equiv-1\pmod{q_m}$.
The displayed determinant identity shows that it is $b$, with numerator
$a$. Hence
\begin{equation}\label{unified-selection-parent-stabilization}
 \frac{c_m}{d_m}=\xi\quad\text{eventually if }\xi\in\mathbb Q.
\end{equation}
In this class of interval sequences the following statements are
therefore equivalent:
\[
 \begin{aligned}
 \xi\notin\mathbb Q
 &\ \Longleftrightarrow\ d_m\to\infty
    \ \Longleftrightarrow\ (d_m)\text{ is unbounded}\\
 &\ \Longleftrightarrow\ M_m\to0
    \ \Longleftrightarrow\ S_m\to0
    \ \Longleftrightarrow\ q_mE_m\to0.
 \end{aligned}
\]
These are exact equivalences, not an independent proof of any one of
them for $\xi=\gamma$.

\paragraph{Application to the canonical gamma intervals.}
Take $\alpha_m=r_m$, $\beta_m=r_{m+1}$ and $\xi=\gamma$.
Lemma~\ref{stage3-canonical-baseline} supplies the stronger condition
$E_m^+/E_m=E_{m+1}/E_m\to0$, so the theorem applies. The rational upper
quantity in Theorem~\ref{next-cutoff-selection} is
\[
 C_m=\frac{r_{m+1}-r_m}{1-\rho_m},\qquad
 \frac{C_m}{E_m}
 =\frac{1-E_{m+1}/E_m}{1-\rho_m}\longrightarrow1.
\]
Thus the original rational upper certificate, the selected primitive
errors, and the true interval optima all have the same rationality
dichotomy. Choosing, for example, $\tau_m=1/(m+2)$ in the effective
relative optimizer additionally gives an explicit positive primitive
sequence whose ratio to the true interval optimum tends to one.
Every step terminates using rational comparisons. No uniform running-time
bound and no convergence to zero for gamma are inferred.

\begin{theorem}[Least denominator need not be asymptotically primitive-optimal]
\label{unified-selection-fibonacci-separation}
Let $\varphi=(1+\sqrt5)/2$, $\xi=\varphi^{-1}$, and let
$F_0=0$, $F_1=1$, $F_{k+2}=F_{k+1}+F_k$ be the Fibonacci numbers.
Put
\[
 r_n=\frac{F_{2n}}{F_{2n+1}}\qquad(n\ge1).
\]
For any integers $N>n\ge1$, the unique least-denominator fraction in
$[r_n,r_N]$ is $r_n$, whereas the unique primitive-error minimizer there
is $r_N$. Their exact errors are
\begin{equation}\label{unified-selection-fibonacci-errors}
 F_{2n+1}\xi-F_{2n}=\varphi^{-(2n+1)},\qquad
 F_{2N+1}\xi-F_{2N}=\varphi^{-(2N+1)}.
\end{equation}
Consequently the quotient of the least-denominator error by the true
primitive minimum is exactly
\begin{equation}\label{unified-selection-fibonacci-ratio}
 \varphi^{2(N-n)}.
\end{equation}
In particular, with $n_m=2^m$ and
$I_m=[r_{n_m},r_{n_{m+1}}]$, one has
\[
 \frac{\xi-r_{n_{m+1}}}{\xi-r_{n_m}}\longrightarrow0,
 \qquad
 \frac{\text{least-denominator primitive error on }I_m}
      {\text{least primitive error on }I_m}
 =\varphi^{2^{m+1}}\longrightarrow\infty.
\]
Thus even the stronger contraction used in the canonical gamma
construction does not, by itself, make least-denominator selection
asymptotically primitive-optimal.
\end{theorem}

\begin{proof}
The recurrence and $1-\xi=\xi^2$ give
\begin{equation}\label{unified-selection-fibonacci-identity}
 F_{k+1}\xi-F_k=(-1)^k\xi^{k+1}.
\end{equation}
For example the left side at $k+1$ is $-\xi$ times its value at $k$,
and its value at $k=0$ is $\xi$. This proves the identity by induction.
Consecutive Fibonacci numbers are coprime, so all fractions in the
statement are reduced. The even-index identity puts $r_n$ strictly below
$\xi$, while Cassini's identity gives
\[
 r_{n+1}-r_n=\frac1{F_{2n+1}F_{2n+3}}>0.
\]

Use $r_0=0$ for the next parent computation. The fractions
\[
 \ell_n=\frac{F_{2n-2}}{F_{2n-1}},\qquad
 u_n=\frac{F_{2n-1}}{F_{2n}}
\]
are Farey neighbours, their mediant is $r_n$, and
\[
 \ell_n<r_n\le r_N<\xi<u_n.
\]
The last strict inequality follows from
\eqref{unified-selection-fibonacci-identity} at the odd index $2n-1$.
Every reduced fraction strictly between two Farey neighbours has
denominator at least the sum of their denominators. Equality is possible
only for their mediant: the coefficient representation in the two
unimodular columns consists of positive integers. Here that sum is
$F_{2n+1}$. Since $r_n$ belongs to the target interval, it is the unique
least-denominator fraction there.

For the primitive minimum, the fractions
\[
 \frac ab=r_N=\frac{F_{2N}}{F_{2N+1}},\qquad
 \frac cd=\frac{F_{2N+1}}{F_{2N+2}}
\]
are Farey neighbours with $bc-ad=1$. Their mediant is
$r_{N+1}$, and
\[
 r_N<r_{N+1}<\xi<c/d.
\]
The three-fraction principle of
Theorem~\ref{stage3-farey-three-fractions} therefore gives
\[
 q\xi-p\ge b\xi-a
 \quad\text{for every reduced }p/q<r_{N+1},
\]
with equality only at $p/q=r_N$. The whole interval $[r_n,r_N]$ is
within this range, proving its unique primitive optimum.

Equation~\eqref{unified-selection-fibonacci-identity} now proves
\eqref{unified-selection-fibonacci-errors}, and their quotient is
\eqref{unified-selection-fibonacci-ratio}. For the final interval
sequence, write $n=n_m$ and $N=n_{m+1}$. Its raw-gap ratio is exactly
\[
 \frac{\xi-r_N}{\xi-r_n}
 =\varphi^{-2(N-n)}\frac{F_{2n+1}}{F_{2N+1}}
 <\varphi^{-2(N-n)}\longrightarrow0.
\]
Since $N-n=2^m$, the asserted unbounded error ratio follows.
\end{proof}

\paragraph{What the comparison adds.}
The Fibonacci example does not assert the same behavior for gamma.
It establishes that an asymptotic optimality claim cannot be deduced
merely from rational interval endpoints, strict approach from below,
and a raw-gap ratio tending to zero. Exact primitive optimization can
therefore improve finite errors substantially even when a
least-denominator selector already has the correct limiting dichotomy.
Conversely, the first theorem shows that switching to a relative
optimizer does not eliminate the unresolved rationality branch.
An independent all-index arithmetic estimate is still required to
prove primitive-error decay for Euler's constant.

\section{Primitive obstructions for dyadic targets}

\subsection{Canonical endpoints}

\begin{theorem}[Uniform primitive-cost divergence of the canonical family]
\label{next-cutoff-cost}
Let $r_{m,t}=p_{m,t}/q_{m,t}$ be reduced with $q_{m,t}>0$.
Except for $(m,t)=(0,0)$, set
\[
 N=2^{t+m+1},\quad k=t+m+1,\quad
 B_{m,t}=tP_m(1)+P_m'(1)>0,\quad
 \mathcal P_N=\prod_{N/2<\ell\le N,\ \ell\ {\rm prime}}\ell.
\]
Then
\begin{equation}\label{next-cutoff-actual-denominator}
 \frac{\mathcal P_N}{\gcd(\mathcal P_N,B_{m,t})}
 \mid q_{m,t}
\end{equation}
and
\begin{equation}\label{next-cutoff-cost-bound}
 q_{m,t}(\gamma-r_{m,t})
 \ge\frac{\mathcal P_N}
 {3^{15}k\,2^{2k(k-1)}}.
\end{equation}
In particular the primitive errors tend to $+\infty$ along every
sequence $m+t\to\infty$, although the raw errors tend to zero and the
approximants increase strictly in both coordinates.
\end{theorem}

\begin{proof}
The coefficients of $P_m(1)^2 C_{m,t}$ are integers. Its top
coefficient is
\[
 -4^{m(m+1)/2}\{tP_m(1)+P_m'(1)\}.
\]
For an odd prime $\ell\in(N/2,N]$ not dividing $B_{m,t}$, only the top
harmonic sum $H_N$ contains a denominator divisible by $\ell$, and it
contains precisely the term $1/\ell$. All correction terms $1/(2n)$
have power-of-two denominators. Thus the integer-coefficient harmonic
sum has $\ell$-valuation $-1$. Division by $P_m(1)^2$ can only decrease
that valuation. Every such prime therefore divides the actual reduced
denominator, proving \eqref{next-cutoff-actual-denominator}. Here
$N\ge4$, so all primes in the interval are odd.

Moreover,
\[
 B_{m,t}=(t+s_m)P_m(1)
 \le k\,4^{m(m+1)/2}.
\]
Use this in \eqref{next-cutoff-actual-denominator} and multiply by the
lower bound in \eqref{next-cutoff-two-sided}. Since $a=2^t$,
\[
 q_{m,t}E_{m,t}\ge
 \frac{\mathcal P_N}
 {3^{15}k\,2^{2(m+1)(t+m)}}
 \ge
 \frac{\mathcal P_N}{3^{15}k\,2^{2k(k-1)}}.
\]
The prime number theorem \cite{dlmf-primes} gives
$\log\mathcal P_N=N/2+o(N)$, whereas the logarithm of the remaining
denominator in this bound is $O((\log N)^2)$. This proves the uniform
divergence. All asymptotic estimates here depend only on $N$; no
parameter-dependent short-interval assertion is used.

The conclusion concerns this explicitly specified canonical family.
It does not exclude rational choices of $\kappa$ in the larger
compressed class, different moment prescriptions, or a target-changing
operation with a separately proved primitive-error bound.
\end{proof}

\subsection{Odd-denominator cost of nonperturbative interpolation}
\label{stage3-arithmetic-section}

We examine arbitrary rational interpolation between consecutive canonical
targets. The interpolation coefficient is allowed to have arbitrarily
large numerator and denominator. The result excludes every dyadic rational
interpolation, including choices whose heights grow without any restriction.
For general coefficients it gives an explicit necessary cost in the odd
part of their denominator. It does not exclude arbitrary rational
interpolation.

Recall the canonical polynomials and rational targets
\[
 P_m(x)=\prod_{j=1}^m(4^jx-1),\quad A_m=P_m(1),\quad
 B_m=P_m'(1),\quad s_m=B_m/A_m,
\]
\[
 C_m(x)=\frac{P_m(x)}{A_m}[1-s_m(x-1)]=\sum_i c_{m,i}x^i,
 \qquad
 r_m=\sum_i c_{m,i}f_i,
 \quad f_i=H_{2^i}-\frac1{2^{i+1}}.
\]
Here $P_0=A_0=C_0=1$, $B_0=s_0=0$, and $r_0=1/2$.
The exact positive-kernel identity is
\[
 E_m:=\gamma-r_m
 =\int_0^\infty
 \frac{u^m\{u+4^{m+1}+s_m(4^{m+1}-1)\}}
 {\prod_{i=0}^{m+1}(u+4^i)}\,
 \frac{du}{e^{2\pi\sqrt u}-1}>0.
\]
It also gives the uniform elementary bound
\begin{equation}\label{stage3-arithmetic-raw-lower}
 E_m\ge 3^{-15}2^{-m(m+1)}.
\end{equation}
Indeed, restrict the integral to $1\le u\le2$. Its numerator is at
least $4^{m+1}$, its denominator is less than
$27\,4^{(m+1)(m+2)/2}$, and the measure density is greater than
$3^{-12}$ on that interval.

\begin{theorem}[Exact two-adic arithmetic of the canonical targets]
\label{stage3-arithmetic-two-adic}
For every integer $m\ge0$,
\[
 v_2(r_m)=-1,\qquad v_2(r_{m+1}-r_m)=2m+1.
\]
Let $0<\lambda=u/v\le1$ be reduced, and write $v=2^L H$ with $H$ odd.
If $q$ is the reduced denominator of
$r=r_m+\lambda(r_{m+1}-r_m)$, then
\[
 v_2(q)=
 \begin{cases}
 1,&0\le L\le2m+1,\\
 0,&L=2m+2,\\
 L-2m-1,&L\ge2m+3.
 \end{cases}
\]
In particular $2^L\le2^{2m+2}q$ for every such interpolation coefficient.
\end{theorem}

\begin{proof}
Write $P_m(x)=\sum_i a_i x^i$, taking coefficients outside its support
to be zero. For $0\le i\le m$,
\[
 v_2(a_i)=i(i+1).
\]
Indeed, in the elementary symmetric sum of $4,4^2,\ldots,4^m$, the
unique term of least two-adic valuation is the product $4^1\cdots4^i$.
The number $A_m$ is odd, and for $m\ge1$ the term $a_1$ uniquely has
least valuation in $B_m=\sum_i i a_i$, so $v_2(B_m)=2$.

For $m\ge1$ the coefficients of $C_m$ satisfy
\[
 c_{m,i}=\frac{(A_m+B_m)a_i-B_ma_{i-1}}{A_m^2}.
\]
Consequently $v_2(c_{m,0})=0$, $v_2(c_{m,1})\ge2$, and
\[
 v_2(c_{m,i})=i^2-i+2\quad(2\le i\le m+1).
\]
The unique term of least valuation in $H_{2^i}$ is $1/2^i$, hence
$v_2(f_i)=-i-1$. The $i=0$ term of $r_m$ therefore has valuation $-1$,
the $i=1$ term has valuation at least $0$, and every $i\ge2$ term has
valuation $(i-1)^2\ge1$. This proves $v_2(r_m)=-1$; $m=0$ is direct.

Put $D_m(x)=P_m(x)(x-1)^2/A_m=\sum_i d_i x^i$ and
$Z_m=\sum_i d_i f_i$. From
$A_md_i=a_i-2a_{i-1}+a_{i-2}$ one gets
\[
 v_2(d_0)=0,\quad v_2(d_1)=1,\quad v_2(d_2)=0,
 \qquad v_2(d_i)=(i-2)(i-1)\quad(3\le i\le m+2).
\]
In $Z_m$, the $i=2$ term has valuation $-3$. The terms $i=0,1$ have
valuation $-1$, the term $i=3$ (when present) has valuation $-2$, and
all terms $i\ge4$ have valuation at least $1$. Thus $v_2(Z_m)=-3$.

Set $R=4^{m+1}$. Direct polynomial subtraction gives
\[
 C_{m+1}-C_m=-\kappa_mD_m,\qquad
 \kappa_m=\frac{R\{R+s_m(R-1)\}}{(R-1)^2}.
\]
For $m=0$ this coefficient has valuation $4$. For $m\ge1$,
$v_2(s_m)=2$ and $v_2(R)\ge4$, so $v_2(\kappa_m)=2m+4$.
Therefore $r_{m+1}-r_m=-\kappa_mZ_m$ has valuation $2m+1$.

If $L>0$, reducedness makes $u$ odd, so the added term has valuation
$2m+1-L$. If this is greater or less than $-1$, comparison with
$v_2(r_m)=-1$ gives the displayed denominator exponent. At equality,
the sum of two rational numbers of valuation $-1$ is two-adically integral,
so the denominator is odd. If $L=0$, the added term has valuation at
least $2m+1$, and the first branch again applies.
\end{proof}

\begin{theorem}[A necessary odd-denominator cost for every rational interpolation]
\label{stage3-arithmetic-odd-cost}
For $m\ge0$, put $k=m+1$, $N=2^{m+2}$, and
\[
 \mathcal P_N=\prod_{\substack{N/2<\ell\le N\\\ell\ \mathrm{prime}}}\ell,
 \quad T_m=A_kB_k>0,
 \quad G_m=\frac{\mathcal P_N}{\gcd(\mathcal P_N,T_m)}.
\]
Let $0<\lambda=u/v\le1$ be reduced, let $H$ be the odd part of $v$,
and write $r=r_m+\lambda(r_{m+1}-r_m)=p/q$ in lowest terms.
For every prime $\ell\mid G_m$, one has the exact formula
\[
 v_\ell(q)=\max\{0,1+v_\ell(v)-v_\ell(u)\}.
\]
In particular,
\begin{equation}\label{stage3-arithmetic-prime-divisibility}
 G_m\mid qu,
 \qquad
 q^2H\ge\frac{G_m}{2^{2m+2}}.
\end{equation}
Since
\[
 T_m<(m+2)2^{2(m+1)(m+2)},
\]
the odd denominator must satisfy the completely explicit bound
\begin{equation}\label{stage3-arithmetic-required-odd-height}
 H\ge
 \frac{\mathcal P_N}
 {(m+2)2^{2(m+1)(m+3)}q^2}.
\end{equation}
No smallness assumption on $\lambda$ is used.
\end{theorem}

\begin{proof}
Fix a prime $\ell\mid G_m$. It is odd and does not divide $A_kB_k$.
Since $A_m\mid A_k$, the lower target $r_m$ is $\ell$-integral:
all its harmonic nodes are at most $N/2$, and its coefficient
denominators divide $A_m^2$. The top coefficient of $A_k^2C_k$ is
\[
 -4^{k(k+1)/2}B_k.
\]
It is an $\ell$-unit, and only its harmonic term $H_N$ contains the
single nonintegral summand $1/\ell$. The correction $1/(2N)$ is an
$\ell$-unit. Hence $v_\ell(r_{m+1})=-1$ and
$v_\ell(r_{m+1}-r_m)=-1$.

If $v_\ell(\lambda)<1$, its product with the increment has negative
valuation and dominates the integral base $r_m$. Otherwise both terms
are integral. This proves the exact denominator formula. It implies
$\ell\mid qu$ for every $\ell\mid G_m$, hence $G_m\mid qu$.
The interval condition gives $u\le v=2^LH$, and the preceding theorem
gives $2^L\le2^{2m+2}q$. Combining them yields
$G_m\le qu\le2^{2m+2}q^2H$.

Finally $A_k<2^{k(k+1)}$ and
$B_k=s_kA_k<(k+1)A_k$, because
$s_k=k+\sum_{j=1}^k(4^j-1)^{-1}<k+1$.
This proves the bound on $T_m$ and then
\eqref{stage3-arithmetic-required-odd-height}.
\end{proof}

\begin{theorem}[All dyadic interpolations fail, with unrestricted height]
\label{stage3-arithmetic-dyadic-obstruction}
For the preceding rational interpolation with $0<\lambda\le1$,
\begin{equation}\label{stage3-arithmetic-complete-cost}
 q(\gamma-r)\ge
 \frac{\sqrt{\mathcal P_N}}
 {3^{15}\sqrt{(m+2)H}\,2^{(m+1)(2m+5)}}.
\end{equation}
Consequently, along every sequence $m\to\infty$ for which the odd part
$H=H_m$ of the reduced interpolation denominator satisfies
$\log H_m=o(2^{m+2})$, the primitive errors tend to $+\infty$.
More quantitatively, with $N_m=2^{m+2}$,
\[
 \liminf_{m\to\infty}\frac{\log(q(\gamma-r))}{N_m}
 \ge\frac14-\frac12
       \limsup_{m\to\infty}\frac{\log H_m}{N_m},
\]
whenever the right side is finite. Thus divergence also holds when the
last limsup is strictly less than $1/2$.

In particular every choice of dyadic rational coefficients
$\lambda_m\in[0,1]\cap\mathbb Z[1/2]$ fails the primitive-smallness
condition, even if their heights grow arbitrarily rapidly.
For any rational-interpolation sequence with primitive error tending
to zero, it is necessary that
\[
 \liminf_{m\to\infty}\frac{\log H_m}{2^{m+2}}\ge\frac12.
\]
The conclusion is a necessary condition, not a construction of weights
satisfying it with successful primitive errors.
\end{theorem}

\begin{proof}
The error is the convex combination
$(1-\lambda)E_m+\lambda E_{m+1}$, so
\eqref{stage3-arithmetic-raw-lower} gives
\[
 \gamma-r\ge3^{-15}2^{-(m+1)(m+2)}.
\]
Combine this with
$q\ge\sqrt{G_m/(2^{2m+2}H)}$ and
$G_m\ge\mathcal P_N/T_m$ to obtain
\eqref{stage3-arithmetic-complete-cost}.
The prime number theorem \cite{dlmf-primes} gives
$\log\mathcal P_N=N/2+o(N)$, while $m^2=o(N)$. Taking logarithms proves
the quantitative bound and the stated divergence.

For dyadic coefficients $H=1$. The case $\lambda=0$, omitted from the
divisibility theorem because its reduced numerator is zero, gives the
canonical target $r_m$ itself. For $m\ge1$ this is the preceding
interval's interpolation with coefficient $1$, so the same proved bound
at index $m-1$ already gives divergence. Therefore allowing zero coefficients at any
subsequence does not change the all-dyadic conclusion.

Finally, if the primitive errors tend to zero, then they are eventually
at most $1$. Squaring \eqref{stage3-arithmetic-complete-cost} gives
\[
 H_m\ge
 \frac{\mathcal P_{N_m}}
 {3^{30}(m+2)2^{2(m+1)(2m+5)}}
\]
at all sufficiently large nonzero-interpolation indices. Zero
interpolation cannot occur infinitely often in a successful sequence.
Taking logarithms proves the necessary liminf bound.
\end{proof}

\paragraph{Scope.}
This argument permits completely nonperturbative weights and arbitrary
two-adic denominator growth. It does not restrict the height of the odd
denominator to be small by assumption and then declare all interpolation
impossible: instead it proves how large that odd part must become for a
possible successful sequence. A least-denominator target chosen by a
Stern--Brocot algorithm may require precisely such enormous odd
interpolation denominators, so that selection procedure is not excluded.
The exact arithmetic above does not prove its primitive errors tend to zero.



\section{Constructive extensions of the positive square map}
\label{sea-square-section}

The results in this section are auxiliary tools for choosing polynomial
square forms.  They do not prove the irrationality of Euler's constant,
and no claim of historical novelty is made.  In particular, a degree bound
conditional on separation from a rational number is not a proof of that
separation.

Put $H_0=0$, $H_j=\sum_{h=1}^j1/h$, and
\[
 a_j=j!(1+jH_j),\qquad b_j=jj!.
\]
The moment normalization is fixed by the exact identity
\[
 \int_0^\infty t^j(t-1)\log t\,e^{-t}\,dt
 =\Gamma'(j+2)-\Gamma'(j+1)
 =j!\bigl(1+jH_j-j\gamma\bigr).
\]
For polynomials $f,g\in\mathbb Q[t]$, define the rational bilinear forms
\[
 \mathcal A(f,g)=\sum_{i,j}f_i g_j a_{i+j},\qquad
 \mathcal B(f,g)=\sum_{i,j}f_i g_j b_{i+j},
\]
and write $A(q)=\mathcal A(q,q)$, $B(q)=\mathcal B(q,q)$.  Thus
\begin{equation}
 I(q):=A(q)-\gamma B(q)
 =\int_0^\infty q(t)^2(t-1)\log t\,e^{-t}\,dt>0
 \quad(q\ne0).
 \label{sea-square-positive}
\end{equation}
For integer $q\ne0$, set $g(q)=\gcd(A(q),B(q))>0$ and
$E(q)=I(q)/g(q)$.  For rational $q$, first clear coefficient denominators;
the resulting primitive pair and $E(q)$ are independent of that choice.
For a real parameter $\theta$, write
\[
 L_\theta(q)=A(q)-\theta B(q)
 =\int_0^\infty q(t)^2(t-1)(\log t+\gamma-\theta)e^{-t}\,dt.
\]
We use $1/2<\gamma<7/12$, which follows directly from the positive Binet
integral $\gamma-1/2=\int h_1\,d\mu<\int d\mu=1/12$.
For nonnegative integer $\alpha$, our normalization is
\[
 L_j^{(\alpha)}(t)
   =\sum_{h=0}^j\frac{(-1)^h}{h!}\binom{j+\alpha}{j-h}t^h.
\]

\begin{theorem}[A quantitative polynomial lift of exponential square tests]
\label{sea-square-exponential-lift}
Let $m=2k\ge2$, $\lambda>1$ be rational, and $0\le\theta\le1$.  Define
\[
 \alpha=m-1,\quad \sigma=\frac{\lambda-1}{2},\quad
 \varrho=\frac{\lambda-1}{\lambda+1},\quad
 c=\left(\frac{2}{\lambda+1}\right)^m,
 \qquad f(t)=t^k e^{-\sigma t}.
\]
For $N\ge0$, the polynomial
\[
 q_N(t)=c\,t^k\sum_{j=0}^N\varrho^jL_j^{(\alpha)}(t)
 \in\mathbb Q[t]
\]
has degree $N+k$.  Put
\[
 z=\frac{1+\varrho^2}{2},\qquad \tau=\frac{\varrho^2}{z},\qquad
 \mathcal S_\alpha(z)=
 \left[1+64\left(z\frac{d}{dz}+\alpha+4\right)^3\right]
 (1-z)^{-\alpha-1},
\]
\[
 \mathcal E_N=6c^2\alpha!\,\tau^{N+1}\mathcal S_\alpha(z),\qquad
 \mathcal F=6\left(\frac{(m+2)!}{\lambda^{m+3}}
                    +\frac{(m-1)!}{\lambda^m}\right).
\]
These bounds are positive rational numbers and $\mathcal E_N\to0$.
The exact exponential test is
\begin{equation}
 L_\theta(f)=\frac{m!}{\lambda^{m+2}}
 \left[(m+1-\lambda)(H_m-\log\lambda-\theta)+1\right],
 \label{sea-square-exponential-test}
\end{equation}
and the entire-ray error satisfies
\begin{equation}
 |L_\theta(q_N)-L_\theta(f)|
 \le \mathcal E_N+2\sqrt{\mathcal E_N\mathcal F}.
 \label{sea-square-lift-error}
\end{equation}
Consequently a certified negative exponential test can be converted into a
negative rational polynomial test by a terminating search in $N$.  Clearing
the coefficients then gives an integer polynomial with the same sign.
\end{theorem}

\begin{proof}
The Laguerre generating function and orthogonality \cite{dlmf-laguerre} give
\[
 e^{-\sigma t}=c\sum_{j\ge0}\varrho^jL_j^{(\alpha)}(t),\qquad
 h_j:=\int_0^\infty t^\alpha e^{-t}
                L_j^{(\alpha)}(t)^2\,dt=\frac{(j+\alpha)!}{j!}.
\]
These identities follow, respectively, from
$(1-u)^{-\alpha-1}\exp(-tu/(1-u))$ and the Rodrigues formula.
In the orthonormal Laguerre basis, multiplication by $t$ has diagonal
$2j+\alpha+1$ and off-diagonal
$-\sqrt{(j+1)(j+\alpha+1)}$.  The sum of absolute values in its $j$th row
is at most $4(j+\alpha+1)$.  Since the matrix is tridiagonal, the row sum
for the absolute values of its cube is at most
$64(j+\alpha+4)^3$.  Symmetry and $2|uv|\le u^2+v^2$ therefore imply
\[
 \int_0^\infty t^\alpha(1+t^3)e^{-t}
       \left|\sum_j d_jL_j^{(\alpha)}(t)\right|^2dt
 \le\sum_jh_j[1+64(j+\alpha+4)^3]|d_j|^2.
\]
First this holds for finite sums.  Applying it to differences of finite
tails proves convergence in the displayed weighted norm.  The generating
function identifies that limit with the exponential, so this also justifies
the infinite-tail estimate without a compactness or tail assumption.

For $\theta\in[0,1]$, the signed density factor obeys
\[
 |(t-1)(\log t+\gamma-\theta)|\le W(t):=6(t^2+t^{-1}).
\]
Indeed $|\gamma-\theta|<1$ and $|\log t|\le t+t^{-1}$; bounding separately
on $t\le1$ and $t\ge1$ proves the stated inequality.  With
$\|u\|_W^2=\int W(t)|u(t)|^2e^{-t}dt$, the preceding Laguerre estimate gives
\[
 \|q_N-f\|_W^2
 \le6c^2\alpha!\sum_{j>N}
 \binom{j+\alpha}{\alpha}\varrho^{2j}
       [1+64(j+\alpha+4)^3]
 \le\mathcal E_N.
\]
Here $\varrho^{2j}=\tau^jz^j\le\tau^{N+1}z^j$, and the complete sum in
$z$ equals $\mathcal S_\alpha(z)$.  Also $\|f\|_W^2=\mathcal F$.
Cauchy--Schwarz applied to $(q_N-f)(q_N+f)$ proves
\eqref{sea-square-lift-error}.

Finally,
\[
 \int_0^\infty t^j\log t\,e^{-\lambda t}dt
 =\frac{j!}{\lambda^{j+1}}(H_j-\gamma-\log\lambda).
\]
Substitution and $H_{m+1}=H_m+1/(m+1)$ give
\eqref{sea-square-exponential-test}.  This evaluation also shows explicitly
why its only unevaluated logarithm is the elementary $\log\lambda$;
the eventual polynomial test has exactly the rational coefficients
$A(q_N)-\theta B(q_N)$ when $\theta$ is rational.
\end{proof}

\begin{theorem}[An effective degree bound for exclusion and representation]
\label{sea-square-effective-degree}
Let $0<\varepsilon\le1/3$ and put
\[
 m=2\left\lceil\frac{48}{\varepsilon^2}\right\rceil,
 \qquad N=8m^2,\qquad D=N+\frac m2<77000\varepsilon^{-4}.
\]
For every real $\theta\in[0,1]$ with
$|\theta-\gamma|\ge\varepsilon$, there is a rational polynomial
$q$ of degree at most $D$ such that $L_\theta(q)<0$.
For rational $\theta$ this polynomial has an explicit construction using
only rational arithmetic and certified bounds for $\log m$.

Let
\[
 S_d=\{X\in\mathbb R:[a_{i+j}-Xb_{i+j}]_{0\le i,j\le d}
                     \text{ is positive semidefinite}\}.
\]
Then
\begin{equation}
 S_D\subset(\gamma-\varepsilon,\gamma+\varepsilon).
 \label{sea-square-effective-interval}
\end{equation}
For every rational $r$ with $|r-\gamma|\ge\varepsilon$, there is moreover
an integer polynomial $q\ne0$ of degree at most $D$ such that
$B(q)\ne0$ and $A(q)/B(q)=r$.  A search restricted to this degree therefore
terminates.  No coefficient-height bound is asserted.
\end{theorem}

\begin{proof}
Choose rational $\widetilde h$ with
\[
 |\widetilde h-(H_m-\log m-\theta)|\le1/m,
 \qquad s=\widetilde h/2,\qquad\lambda=m(1+s).
\]
For rational $\theta$ this uses no digits or assumption about $\gamma$.
For example, for $u=(m-1)/(m+1)$ the partial sums of
$2\sum_{j\ge0}u^{2j+1}/(2j+1)$ approximate $\log m$ with remainder at most
$2u^{2J+3}/((2J+3)(1-u^2))$ after index $J$.
The elementary estimate
\[
 0<H_m-\log m-\gamma<1/m
\]
follows by summing
$0<\log(1+1/j)-1/(j+1)<1/(j(j+1))$ for $j\ge m$.
Consequently $|s|\le1/3$ and $|s|\ge\varepsilon/4$.
In particular $2m/3\le\lambda\le4m/3$ and $\lambda>1$.

Write $h=H_m-\log m-\theta$ and $e=h-2s$, so $|e|\le1/m$.
On $[-1/3,1/3]$ the function $g(s)=2s-\log(1+s)$ satisfies
\[
 s g(s)\ge\tfrac12s^2,\qquad |g(s)|\le\tfrac54|s|,
\]
as follows from $1/2\le g'(s)\le5/4$.  The bracket in
\eqref{sea-square-exponential-test} is thus
\[
 1+(1-ms)(g(s)+e)
 \le1+\tfrac94|s|+1/m-\tfrac m2s^2
 <2-\tfrac m2s^2\le-1.
\]
It follows that $L_\theta(f)\le-\eta$, where
$\eta=m!/\lambda^{m+2}>0$ is rational.

We record estimates establishing a fully explicit truncation order.
With the notation of Theorem~\ref{sea-square-exponential-lift}, $m\ge96$
and elementary binomial generating functions give
\[
 \mathcal S_{m-1}(z)\le257m^3(1-z)^{-m-3}.
\]
For clarity, use
$(j+m+3)^3\le2(j+m)^3$ for $m\ge32$ and
\[
 (j+m)^3\binom{j+m-1}{m-1}
 \le m(m+1)(m+2)\binom{j+m+2}{m+2},
\]
then $m(m+1)(m+2)\le2m^3$.
Since
\[
 1-z=\frac{2\lambda}{(\lambda+1)^2},\qquad
 \tau=1-\frac{2\lambda}{\lambda^2+1},
\]
one obtains
\[
 \frac{\mathcal E_N}{\eta}
 <60000m^7 2^m\exp\left(-\frac{3(N+1)}{4m}\right),
 \qquad \frac{\mathcal F}{\eta}\le24m.
\]
For the first inequality, use
\[
 6c^2(m-1)!\mathcal S_{m-1}(z)
 <30000m^6(m-1)!(2/\lambda)^m
\]
and $\tau^{N+1}\le\exp(-3(N+1)/(4m))$.
For $N=8m^2$,
\[
 \mathcal E_N<\frac{\eta}{1536m}.
\]
Indeed the logarithm of $92160000m^8 2^m e^{-6m}$ is negative for $m\ge96$;
the inequalities
\[
 \log(92160000)<27,\qquad 8\log m\le2m,\qquad m\log2<m
\]
suffice.
Thus
\[
 \mathcal E_N+2\sqrt{\mathcal E_N\mathcal F}<\eta/2,
\]
and the constructed $q_N$ has $L_\theta(q_N)<-\eta/2$.

Each $S_d$ is a convex interval, contains $\gamma$, and the intervals are
nested as $d$ grows.  The tests at $\theta=0,1$ exclude both endpoints,
since their distances from $\gamma$ exceed $1/3$.  Therefore $S_D\subset(0,1)$;
the other tests give \eqref{sea-square-effective-interval}.
The degree estimate follows from $m\le98\varepsilon^{-2}$ and
$8\cdot98^2+49<77000$.

For rational $r\in[0,1]$ at the stated distance, the rational quadratic
form $L_r$ on polynomials of degree at most $D$ has a negative vector and
the positive vector $1$.  For $r<0$, use the test negative at $0$:
its $B$ is negative by \eqref{sea-square-positive}, so it remains negative
at every $r<0$.  Likewise the test negative at $1$ remains negative for
$r>1$.  Thus the form is indefinite in every case.
Its dimension is at least five.  If singular, a nonzero rational kernel
vector is isotropic.  Otherwise Meyer's theorem for indefinite rational
quadratic forms in at least five variables supplies a rational isotropic
vector.  Clearing denominators gives $q\in\mathbb Z[t]\setminus\{0\}$.
Equation~\eqref{sea-square-positive} forces $B(q)\ne0$, and hence
$A(q)/B(q)=r$.  Enumeration of integer coefficient vectors of the bounded
degree is a terminating exact algorithm; no bound for its running time
or the coefficient height is implied.
\end{proof}

\paragraph{A quantitative enclosure and its scope.}
For every integer $d\ge6237000$, the same argument yields the explicit
enclosure
\[
 S_d\subset\left(\gamma-(77000/d)^{1/4},
                 \gamma+(77000/d)^{1/4}\right).
\]
The condition $|r-\gamma|\ge\varepsilon$ remains essential in the
representation theorem.  Applying it to an arbitrary rational $r$ without
first establishing that condition would assume the desired exclusion.


\section{Exact arithmetic in Laguerre coordinates}
\label{next-square-section}

This section makes the coefficient arithmetic of the positive square map
explicit in the ordinary Laguerre basis. It provides an optimal universal
normalization, a local prime-residue rule, and a full primitive-cost
classification for fixed polynomial coefficient schedules. It does not
exclude arbitrary growing sequences of coefficient vectors.

Let
\[
 L_j(t)=\sum_{h=0}^j\frac{(-1)^h}{h!}\binom jh t^h,\qquad
 H_0=0,\quad H_j=\sum_{h=1}^j\frac1h.
\]
For a rational polynomial $Q$, define the rational numbers $A(Q),B(Q)$ by
\[
 A(Q)-\gamma B(Q)
 =I(Q):=\int_0^\infty Q(t)^2(t-1)\log t\,e^{-t}\,dt,\qquad
 B(Q)=\int_0^\infty Q(t)^2(t-1)e^{-t}\,dt.
\]
Their rationality also follows from the matrices below. For every nonzero
$Q$, $I(Q)>0$. Clear denominators in $(A(Q),B(Q))$ by any positive integer
$D$ and let $g=\gcd(DA(Q),DB(Q))$. The primitive error is
$E(Q)=DI(Q)/g$, independent of the choice of $D$ or nonzero rational
scaling of $Q$. If $B(Q)\ne0$ and $q$ is the reduced denominator of
$A(Q)/B(Q)$, then
\begin{equation}
 E(Q)=q\,\frac{I(Q)}{|B(Q)|}.
 \label{next-square-primitive-error}
\end{equation}
If $B(Q)=0$, then $E(Q)=1$.

\begin{theorem}[Optimal half-lcm normalization in Laguerre coordinates]
\label{next-square-normalization}
Let $Q_c=\sum_{j=0}^n c_jL_j$ with $c_j\in\mathbb Z$. Then $B(Q_c)$
is an integer. For $n\ge2$, the smallest positive integer $D_n$ such that
$D_nA(Q_c)\in\mathbb Z$ for every $c\in\mathbb Z^{n+1}$ is
\[
 D_n=\frac{\operatorname{lcm}(1,\ldots,n)}2.
\]
For $n=0,1$ it is $D_n=1$. The exact symmetric Gram matrices are
\begin{align}
 B_{ii}&=2i,& B_{i,i+1}&=-(i+1),&
 B_{ij}&=0\quad(|i-j|\ge2),\label{next-square-b-gram}\\
 A_{ii}&=2iH_i+2i+1,\qquad&
 A_{i,i+1}&=-(i+1)H_i-\frac{3i}{2}-1,\label{next-square-a-near}\\
 A_{ij}&=\frac{i+j+(j-i)^2}{(j-i)((j-i)^2-1)}
 &&(j-i\ge2).\label{next-square-a-far}
\end{align}
Thus $(D_nA(Q_c),D_nB(Q_c))$ is an exact integer coefficient pair,
without requiring factorial clearing of the polynomial coefficients.
Its gcd must still be removed to obtain its primitive pair.
\end{theorem}

\begin{proof}
The Laguerre recurrence is
\[
 tL_i=(2i+1)L_i-iL_{i-1}-(i+1)L_{i+1},
\]
and the $L_i$ are orthonormal for $e^{-t}dt$. This gives
\eqref{next-square-b-gram}. The logarithmic Gram matrix is
\[
 \int_0^\infty \log t\,L_i(t)L_j(t)e^{-t}\,dt
 =
 \begin{cases}
 H_i-\gamma,&i=j,\\
 -1/|i-j|,&i\ne j.
 \end{cases}
\]
One proof differentiates generalized Laguerre orthogonality at parameter
zero. The norm derivative is $H_i-\gamma$, and
\[
 \left.\partial_\alpha L_i^{(\alpha)}\right|_{\alpha=0}
 =\sum_{j<i}\frac{L_j}{i-j}.
\]
The last identity follows by differentiating the Laguerre generating
function; see \cite{dlmf-laguerre}. All expressions are finite polynomial
sums and the differentiated integrals have an integrable majorant.
Multiplying this logarithmic matrix by the recurrence for $t-1$, including
the $L_{n+1}$ term when necessary, gives
\eqref{next-square-a-near}--\eqref{next-square-a-far}.

Write $d_n=\operatorname{lcm}(1,\ldots,n)$. A quadratic form is
integer-valued on all integer vectors exactly when its diagonal
coefficients and twice its off-diagonal coefficients are integers.
For $k=j-i\ge2$, the useful partial-fraction identity is
\[
 2A_{ij}=\frac{2(i+1)}{k-1}-\frac{4i}{k}+\frac{2i}{k+1}.
\]
If $i\ge1$, all displayed denominators are at most $n$. If $i=0$, only
$2/(k-1)$ remains. Together with \eqref{next-square-a-near}, this shows
that $d_n/2$ clears the quadratic values.

For sharpness, testing the vectors $e_i$ and $e_i+e_j$ forces every
denominator in $A_{ii}$ and $2A_{ij}$. Row zero supplies
$2/(k-1)$ for $2\le k\le n$, so for $n\ge3$ it forces
$d_{n-1}/2$. Induct on $n$, starting at $n=2$. A new factor is needed
only when $n$ is a prime power. For $n\ge3$,
\[
 2A_{1n}=\frac{2[n(n-1)+2]}{(n-1)n(n-2)}.
\]
If $n=p^a$ for an odd prime $p$, this has $p$-adic valuation $-a$.
If $n=2^a$ with $a\ge2$, its valuation at $2$ is $1-a$. These are
exactly the new prime-power factors of $d_n/2$, proving optimality.
\end{proof}

For finite vectors $c,d$, pad both by zero and write
\[
 \mathcal B(c,d)=
 \sum_{r\ge0}2r\,c_rd_r
 -\sum_{r\ge0}(r+1)(c_rd_{r+1}+c_{r+1}d_r).
\]
Thus $B(Q_c)=\mathcal B(c,c)$.

\begin{theorem}[Exact residue at a prime in the upper half]
\label{next-square-prime-block}
Let $p$ be an odd prime with $n/2<p\le n$, and let $s=n-p$. For an integer
coefficient vector $c=(c_0,\ldots,c_n)$ put
\[
 d_r=c_{p+r}\ (0\le r\le s),\qquad
 h_r=c_r\ (0\le r\le s+1),
\]
and pad these vectors by zero. Then $pA(Q_c)$ is $p$-integral and
\begin{equation}
 pA(Q_c)\equiv
 \mathcal B(d,d)-2\mathcal B(h,d)\pmod p.
 \label{next-square-prime-residue}
\end{equation}
If this residue is nonzero and $B(Q_c)\ne0$, the prime $p$ divides the
actual reduced denominator of $A(Q_c)/B(Q_c)$.
\end{theorem}

\begin{proof}
The normalization theorem shows that $v_p(A(Q_c))\ge-1$.
In the diagonal terms, the residue of $pA_{ii}$ is $2i$ when $i\ge p$,
and zero otherwise. In the adjacent cross terms, the residue of
$2pA_{i,i+1}$ is $-2(i+1)$ when $i\ge p$, and zero otherwise.
These give $\mathcal B(d,d)$.

For a gap $k=j-i\ge2$, the only potentially nonzero residues of
$2pA_{ij}$ occur at $k=p-1,p,p+1$. They are, respectively,
\[
 2i,\qquad -4i,\qquad 2(i+1).
\]
The possible exceptional gap $2p-1$ has $i=0$ and residue zero.
Collecting the three contributing diagonals gives
\[
 2\sum_{r=0}^{s}(r+1)h_{r+1}d_r
 -4\sum_{r=0}^{s}r h_rd_r
 +2\sum_{r=1}^{s}r h_{r-1}d_r
 =-2\mathcal B(h,d).
\]
At $n=2p-1$ the apparent overlapping endpoint has coefficient $s+1=p$
and vanishes modulo $p$, so the same formula applies.

If the residue is nonzero, $v_p(A(Q_c))=-1$. Since $B(Q_c)$ is an
integer, $v_p(A(Q_c)/B(Q_c))=-1-v_p(B(Q_c))<0$. This proves survival
in the reduced denominator, not merely in a chosen clearing factor.
\end{proof}

\begin{theorem}[Complete primitive-cost classification of fixed polynomial schedules]
\label{next-square-polynomial-filter}
Fix a nonzero polynomial $P(X,Y)\in\mathbb Q[X,Y]$ and put
\[
 Q_N(t)=\sum_{k=0}^N P(N,k)L_k(t).
\]
These polynomials are nonzero for all sufficiently large $N$.
The quantity $B_N=B(Q_N)$ is a rational polynomial in $N$.
Exactly one of the following holds:
\begin{enumerate}
\item $B_N$ is identically zero. Then $E(Q_N)=1$ for all sufficiently
large $N$.
\item $B_N$ is not identically zero. Then it is eventually nonzero and
\[
 \liminf_{N\to\infty}\frac{\log E(Q_N)}{N}\ge\frac12.
\]
\end{enumerate}
In particular, no fixed polynomial coefficient schedule of this form
has primitive error tending to zero.
\end{theorem}

\begin{proof}
Clear the fixed coefficient denominators of $P$ once. This scales each
$Q_N$ by the same nonzero rational factor and leaves its primitive error
unchanged, so assume $P\in\mathbb Z[X,Y]$.

Write $c_k=P(N,k)$ and define the integer-valued rational polynomial
\begin{equation}
 F_P(N)=2(N+1)P(N,N)P(N,N+1)-B_N.
 \label{next-square-F-polynomial}
\end{equation}
Both $B_N$ and $F_P(N)$ are polynomials because sums of polynomial
functions over $k=0,\ldots,N$ are polynomials in $N$. We first prove
that $F_P$ is nonzero; it is essential not to infer this merely from
nonconstancy in the second variable.

Let $d$ be the total degree of $P$. Its highest homogeneous part gives
a nonzero real polynomial $f$ such that, uniformly for $x\in[0,1]$,
\[
 P(N,Nx)=N^d f(x)+O_P(N^{d-1}).
\]
The exact discrete energy identity is
\begin{equation}
 B_N=(N+1)c_N^2-\sum_{k=0}^N c_k^2
       +\sum_{k=0}^{N-1}(k+1)(c_k-c_{k+1})^2.
 \label{next-square-energy}
\end{equation}
The last sum is $O_P(N^{2d})$, while the middle sum is a polynomial
Riemann sum. Hence
\[
 B_N=N^{2d+1}
 \left[f(1)^2-\int_0^1 f(x)^2\,dx\right]+O_P(N^{2d}),
\]
and therefore
\begin{equation}
 F_P(N)=N^{2d+1}
 \left[f(1)^2+\int_0^1 f(x)^2\,dx\right]+O_P(N^{2d}).
 \label{next-square-F-leading}
\end{equation}
The coefficient in brackets is strictly positive. Thus $F_P$ has
degree exactly $2d+1$ and is nonzero for all sufficiently large $N$.

Put $r=\deg_Y P$. Let $p$ be a prime with
\[
 N/2<p\le N,\qquad p>2r+2,
\]
and let $s=N-p$. In \eqref{next-square-prime-residue},
$P(N,p+j)\equiv P(N,j)\pmod p$. Thus the residue becomes
\[
 pA(Q_N)\equiv
 2(s+1)P(N,s)P(N,s+1)-B_s(P(N,\cdot))\pmod p.
\]
For fixed $N$ the right side is a rational polynomial in $s$. Its
coefficient denominators have no prime factors above $2r+2$:
expand the polynomial summands of degree at most $2r+1$ in the Newton
binomial basis and sum them. Since $s\equiv N\pmod p$, we may replace
$s$ by $N$ and conclude
\begin{equation}
 pA(Q_N)\equiv F_P(N)\pmod p.
 \label{next-square-poly-prime}
\end{equation}
When $B_N\ne0$, every prime in $(N/2,N]$ not dividing $F_P(N)$ therefore
survives in the reduced denominator $q_N$ of $A(Q_N)/B_N$.
For all sufficiently large $N$ the fixed lower restriction $p>2r+2$
is automatic. Since the product of the distinct excluded primes divides
the nonzero integer $F_P(N)$,
\[
 \log q_N\ge
 \vartheta(N)-\vartheta(N/2)-\log|F_P(N)|
 =\left(\frac12+o(1)\right)N.
\]
Here $\vartheta(x)=\sum_{p\le x}\log p$ and
$\vartheta(x)\sim x$ is the prime number theorem
\cite{dlmf-primes}.

We must also control the real error; denominator growth alone would
not suffice. If $Q$ has degree $j$ and highest Laguerre coefficient
$c_j$, then
\begin{equation}
 I(Q)\ge e^{-2}\log2\,(j+1)c_j^2.
 \label{next-square-real-lower}
\end{equation}
Indeed, restrict the positive integral to $t\ge2$ and use
$(t-1)\log t\ge(t-2)\log2$. After $t=x+2$, the leading monomial
coefficient is $(-1)^jc_j/j!$. The monic degree-$j$ orthogonal
polynomial for $xe^{-x}dx$ has squared norm $j!(j+1)!$, so the
orthogonal decomposition of $Q(x+2)$ gives
\[
 I(Q)\ge e^{-2}\log2
 \left(\frac{c_j}{j!}\right)^2j!(j+1)!,
\]
which is \eqref{next-square-real-lower}. The norm identity follows
directly from the generalized Laguerre Rodrigues formula by $j$
integrations by parts.

For all sufficiently large $N$, $P(N,Y)$ is a nonzero polynomial of
degree at most $r$. It cannot vanish at the $r+1$ distinct integers
$N,N-1,\ldots,N-r$. Consequently $\deg Q_N\ge N-r$, and its highest
Laguerre coefficient is a nonzero integer. Thus $I(Q_N)$ has a positive
linear lower bound. Since the nonzero polynomial $B_N$ has polynomial
size, $I(Q_N)/|B_N|$ is bounded below by an inverse power of $N$.
Equation \eqref{next-square-primitive-error} now proves the exponential
primitive-error lower growth.

If $B_N$ is identically zero, positivity gives $A(Q_N)>0$ whenever
$Q_N\ne0$. Primitive normalization of the rational pair $(A(Q_N),0)$
then gives error exactly one.
\end{proof}

\paragraph{Scope and an instructive zero-coefficient case.}
Nonconstant coefficient profiles can satisfy $B_N\equiv0$. For example,
$P(N,k)=3k-(2N+1)$ has this property; the first branch of the theorem
includes it. The proof of \eqref{next-square-F-leading} remains valid
in this case and does not assume that $B_N$ is nonzero.

Let
\[
 \mathcal D_t=t\frac{d^2}{dt^2}+(1-t)\frac d{dt},
 \qquad \mathcal D_tL_k=-kL_k.
\]
The theorem covers every fixed polynomial filter
$P(N,-\mathcal D_t)$ applied to the zero-location kernel
$\sum_{k=0}^N L_k(t)$. It therefore covers polynomial tapers of any
fixed order and every fixed polynomial dependence on $N$.
Coefficients such as $L_k(x_N)$ at nonzero $x_N$, filters whose degree
grows with $N$, and arbitrary coefficient sequences are outside its
scope.


\section{Shared-cutoff support compression}

\begin{theorem}[Compression at a shared cutoff and its primitive cost]
\label{sea-arith-cutoff-compression}
For $m\ge0$ define
\[
 T_m(x)=(x-1)^2(2x-1)\prod_{k=1}^m(4^kx-1),\qquad
 d=m+3,\qquad s=1+m(m+1).
\]
For every $C\in\mathbb Z[x]$, division gives unique
$S,R\in\mathbb Z[1/2][x]$ with
\[
 C=T_mS+R,\qquad \deg R<d.
\]
If $\deg C=L\ge d$, all coefficient denominators of $S$ and $R$ divide
$2^{s(L-d+1)}$. Moreover,
\begin{align}
 R(1)&=C(1),& R'(1)&=C'(1),& R(1/2)&=C(1/2),\notag\\
 R(4^{-k})&=C(4^{-k})&&&&(1\le k\le m).
 \label{sea-arith-compression-invariants}
\end{align}

Let $a=2^t$, where $t\ge0$ is an integer, and assume
\begin{equation}
 Q=C(1)\ne0,\qquad C'(1)+tQ=0.
 \label{sea-arith-whole-log-condition}
\end{equation}
For $C(x)=\sum c_ix^i$ and $R(x)=\sum r_ix^i$, define the corrected
harmonic targets
\begin{equation}
 r_C=\frac1Q\sum_i c_i\left(H_{a2^i}-\frac1{2a2^i}\right),\qquad
 r_R=\frac1Q\sum_i r_i\left(H_{a2^i}-\frac1{2a2^i}\right).
 \label{sea-arith-corrected-targets}
\end{equation}
Both give log-free gamma forms. Their target difference is generally
nonzero, but it satisfies the exact identity and bound
\begin{align}
 \Delta:=r_C-r_R
   &=\frac1Q\sum_i[x^i](T_mS)H_{a2^i},\notag\\
 |\Delta|&\le\Delta_{\max}:=
 \frac{\mu_m a^{-2m-2}}{|Q|}
 T_m^{\mathrm{abs}}(x_0)S^{\mathrm{abs}}(x_0),\qquad
 x_0=4^{-m-1},
 \label{sea-arith-compression-delta}
\end{align}
where $P^{\mathrm{abs}}(x)=\sum_i|[x^i]P|x^i$.

Let $E$ be the least nonnegative integer with $2^ER\in\mathbb Z[x]$,
and put
\[
 N_0=a2^{d-1},\qquad D_0=2\operatorname{lcm}(1,\ldots,N_0),\qquad
 Z=2^ED_0\sum_i r_i\left(H_{a2^i}-\frac1{2a2^i}\right)\in\mathbb Z.
\]
The complete reduced denominator of $r_R$ is exactly
\begin{equation}
 q_R=\frac{|Q|2^ED_0}{\gcd(|Z|,|Q|2^ED_0)}.
 \label{sea-arith-compressed-final-q}
\end{equation}
In particular,
\[
 \operatorname{odd}(q_R)\mid
 \operatorname{odd}(Q)\operatorname{odd}
        \bigl(\operatorname{lcm}(1,\ldots,N_0)\bigr).
\]
Every odd prime above $N_0$ not dividing $Q$ is therefore absent after
compression, even if it appeared in the original target $r_C$.

If independently certified rational bounds satisfy
\[
 0<\ell\le\gamma-r_C\le u,\qquad \Delta_{\max}<\ell,
\]
then, with $P_R=q_Rr_R\in\mathbb Z$,
\begin{equation}
 0<\ell-\Delta_{\max}\le\gamma-r_R\le u+\Delta_{\max},\qquad
 0<q_R\gamma-P_R\le q_R(u+\Delta_{\max}).
 \label{sea-arith-compressed-error-transfer}
\end{equation}
This is a checkable transfer of a positive error certificate. It is not
an assertion that the last upper bound tends to zero.
\end{theorem}

\begin{proof}
The leading coefficient of $T_m$ is
\[
 2\prod_{k=1}^m4^k=2^{1+m(m+1)}=2^s.
\]
Long division therefore introduces only powers of $2$ in denominators.
Each of the at most $L-d+1$ steps increases the necessary power by at
most $s$, giving the stated bound. The double root at $1$ and the other
listed roots give \eqref{sea-arith-compression-invariants}. In particular
the entire condition \eqref{sea-arith-whole-log-condition} is preserved.
It is essential to include $tQ$: for $a>1$, cancellation of $C'(1)$
alone would leave $Q\log a$.

Subtract the corrected targets. Since $(T_mS)(1/2)=0$, the inverse first
moment vanishes, and the $1/(2n)$ terms cancel in their difference.
Also $(T_mS)(1)=(T_mS)'(1)=0$, so all gamma and logarithmic terms cancel
when \eqref{stage3-binet} is summed over its coefficients.
For $z>0$, the exact finite expansion is
\[
 \Phi(z)=\sum_{k=0}^{m-1}(-1)^k\mu_kz^{-k-1}
    +(-1)^mz^{-m}\int\frac{u^m}{u+z}\,d\mu(u).
\]
For $m=0$ the sum is empty. The remainder has absolute value at most
$\mu_mz^{-m-1}$. Each of the first $m$ terms vanishes in the coefficient
sum because $(T_mS)(4^{-k})=0$ for $1\le k\le m$. Thus
\[
 |\Delta|\le
 \frac{\mu_ma^{-2m-2}}{|Q|}(T_mS)^{\mathrm{abs}}(x_0)
 \le\frac{\mu_ma^{-2m-2}}{|Q|}
        T_m^{\mathrm{abs}}(x_0)S^{\mathrm{abs}}(x_0),
\]
which is \eqref{sea-arith-compression-delta}.

All surviving indices satisfy $a2^i\le N_0$. Each corrected harmonic
denominator divides $D_0$, and $2^E$ clears all $r_i$. Hence $Z$ is an
integer, and reduction of $Z/(Q2^ED_0)$ gives
\eqref{sea-arith-compressed-final-q}. The odd-part assertion follows
immediately. Finally
$\gamma-r_R=(\gamma-r_C)+\Delta$, so
\eqref{sea-arith-compressed-error-transfer} follows from the triangle
inequality and the exact value of $q_R$.
\end{proof}

\paragraph{An infinite shared-prime cancellation with positive errors.}
Set the external base $a=1$, take $m=0$, and put
\[
 C_t(x)=1+x^tT_0(x),\qquad t\ge1.
\]
The remainder is $R=1$. Its corrected target is exactly
$H_1-1/2=1/2$, with final denominator $2$. The original target is
\[
 r_{C_t}=\frac12+V_t,\qquad
 V_t=\sum_{i=0}^3[x^i]T_0(x)H_{2^{t+i}}.
\]
Since $T_0^{\mathrm{abs}}(1/4)=75/32$, the finite Binet remainder gives
\[
 |V_t|\le\frac{25}{128\,4^t}<\frac3{40}.
\]
On the other hand $h_1(u)>1-u$ for $u>0$, so
\[
 \gamma-\frac12>\mu_0-\mu_1=\frac3{40}.
\]
Every original target and every compressed target therefore has a
strictly positive gamma error.

For any prime $\ell$ in $(2^{t+2},2^{t+3}]$, only the last harmonic sum
in $V_t$ contains a denominator divisible by $\ell$, and its coefficient
is $2$. It contains exactly the single term $1/\ell$ with such a
denominator. Thus $v_\ell(r_{C_t})=-1$. All these primes disappear after
compression. For instance the exact reduced denominators for $t=1,2$
are respectively $90090$ and $9025222055850$ before compression, and
$2$ afterward.

This example explicitly removes a block of primes already present in
the base. It lies outside an obstruction whose hypothesis requires that
the base avoid that block. At the same time, it changes the target to
$1/2$, leaving the error $\gamma-1/2$ fixed. It therefore proves the
arithmetic cancellation mechanism without claiming the missing
small-error estimate needed for irrationality.


\section{Rational localization at the double zero of the positive weight}
\label{stage3-localization-section}

The positive square weight $(t-1)\log t$ vanishes quadratically at $t=1$.
This section constructs rational polynomials localized at points tending
to that zero. The coefficients depend on Laguerre values and lie outside
the fixed-polynomial-schedule obstruction. The real integral becomes
small with an explicit order, and its full primitive cost is expressed
in terms of the actual reduced denominator. A one-step change of the
cutoff also removes a specific adverse prime factor exactly.

Let $L_k^{(1)}$ denote the generalized Laguerre polynomial with parameter
$1$, and let $L_k=L_k^{(0)}$. For an integer $M\ge2$ and
$\sigma\in\{0,-1\}$, define
\[
 N=M^2+\sigma,\qquad x=1-\frac1M,\qquad H=N+1,\qquad
 w_k=\left(1-\frac{k}{H}\right)^2\quad(0\le k\le N),
\]
\begin{equation}
 p_{N,x}(t)=
 \sum_{k=0}^N\frac{w_kL_k^{(1)}(x)L_k^{(1)}(t)}{k+1},
 \qquad Q_M(t)=t\,p_{N,x}(t).
 \label{stage3-localization-family}
\end{equation}
These polynomials have rational coefficients, constructed without using
a value or an assumption about the arithmetic nature of $\gamma$.
Write
\[
 I_M=A_M-\gamma B_M
 =\int_0^\infty Q_M(t)^2(t-1)\log t\,e^{-t}dt>0.
\]
The positivity is strict because $Q_M$ is nonzero. Indeed its ordinary
Laguerre coefficient at index zero is $w_0L_0^{(1)}(x)=1$.

\begin{theorem}[Exact rational certificates for the localized family]
\label{stage3-localization-sandwich}
Use the inner product
$\langle f,g\rangle_1=\int_0^\infty f(t)g(t)t e^{-t}dt$ and put
\[
 H_j=\langle p_{N,x},t^j p_{N,x}\rangle_1,\qquad
 T_M=H_2+H_1,\qquad U_M=H_3-H_2-H_1+H_0.
\]
Then all these numbers are rational, $B_M=H_2-H_1$, and
\begin{equation}
 \frac{2B_M^2}{T_M}\le I_M\le\frac{U_M}{2}.
 \label{stage3-localization-rational-sandwich}
\end{equation}
When $B_M\ne0$, let $q_M$ be the actual reduced denominator of $A_M/B_M$.
The complete primitive error $\mathcal E_M$ satisfies the entirely
rational certificate
\begin{equation}
 \frac{2q_M|B_M|}{T_M}
 \le\mathcal E_M
 \le\frac{q_MU_M}{2|B_M|}.
 \label{stage3-localization-primitive-sandwich}
\end{equation}

For exact integerization, put
\[
 F_M=H^2N!M^N,\qquad
 D_N=\frac{\operatorname{lcm}(1,\ldots,N+1)}2,\qquad
 S_M=D_NF_M^2.
\]
Then $S_MA_M,S_MB_M$ are integers. If
$g_M=\gcd(S_MA_M,S_MB_M)>0$, the identities
\[
 \mathcal E_M=\frac{S_MI_M}{g_M},\qquad
 q_M=\frac{S_M|B_M|}{g_M}\quad(B_M\ne0)
\]
use the complete gcd, not merely a chosen clearing factor.
\end{theorem}

\begin{proof}
For every $t>0$,
\[
 \frac{2(t-1)^2}{t+1}
 \le(t-1)\log t
 \le\frac{(t-1)^2(t+1)}{2t}.
\]
For the lower inequality use
$\log t=2\operatorname{artanh}((t-1)/(t+1))$.
For the upper one, the derivative of
$(t-t^{-1})/2-\log t$ is $(t-1)^2/(2t^2)\ge0$, and its value at $1$
is zero. Multiplication by $Q_M(t)^2=t^2p_{N,x}(t)^2$ and integration
gives the upper bound $U_M/2$. Cauchy--Schwarz gives
\[
 B_M^2
 \le
 \left(\int_0^\infty
       \frac{t(t-1)^2}{t+1}p_{N,x}(t)^2t e^{-t}dt\right)
 \left(\int_0^\infty t(t+1)p_{N,x}(t)^2t e^{-t}dt\right)
 \le\frac{I_MT_M}{2}.
\]
This proves \eqref{stage3-localization-rational-sandwich}.
Since $I_M>0$,
\[
 \mathcal E_M=q_M\left|\gamma-\frac{A_M}{B_M}\right|
             =q_M\frac{I_M}{|B_M|},
\]
which gives \eqref{stage3-localization-primitive-sandwich}.

The elementary identity
\[
 tL_k^{(1)}(t)=(k+1)(L_k(t)-L_{k+1}(t))
\]
shows that, with $z_k=w_kL_k^{(1)}(x)$,
\[
 Q_M=\sum_{k=0}^N z_k(L_k-L_{k+1}).
\]
Its ordinary Laguerre coefficient vector is
$(z_0,z_1-z_0,\ldots,z_N-z_{N-1},-z_N)$.
The finite coefficient formula for $L_k^{(1)}(x)$ shows that $F_M$
clears this vector. Lemma~\ref{next-square-normalization}, applied to an integer
Laguerre vector of length $N+2$, proves that $S_M$ clears $(A_M,B_M)$.
The remaining identities follow by complete reduction of this pair.
\end{proof}

\begin{theorem}[Uniform localization and its exact arithmetic threshold]
\label{stage3-localization-asymptotic}
For either $\sigma=0$ or $\sigma=-1$ in
\eqref{stage3-localization-family}, let
\[
 \kappa=\frac{128e}{315\pi}>0.
\]
Then, as $M\to\infty$,
\[
 B_M=-\kappa+O\left(\frac{\log M}{M}\right),\qquad
 T_M=2\kappa M+O(\log M),\qquad
 U_M=O(M^{-1}).
\]
Consequently
\[
 I_M=\Theta(M^{-1}),\qquad
 \gamma-\frac{A_M}{B_M}=\Theta(M^{-1})
\]
with the second quantity positive for all sufficiently large $M$.
For the actual reduced denominator,
\begin{equation}
 \mathcal E_M=\Theta\!\left(\frac{q_M}{M}\right).
 \label{stage3-localization-threshold}
\end{equation}
Thus this particular family would satisfy primitive-error decay exactly
if $q_M=o(M)$; equivalently, $Mg_M/S_M\to\infty$.
Neither denominator condition is assumed or proved here.
\end{theorem}

\begin{proof}
All estimates below are uniform for $x$ in the compact interval $[1/2,1]$.
We use the established Laguerre asymptotic
\cite{dlmf-laguerre-asymptotic}
\begin{equation}
 L_k^{(\alpha)}(x)=
 \frac{e^{x/2}k^{\alpha/2-1/4}}
      {\sqrt{\pi}\,x^{\alpha/2+1/4}}
 \left[
 \cos\!\left(2\sqrt{kx}
       -\left(\frac{\alpha}{2}+\frac14\right)\pi\right)
 +O(k^{-1/2})\right]
 \label{stage3-localization-PR}
\end{equation}
for $\alpha=0,1$. Its compact-$x$ uniformity is essential because
$x=1-1/M$ varies with $M$.

Let $\phi_k=L_k^{(1)}/\sqrt{k+1}$, an orthonormal basis for the inner
product $\langle\ ,\ \rangle_1$. Multiplication by $t$ has Jacobi matrix
\[
 J_{kk}=2k+2,\qquad
 J_{k,k+1}=-\sqrt{(k+1)(k+2)}.
\]
The coefficient vector of $p_{N,x}$ is
$c_k=w_kL_k^{(1)}(x)/\sqrt{k+1}$ for $k\le N$, and zero otherwise.
Write
\[
 S=\|c\|^2,\qquad g=(J-x)c,\qquad \delta=\langle c,g\rangle.
\]

First, squaring \eqref{stage3-localization-PR} for $\alpha=1$ gives
\begin{equation}
 S=\kappa(x)\sqrt N+O(\log(N+2)),\qquad
 \kappa(x)=\frac{128e^x}{315\pi x^{3/2}}.
 \label{stage3-localization-norm}
\end{equation}
Here the nonoscillating sum is
\[
 \frac{e^x}{2\pi x^{3/2}}
 \sum_{k=1}^N k^{-1/2}\left(1-\frac{k}{N+1}\right)^4,
\]
whose leading coefficient is
\[
 \frac{e^x}{2\pi x^{3/2}}
 \int_0^1u^{-1/2}(1-u)^4du
 =\frac{128e^x}{315\pi x^{3/2}}.
\]
The asymptotic remainders sum to $O(\log(N+2))$. For the oscillating
sum, comparison with its integral has the same error bound:
the phase derivative gives an integrable-on-finite-interval bound
$O(1/t)$ for the derivative of the summand. Substitution $t=y^2$
bounds the integral by $O(1)$ using integration by parts. The frequency
$4\sqrt x$ stays bounded away from zero, so all constants are uniform.

Second, the exact Jacobi localization identity is
\[
 \delta=\sum_{k=0}^{N}
 L_k^{(1)}(x)L_{k+1}^{(1)}(x)(w_k-w_{k+1})^2,
 \qquad w_{N+1}=0.
\]
The point bounds from \eqref{stage3-localization-PR} and
$|w_k-w_{k+1}|\le2/(N+1)$ imply
\begin{equation}
 |\delta|=O(N^{-1/2}).
 \label{stage3-localization-mean}
\end{equation}
No sign of the individual products is used.

We also need a weighted residual bound, not just weak concentration.
Put $H=N+1$, $L_{-1}^{(1)}=0$, and
\[
 U_k=L_k(x)+L_{k+1}(x),\qquad
 V_k=L_{k-1}^{(1)}(x)+L_{k+1}^{(1)}(x).
\]
Direct substitution of the quadratic taper gives, for $0\le k\le N$,
\[
 g_k=\sqrt{k+1}
 \left[\frac{2(H-k)}{H^2}U_k-\frac1{H^2}V_k\right],
 \qquad
 g_{N+1}=-\frac{\sqrt{N+2}}{H^2}L_N^{(1)}(x).
\]
All later coefficients vanish. The identities
\[
 L_{k+1}^{(1)}-L_k^{(1)}=L_{k+1},\qquad
 L_{k+1}-L_k=-\frac{x}{k+1}L_k^{(1)}
\]
and \eqref{stage3-localization-PR} give
\[
 U_k=O((k+1)^{-1/4}),\quad
 U_{k+1}-U_k=O((k+1)^{-3/4}),
\]
\[
 V_k=O((k+1)^{1/4}),\quad
 V_{k+1}-V_k=O((k+1)^{-1/4}).
\]
It follows that $g_k=O((k+1)^{1/4}/H)$ in the bulk, while
$g_N,g_{N+1}=O(H^{-5/4})$ at the endpoint. Thus
$\|g\|^2=O(N^{-1/2})$.

For the additional factor of $J$, use the exact nonnegative energy
identity
\[
 \langle g,Jg\rangle
 =g_0^2+\sum_{k\ge0}
   \bigl(\sqrt{k+1}\,g_k-\sqrt{k+2}\,g_{k+1}\bigr)^2.
\]
The displayed difference bounds show that the bulk differences are
$O((k+1)^{1/4}/H)$; the two endpoint contributions are
$O(H^{-3/4})$ before squaring. Therefore
\begin{equation}
 \langle g,(J+1)g\rangle=O(N^{-1/2}).
 \label{stage3-localization-weighted-residual}
\end{equation}
This controls the required unbounded observable on the full half-line.

The exact moment identities now give
\[
 B_M=(x^2-x)S+(2x-1)\delta+\|g\|^2,
\]
\[
 T_M=(x^2+x)S+(2x+1)\delta+\|g\|^2.
\]
Since $\sqrt N=M+O(M^{-1})$, $x=1-M^{-1}$, and $\kappa(1)=\kappa$,
these prove the stated asymptotics of $B_M$ and $T_M$.
Finally,
\[
 U_M=\langle (J-1)c,(J+1)(J-1)c\rangle
 \le2\langle g,(J+1)g\rangle
      +2(x-1)^2\langle c,(J+1)c\rangle
 =O(M^{-1}).
\]
The rational sandwich proves $I_M=\Theta(M^{-1})$:
its lower bound is asymptotic to $\kappa/M$, and its upper bound is
$O(M^{-1})$. Since $|B_M|\to\kappa$, both remaining assertions follow.
\end{proof}

\begin{theorem}[An adverse prime factor and its removal by one cutoff step]
\label{stage3-localization-primes}
For the family above:
\begin{enumerate}
\item If $\sigma=0$ and $M=p^e$ is any prime power, then
\[
 v_p(A_M/B_M)=-2e.
\]
Thus the actual reduced denominator contains $M^2$. Along every
prime-power sequence $M\to\infty$, the primitive error is bounded below
by a positive constant times $M$.
\item If $\sigma=-1$ and $M=p^e$ with $p$ odd, then
\[
 v_p(q_M)\le2e-2.
\]
For $p\ge5$ this improves to
\[
 v_p(q_M)\le\max(0,2e-3).
\]
In particular, at $M=p$ every odd prime $p$, including $p=3$, is absent
from the actual reduced denominator.
\end{enumerate}
These statements concern the indicated prime factor only. The second
case does not bound the cost of all remaining denominator factors.
\end{theorem}

\begin{proof}
Put $u_k=L_k-L_{k+1}$ and $z_k=w_kL_k^{(1)}(x)$, so
$Q_M=\sum_{k=0}^N z_ku_k$. For $M=p^e$,
$v_p(x)=-e$. The finite Laguerre formula gives
\[
 L_k^{(1)}(x)=\frac{(-x)^k}{k!}\bigl(1+p^e\mathbb Z_p\bigr).
\]
Every nonleading term divided by the leading term is an integer times
$x^{-j}$ for some $j\ge1$. In particular the displayed leading unit is
nonzero. Consequently, for $1\le j\le N$,
\begin{equation}
 v_p\!\left(\frac{z_{N-j}}{z_N}\right)
 =je+v_p\!\left(\frac{N!}{(N-j)!}\right)+2v_p(j+1).
 \label{stage3-localization-z-valuation}
\end{equation}

The ordinary Laguerre Gram matrices give exactly
\begin{align}
 B(u_N)&=6N+4,\label{stage3-localization-top-B}\\
 A(u_N)&=(6N+4)H_N+7N+8,\label{stage3-localization-top-A}\\
 2\mathcal A(u_N,u_{N-1})
 &=-2(4N+1)H_N-\frac{32N+7}{3}.
 \label{stage3-localization-adjacent-A}
\end{align}
Here $\mathcal A$ denotes the bilinear form associated with $A$.

First take $N=p^{2e}$. Every ratio in
\eqref{stage3-localization-z-valuation} has valuation at least $3e$.
Since $H_{p^{2e}}$ has the unique summand $1/p^{2e}$ of lowest
$p$-adic valuation,
\[
 v_p(H_{p^{2e}})=-2e.
\]
For odd $p$, \eqref{stage3-localization-top-B} is a unit and
\eqref{stage3-localization-top-A} has valuation $-2e$.
For $p=2$, their valuations are respectively $2$ and $2-2e$.

Write $Q_M/z_N=u_N+\sum_{j\ge1}(z_{N-j}/z_N)u_{N-j}$.
In integer Laguerre coordinates of degree $N+1$, the denominators of
quadratic coefficients divide
$\operatorname{lcm}(1,\ldots,N+1)/2$. Its valuation is $2e$ for odd
$p$ and $2e-1$ for $p=2$. Thus the perturbation in $A$ has valuation at
least $e$ for odd $p$ and $e+1$ for $p=2$. The perturbation in $B$ has
valuation at least $3e$. Both are strictly higher than the corresponding
leading valuations. Division gives $v_p(A_M/B_M)=-2e$ in every case.
The real lower bound then follows from
\eqref{stage3-localization-threshold}.

Now take $N=p^{2e}-1$ with $p$ odd. Pairing $j$ and $p^a-j$ in the
sum over $p$-units shows
\[
 H_{p^a-1}=\frac1pH_{p^{a-1}-1}+V_a,\qquad v_p(V_a)\ge a.
\]
Also $v_p(H_{p-1})\ge1$. Hence
\[
 v_p(H_N)\ge2-2e.
\]
For $p\ge5$, pairing modulo $p^2$ and using
$\sum_{j=1}^{p-1}j^{-2}\equiv0\pmod p$ gives
$v_p(H_{p-1})\ge2$, improving this to $v_p(H_N)\ge3-2e$.

The leading value $B(u_N)=6p^{2e}-2$ is a $p$-adic unit.
Every ratio in \eqref{stage3-localization-z-valuation} has valuation
at least $je$, so all changes to $B$ are divisible by $p$.
Thus the normalized $B$ remains a unit.

For $A$, all linear perturbation terms with $j\ge2$ and all quadratic
perturbation terms are $p$-integral, because their numerator valuations
are at least $2e$ and the global quadratic coefficient denominators
have valuation at most $2e$. Only $j=1$ needs a local check.
For $p\ge5$, \eqref{stage3-localization-adjacent-A} has valuation at least
$\min(0,3-2e)$, and multiplication by $z_{N-1}/z_N$, of valuation $e$,
cannot worsen that bound.

For $p=3$, $v_3(4N+1)=1$. The adjacent coefficient in
\eqref{stage3-localization-adjacent-A} therefore has valuation at least
$\min(3-2e,-1)$. After multiplication by the ratio, its valuation is at
least $\min(3-e,e-1)\ge2-2e$. This includes $p=3,e=1$.
The leading value \eqref{stage3-localization-top-A} satisfies the same
required lower bounds. Since the normalized $B$ is a unit, the stated
bounds for the reduced denominator follow.
\end{proof}

\paragraph{What the one-step repair accomplishes.}
The two cutoffs have the same rigorously proved real error order
$\Theta(1/M)$. At $M=p$, changing $N=p^2$ to $N=p^2-1$ removes the
entire selected prime-power denominator $p^2$ rather than merely
changing a clearing factor. The exact checks give the following
comparison:
\[
\begin{array}{c|cc|cc}
 M=p&
 \text{digits of }q\ (N=p^2)&
 \text{digits of }q\ (N=p^2-1)&
 v_p(q)\ (N=p^2)&v_p(q)\ (N=p^2-1)\\ \hline
 3&20&17&2&0\\
 5&86&79&2&0\\
 7&184&177&2&0
\end{array}
\]
The remaining denominators in these examples are large.
No estimate $q_M=o(M)$ is inferred from this local improvement.


% Additional synthesis result, independently derived 31 October 2026.
% Uses the localized family and its proved arithmetic estimates in the
% unified manuscript.  No source manuscript is modified by this fragment.

\subsection{The repaired cutoff still fails on two fixed-prime towers}
\label{unified-repaired-towers-section}

The cutoff change $M^2\mapsto M^2-1$ removes the selected prime
entirely when $M$ itself is an odd prime.  It does not have that
effect along prime-power towers with a fixed prime.  The following
additional consequence gives exact valuations on two such towers.
It is a further obstruction within the present construction, not a
claim that every modified localization schedule fails.

\begin{proposition}[Exact denominator costs after repair on the $3$- and $5$-power towers]
\label{unified-repaired-towers}
Use the localized family with $\sigma=-1$, so
$N=M^2-1$, and let $q_M$ be the actual reduced denominator of
$A_M/B_M$.  Then, for every integer $e\ge1$,
\[
 \begin{aligned}
  M=3^e&:\qquad v_3(q_M)=2e-2,\\
  M=5^e&:\qquad v_5(q_M)=\max(0,2e-3).
 \end{aligned}
\]
In particular,
\[
 \frac{q_M}{M}\ge\frac{M}{9}\quad(M=3^e),
 \qquad
 \frac{q_M}{M}\ge\frac{M}{125}\quad(M=5^e, e\ge2).
\]
Together with Theorem~\ref{stage3-localization-asymptotic}, these
bounds imply that the primitive errors tend to $+\infty$ on
both fixed-prime towers, even with the repaired cutoff.
\end{proposition}

\begin{proof}
First we determine the relevant harmonic valuations exactly.
For an odd prime $p$ and an integer $a\ge2$, write
\[
 H_{p^a-1}=\frac1pH_{p^{a-1}-1}+V_{p,a},
 \qquad
 V_{p,a}=\sum_{\substack{1\le j<p^a\\p\nmid j}}\frac1j.
\]
Pairing $j$ with $p^a-j$ gives
\[
 \frac1j+\frac1{p^a-j}
 =\frac{p^a}{j(p^a-j)},
\]
so $v_p(V_{p,a})\ge a$.  The exact initial values
$H_2=3/2$ and $H_4=25/12$ now imply, by induction,
\begin{equation}\label{unified-repaired-harmonic-valuations}
 v_3(H_{3^a-1})=2-a,
 \qquad
 v_5(H_{5^a-1})=3-a
 \qquad(a\ge1).
\end{equation}
Indeed, at each step $a\ge2$, the inherited valuations are
$2-a$ and $3-a$, respectively, both strictly smaller than the
unit-block lower bound $a$.  Thus no cancellation can change
the leading valuation.

We next check that these harmonic valuations survive the complete
quadratic form.  Use the notation
\[
 u_k=L_k-L_{k+1},\qquad
 z_k=\left(\frac{N+1-k}{N+1}\right)^2
             L_k^{(1)}\!\left(1-\frac1M\right).
\]
For $M=p^e$, the leading-term valuation of the Laguerre polynomial
shows $z_N\ne0$.  Normalize the polynomial by this coefficient:
\[
 \frac{Q_M}{z_N}=u_N+\sum_{j=1}^N t_j u_{N-j},
 \qquad t_j=\frac{z_{N-j}}{z_N}.
\]
The exact coefficient estimate proved in
Theorem~\ref{stage3-localization-primes} is
\begin{equation}\label{unified-repaired-ratio-estimate}
 v_p(t_j)=je+
  v_p\!\left(\frac{N!}{(N-j)!}\right)+2v_p(j+1)
 \ge je.
\end{equation}
For $j=1$ and either $p=3$ or $p=5$, equality is $v_p(t_1)=e$,
because $N=p^{2e}-1$ and $2$ are $p$-adic units.

The exact leading coefficients are
\[
 B(u_N)=6N+4=6p^{2e}-2,
 \qquad
 A(u_N)=(6p^{2e}-2)H_N+7p^{2e}+1.
\]
The first is a $p$-adic unit for both primes.
For $e\ge2$, \eqref{unified-repaired-harmonic-valuations} gives
the strictly negative valuations
\begin{equation}\label{unified-repaired-top-valuations}
 v_3(A(u_N))=2-2e,
 \qquad
 v_5(A(u_N))=3-2e.
\end{equation}
There is no cancellation within this leading expression:
the harmonic term has negative valuation, while
$7p^{2e}+1$ is a unit.

All coefficients of the quadratic form $B$ in integer Laguerre
coordinates are integers.  Therefore every change from $B(u_N)$
after substituting \eqref{unified-repaired-ratio-estimate} is
divisible by $p$, and
\[
 v_p\bigl(B(Q_M/z_N)\bigr)=0.
\]
For $A$, the half-lcm clearing factor in degree $N+1=p^{2e}$
shows that every quadratic coefficient has valuation at least
$-2e$ for these odd primes.  Consequently, every linear
perturbation $t_j\,2\mathcal A(u_N,u_{N-j})$ with $j\ge2$
has valuation at least $je-2e\ge0$, and every quadratic
perturbation in the $t_j$ has valuation at least
$(j+k)e-2e\ge0$.

It remains to check the single adjacent linear perturbation.
Its coefficient is exactly
\[
 C_N:=2\mathcal A(u_N,u_{N-1})
 =-2(4N+1)H_N-\frac{32N+7}{3}.
\]
For $p=3$, the factor $4N+1=4\cdot3^{2e}-3$ has valuation
one, while $32N+7=32\cdot3^{2e}-25$ is a unit.  Hence
\[
 v_3(t_1C_N)\ge e+\min(3-2e,-1)
 =\min(3-e,e-1)>2-2e\qquad(e\ge2).
\]
For $p=5$, $4N+1$ and $3$ are units, and
\eqref{unified-repaired-harmonic-valuations} gives
\[
 v_5(t_1C_N)\ge e+(3-2e)=3-e>3-2e
 \qquad(e\ge2).
\]
Thus every perturbation has valuation strictly greater than the
corresponding negative leading valuation in
\eqref{unified-repaired-top-valuations}.  The full normalized
$A$ retains that valuation.  Since the normalized $B$ is a unit
and quadratic homogeneity cancels $z_N^2$ in $A_M/B_M$, this proves
\[
 v_3(A_M/B_M)=2-2e,
 \qquad
 v_5(A_M/B_M)=3-2e
 \qquad(e\ge2).
\]
Reduction of a rational number makes its denominator valuation
the negative part of its valuation, giving the asserted exponents
for $e\ge2$.  At $e=1$, the odd-prime removal conclusion of
Theorem~\ref{stage3-localization-primes} gives exponent zero in
both cases, which is exactly the formula stated here.

The resulting divisibilities are $M^2/9\mid q_M$ when
$M=3^e$, and $M^2/125\mid q_M$ when $M=5^e$ and $e\ge2$.
They imply the displayed bounds for $q_M/M$.
Finally the established estimate
$\mathcal E_M=\Theta(q_M/M)$ gives divergence of the primitive
error on both towers.
\end{proof}

\paragraph{Small exact checks and their role.}
The leading-unit checks at $e=2$ reduce to the finite rational sums
$H_{80}$ and $H_{624}$.  Exact arithmetic gives
\[
 \begin{array}{c|c|c|c}
 (p,M,N)&v_p(H_N)&
 p^{-v_p(A(u_N))}A(u_N)\pmod p&B(u_N)\pmod p\\ \hline
 (3,9,80)&-2&2&1\\
 (5,25,624)&-1&4&3
 \end{array}
\]
The perturbation bounds above show that the same scaled $A$
residues hold for $A(Q_M/z_N)$, and that the $B$ residues also
remain unchanged.  Thus these checks certify denominator exponents
$2$ and $1$, respectively.  The all-index assertions follow from
the valuation proof, not extrapolation from these checks.

\subsection{A fresh-prime test after the cutoff repair}
\label{sec:stage3-fresh-primes}

The removal of a prime dividing the localization parameter does not
control primes of the size of the cutoff.  The following exact test
detects many such primes in the \emph{actual} reduced denominator.
Here $A(Q)$ and $B(Q)$ have their positive-square meaning,
$A(Q)-\gamma B(Q)=\int_0^\infty Q(t)^2(t-1)\log t\,e^{-t}\,dt$.

\begin{proposition}[Fresh-prime denominator test]
\label{prop:stage3-fresh-prime}
Let $M\ge2$, $H=M^2$, $N=H-1$, $x=(M-1)/M$, and
\[
 w_k=\left(\frac{H-k}{H}\right)^2,\qquad
 Q(t)=\sum_{k=0}^{N}w_kL_k^{(1)}(x)
                    \bigl(L_k(t)-L_{k+1}(t)\bigr).
\]
Let $\ell$ be an odd prime with $H/2<\ell\le H$, and set
$s=N-\ell$.  Define the truncated polynomial
\[
 W_{M,\ell}(t)=\sum_{r=0}^{s}w_rL_r^{(1)}(x)
                    \bigl(L_r(t)-L_{r+1}(t)\bigr),\qquad
 R_{M,\ell}=x^2B(W_{M,\ell})\pmod\ell.
\]
Then $s\ge0$, $s\le\ell-2$, the residue on the right is defined,
and the two numbers on the left of
\begin{equation}
 \ell^3A(Q)\equiv\ell^2B(Q)\equiv R_{M,\ell}\pmod\ell
 \label{eq:stage3-fresh-prime-residue}
\end{equation}
are $\ell$-integral.  If $R_{M,\ell}\ne0$, then
\[
 v_\ell(A(Q))=-3,\qquad v_\ell(B(Q))=-2,\qquad
 v_\ell\!\left(\operatorname{den}\frac{A(Q)}{B(Q)}\right)=1.
\]
In particular, if $B(Q)\ne0$, the completely computable integer
\[
 G_M=\prod_{\substack{\ell\text{ odd prime},\ H/2<\ell\le H\\
                         R_{M,\ell}\ne0}}\ell
\]
divides the actual reduced denominator $q_M$.
\end{proposition}

\begin{proof}
Since $H=M^2$ is composite, $\ell\le H-1$, and $H<2\ell$
gives $0\le s\le\ell-2$.  Such a prime divides neither $M$ nor
$M-1$, so $x,H$ are $\ell$-adic units.  The finite Laguerre formula,
Wilson's theorem, Lucas's congruence, and $x^\ell\equiv x$ give
\begin{equation}
 \ell L_{\ell+r}^{(1)}(x)\equiv xL_r^{(1)}(x)\pmod\ell
 \qquad(0\le r\le s).
 \label{eq:stage3-frobenius-block}
\end{equation}
For completeness, in
$L_n^{(1)}(x)=\sum_{j=0}^n(-1)^j\binom{n+1}{j+1}x^j/j!$,
the terms $j<\ell$ disappear after multiplication by $\ell$.
Writing $j=\ell+a$ in the remaining terms gives
\[
 \frac{\ell}{(\ell+a)!}\equiv-\frac1{a!},\qquad
 \binom{\ell+r+1}{\ell+a+1}\equiv\binom{r+1}{a+1},
\]
which proves \eqref{eq:stage3-frobenius-block}.

Write $z_k=w_kL_k^{(1)}(x)$ and let $c$ be the ordinary Laguerre
coefficient vector of $Q$:
\[
 c=(z_0,z_1-z_0,\ldots,z_N-z_{N-1},-z_N).
\]
The vector $\widetilde c=\ell c$ is $\ell$-integral, vanishes
modulo $\ell$ in indices $0,\ldots,\ell-1$, and in indices
$\ell,\ldots,H$ is congruent to $x$ times the \emph{complete}
ordinary Laguerre coefficient vector of $W_{M,\ell}$, including
its final coefficient $-z_s$.  This uses $w_{\ell+r}\equiv w_r$.
The exact Gram form
\[
 B(c)=\sum_{i=0}^{H}2i c_i^2
       -2\sum_{i=0}^{H-1}(i+1)c_ic_{i+1}
\]
therefore gives $B(\widetilde c)\equiv x^2B(W_{M,\ell})$.

All entries of $\ell A$ are $\ell$-integral because $H<2\ell$.
For indices $i,j\ge\ell$, its diagonal residues are $2i$,
its adjacent cross-term residues are $-2(i+1)$, and its
remaining cross-term residues vanish.  Indeed
$\ell H_i\equiv1$ for $\ell\le i\le H$ and
\[
 2A_{ij}=\frac{2(i+1)}{j-i-1}
        -\frac{4i}{j-i}
        +\frac{2i}{j-i+1}\qquad(j-i\ge2).
\]
The only possible exceptional gap within the high block is
$j-i=\ell-1$.  It requires $H=2\ell-1$, $i=\ell$, and
$j=2\ell-1$; the potential residue contains $i\equiv0$
and is zero.  All terms involving a low coefficient vanish
because that coefficient is divisible by $\ell$.
Consequently $\ell A(\widetilde c)\equiv B(\widetilde c)$.
Quadratic homogeneity proves \eqref{eq:stage3-fresh-prime-residue},
and a nonzero residue gives the three exact valuations.
\end{proof}

\paragraph{Scope and a zero-residue example.}
The nonzero-residue condition is sufficient, and is not necessary.
For $M=5$, $N=24$, $\ell=17$, $s=7$, one has
$x\equiv11\pmod{17}$ and $B(W_{5,17})\equiv0\pmod{17}$.
Exact rational arithmetic nevertheless gives
\[
 v_{17}(A(Q))=-2,\qquad v_{17}(B(Q))=0,\qquad
 v_{17}(q_5)=2.
\]
Thus a zero residue does not certify cancellation of the prime.
The accompanying verifier checks the full rational forms and actual
reduced denominators for $2\le M\le11$.  Its separate finite-field
scan over $2\le M\le30$ has $786$ eligible pairs, of which $781$
have nonzero residue.  The five zero pairs $(M,\ell)$ are
$(5,17)$, $(17,197)$, $(19,241)$, $(23,487)$, and $(24,449)$.
These finite data do not establish a lower bound for $G_M$ as
$M\to\infty$.  Combined with the proved localization estimate,
the proposition gives $\mathcal E_M\gg G_M/M$ for sufficiently
large $M$; it supplies no unconditional global growth rate for
$G_M$, and no irrationality proof.



\section{Exact logarithm cancellation and the nonvanishing barrier}
\label{stage3-counting-section}

This section changes the construction. It obtains small log-free integer
forms from a dense block of proper kernels. Positivity is not preserved
by the counting argument. We prove that the same argument also produces
formal-zero combinations, and then give a fixed-degree positive
construction that would work if a specific independent-pair height bound
were established. That height bound is not proved here.

For $n\ge1$, define
\[
 R_n(u)=\frac{\prod_{k=0}^{2n-1}(4u+k^2)}
              {\prod_{j=n}^{3n}(u+j^2)},\qquad
 I_n=\int_0^\infty R_n(u)\,d\mu(u)>0.
\]
The residues, constant coefficient, and prime-log coefficients are
\begin{align*}
 A_{n,j}&=(-1)^{j-n}
 \frac{4j^2(2j+2n-1)!(n+j-1)!}
 {(2j-2n)!(j-n)!(3n-j)!(3n+j)!},\\
 p_n&=\sum_{j=n}^{3n}A_{n,j}\left(H_j-\frac1{2j}\right),\qquad
 E_{n,\ell}=\sum_{j=n}^{3n}A_{n,j}v_\ell(j).
\end{align*}
There is no polynomial part in the partial fraction decomposition, and
$\sum_j A_{n,j}=16^n$. Thus Binet's formula gives exactly
\begin{equation}\label{stage3-proper-linear-form}
 I_n=16^n\gamma-p_n+\sum_{\ell\le3n}E_{n,\ell}\log\ell.
\end{equation}
Only primes occur in the last sum.

\begin{lemma}[Uniform height and raw integral bounds]
\label{stage3-proper-bounds}
For every $n\ge1$ and $n\le j\le3n$,
\[
 2A_{n,j}\in\mathbb Z,\qquad
 2|A_{n,j}|\le6\,2^{12n}.
\]
Moreover
\begin{equation}\label{stage3-proper-U}
 0<I_n<U_n:=
 \frac{(4n)!((n-1)!)^2}{8\,16^n((3n)!)^2},
 \qquad U_{n+1}<U_n,
\end{equation}
and, with $\kappa_0=\log(729/16)$,
\[
 \log U_n=-\kappa_0 n+O(\log(n+1)).
\]
\end{lemma}

\begin{proof}
Put $a=j+n$ and $b=j-n$. With
\[
 B=\binom{2a}{b},\quad C=\binom a{2b},\quad
 B'=\binom{2a-1}{b-1},\quad C'=\binom{a-1}{2b-1},
\]
where $B'=C'=0$ at $b=0$, direct factorial simplification gives
\[
 2|A_{n,j}|=BC+4B'C+B'C'.
\]
This proves integrality. Every product in this expression is at most
$2^{3a}\le2^{12n}$, proving the stated height bound.

For the integral estimate, put
$d\nu(u)=du/(2\sinh(2\pi\sqrt u))$. Then
$1<d\mu/d\nu=1+e^{-2\pi\sqrt u}<2$.
Let $T_j(u)=P_{2j}^{(1)}(2\sqrt u;\pi/2)$ be the even
Meixner--Pollaczek polynomials in their usual normalization.
Their generating function and orthogonality
\cite{dlmf-mp,dlmf-mp-generating} imply
\[
 S_r(u):=\sum_{j=0}^r\binom rj T_j(u)
       =\frac{4^r}{(2r)!}\prod_{k=0}^{r-1}(4u+k^2),
 \qquad \int_0^\infty S_r\,d\nu=\frac1{16}.
\]
For clarity, the generating function of the binomial transform is
$\cosh(4\sqrt u\arcsin\sqrt z)$, whose coefficient gives the displayed
product. In the integral, all $T_j$ with $j>0$ vanish by orthogonality;
the constant norm is $\int d\nu=1/16$. The product is positive for
$u>0$, so $\int S_r d\mu<1/8$.
Since
\[
 R_n(u)=\frac{(4n)!}{16^n}
          \frac{S_{2n}(u)}{\prod_{j=n}^{3n}(u+j^2)},
\]
replacing each denominator factor by $j^2$ gives
\eqref{stage3-proper-U}. A direct quotient estimate gives
\[
 \frac{U_{n+1}}{U_n}
 =\frac{(4n+1)(4n+2)(4n+3)(4n+4)n^2}
 {16(3n+1)^2(3n+2)^2(3n+3)^2}
 \le\frac{16}{729}\left(1+\frac1n\right)^4
 \le\frac{256}{729}<1.
\]
Finally Stirling's formula gives the logarithmic estimate.
\end{proof}

\begin{theorem}[Thin-block cancellation with an exponential bound, possibly zero]
\label{stage3-counting-main}
For integers $N\ge2$, define using integer operations
\[
 h_N=\max\left(1,\left\lfloor\sqrt{\lfloor\log_2N\rfloor}\right\rfloor\right),
 \quad M=N+\left\lceil\frac N{h_N}\right\rceil,
 \quad d=M-N+1,
\]
\[
 D_N=4\operatorname{lcm}(1,\ldots,3M),\qquad
 \delta_* =\log(729/16)-3>0.
\]
For every sufficiently large $N$, there is a nonzero integer vector
$c=(c_N,\ldots,c_M)$ satisfying
\[
 \sum_{n=N}^M c_n E_{n,\ell}=0\quad\hbox{for every prime }\ell,
 \qquad
 \log\max(1,\|c\|_\infty)
       =O\!\left(\frac N{\sqrt{\log N}}\right).
\]
Its exactly log-free form is
\begin{equation}\label{stage3-counting-small-form}
 F_N=D_N\sum_{n=N}^M c_n I_n=Q_N\gamma-P_N,
 \qquad Q_N,P_N\in\mathbb Z,
\end{equation}
and satisfies
\begin{equation}\label{stage3-counting-upper-bound}
 |F_N|\le\exp(-\delta_*N+o(N)).
\end{equation}
This is an upper bound; $F_N$ is allowed to be exactly zero.

Indeed, for every sufficiently large $N$, there also exists a nonzero
integer vector of the same height order for which all prime-log
coefficients, $Q_N$, and $P_N$ are zero simultaneously. Its rational
kernel $\sum c_nR_n$ is nonzero and changes sign on $(0,\infty)$.
\end{theorem}

\begin{proof}
Use the integer log matrix
\[
 V_{\ell,n}=2E_{n,\ell},\qquad
 \ell\le3M\hbox{ prime},\quad N\le n\le M.
\]
It has $r=\pi(3M)$ rows and $d$ columns, and its entries have absolute
value at most the integer
\[
 H_N=6(2M+1)\lceil\log_2(3M)\rceil 2^{12M}.
\]
The prime number theorem \cite{dlmf-primes} and the chosen block give
\[
 M/N\to1,\quad d\asymp\frac N{\sqrt{\log N}},\quad
 r\asymp\frac N{\log N},\quad
 \log H_N=O(N).
\]
Thus $d>r$ for all sufficiently large $N$.

Here is the full counting argument. For any $r$-by-$d$ integer matrix
with entries bounded by $H$, let $K$ be a positive integer such that
\[
 K^{d-r}>(3dH)^r.
\]
The $(K+1)^d$ vectors in $\{0,\ldots,K\}^d$ have images with at most
$(2dKH+1)^r\le(3dKH)^r<K^d$ possible values. Two images coincide.
Subtracting the input vectors gives a nonzero integer null vector of
supremum norm at most $K$. Taking the smallest admissible integer $K$
yields
\[
 \log K\le
 \frac r{d-r}\log(3dH)+O(1)
 =O\!\left(\frac N{\sqrt{\log N}}\right).
\]
This proves exact log cancellation and the height estimate. It also
gives a finite, though not efficient, construction by enumeration.

Since $2A_{n,j}$ is integral, every denominator in $p_n$ divides
$4\operatorname{lcm}(1,\ldots,3M)=D_N$. Thus the coefficients
\[
 Q_N=D_N\sum c_n16^n,\qquad P_N=D_N\sum c_np_n
\]
in \eqref{stage3-counting-small-form} are integers. Positivity and
monotonicity of the raw bounds give
\[
 |F_N|\le D_NdK U_N.
\]
The same prime number theorem gives
$\log D_N=3M+o(M)=3N+o(N)$, while
$\log(dK)=o(N)$. Lemma~\ref{stage3-proper-bounds} therefore proves
\eqref{stage3-counting-upper-bound}. The margin is positive since
$e<3$ gives $e^3<27<729/16$.

To see exactly why nonvanishing is missing, append to $V$ the two
integer rows $(16^n)_{n=N}^M$ and $(D_Np_n)_{n=N}^M$.
Their entries still have logarithmic height $O(N)$: the residue bound,
$H_j\le j$, and $\log D_N=O(N)$ suffice. Replacing $r$ by $r+2$
does not change the asymptotic height estimate. The same counting
argument therefore gives a nonzero vector annihilating all three
types of coefficients. Its integral is exactly zero, as a formal
identity, independently of the arithmetic nature of $\gamma$.

Let $n_*$ be the largest index with nonzero coefficient. The pole at
$u=-(3n_*)^2$ occurs in $R_{n_*}$ with nonzero residue and in none of
the smaller-index kernels. Hence $\sum c_nR_n$ is not the zero
rational function. It is continuous on $(0,\infty)$ and has zero
integral against a strictly positive measure with full support.
It must take both signs there. This proves the last assertion.
\end{proof}

\subsection{A fixed-degree positive construction from two independent pairs}
Let $\mathcal K_N=\ker_{\mathbb Z}V$ and define the integer coefficient
map
\[
 \Phi_N(c)=\left(D_N\sum c_n16^n,\ D_N\sum c_np_n\right).
\]
For $c\in\mathcal K_N$, the corresponding form is
$F_c(\gamma)=Q_c\gamma-P_c$. Define $\sigma_2(N)$ as the least value of
$\max(\|c^{(1)}\|_\infty,\|c^{(2)}\|_\infty)$ among pairs in
$\mathcal K_N$ with
\[
 \det\begin{pmatrix}Q_1&P_1\\Q_2&P_2\end{pmatrix}\ne0.
\]
Set $\sigma_2(N)=\infty$ if no such pair exists. This definition is a
property of explicitly given integer and rational matrices. It does
not compare any computed value to $\gamma$.

\begin{proposition}[A positive quadratic certificate and its exact missing bound]
\label{stage3-counting-independent-pair}
For every pair with nonzero determinant, the integer polynomial
\[
 W_N(X)=(Q_1X-P_1)^2+(Q_2X-P_2)^2
\]
is strictly positive for every real $X$. If, along an infinite sequence
of indices tending to infinity,
\begin{equation}\label{stage3-counting-sigma-target}
 \limsup\frac{\log\sigma_2(N)}N<\delta_*,
\end{equation}
then, choosing a pair attaining $\sigma_2(N)$ at each of those indices,
the resulting polynomials satisfy $W_N(\gamma)\to0$, and $\gamma$ is irrational.
No estimate \eqref{stage3-counting-sigma-target} is proved here.
\end{proposition}

\begin{proof}
Two independent coefficient pairs cannot vanish at the same real
argument. Equivalently the leading coefficient is positive and the
discriminant of $W_N$ is
\[
 -4(Q_1P_2-Q_2P_1)^2<0.
\]
For a pair attaining $\sigma_2(N)$, the raw estimates give
\[
 |F_{c^{(i)}}(\gamma)|\le D_Nd\sigma_2(N)U_N.
\]
The logarithm of this upper bound is
$-\delta_*N+\log\sigma_2(N)+o(N)$. Under
\eqref{stage3-counting-sigma-target} both forms tend to zero, and so
does their sum of squares. If $\gamma=a/b$ were rational, then
\[
 b^2W_N(\gamma)
 =(aQ_1-bP_1)^2+(aQ_2-bP_2)^2\in\mathbb Z_{>0},
\]
contradicting convergence to zero. The degree is the fixed number $2$,
so this integerization introduces the fixed factor $b^2$ only.
\end{proof}

The preceding construction separates positivity from the sign of an
integral kernel. It replaces that sign problem by an exact determinant
test. It does not supply cheap independent pairs: the counting theorem
controls a short null vector, which can lie in $\ker\Phi_N$.

\begin{proposition}[Why dimension and entry height alone do not bound the pair]
\label{stage3-counting-countermodel}
For integers $H\ge2$, $s\ge1$, and any number of free coordinates
$z_1,\ldots,z_k$ with $k\ge1$, impose
\[
 x_i-Hx_{i+1}=0,\quad y_i-Hy_{i+1}=0\quad(0\le i<s),
 \qquad\Phi(x,y,z)=(x_s,y_s).
\]
The constraint matrix has rank $2s$, entries of absolute value at most
$H$, and arbitrarily many independent null vectors of height $1$ with
$\Phi=0$. Appending the two rows of $\Phi$ raises its rank by $2$.
Nevertheless the least height of two null vectors with independent
$\Phi$-images is exactly $H^s$.
\end{proposition}

\begin{proof}
Every null vector has
$x_i=H^{s-i}x_s$ and $y_i=H^{s-i}y_s$. If its height is less than
$H^s$, both $x_s$ and $y_s$ are zero. The free coordinate vectors have
height $1$, and the two choices $(x_s,y_s)=(1,0),(0,1)$ give independent
images at height exactly $H^s$. The rank assertions follow directly
from the independent chain constraints.
\end{proof}

This is not a failure theorem for the specific proper-kernel matrix.
It proves that its additional arithmetic structure must be used:
dimension, entry size, and formal rank alone cannot upgrade the
counting theorem to the required independent-pair estimate.


% Additional synthesis result, independently derived 31 October 2026.
% This fragment uses the theorem/lemma environments of the unified manuscript.
% Only this owned fragment was created; the source manuscripts are unchanged.

\section{An actual-matrix height barrier for independent signed outputs}
\label{unified-signed-height-section}

The counting theorem produces short vectors in the integer kernel of the
prime-log matrix, but it does not provide two independent coefficient
pairs.  The following additional synthesis result shows that this
distinction is unavoidable for the actual matrices considered here:
subexponential-height log-null vectors can have formal outputs of rank
at most one.  The proof combines the common gamma-coefficient factor
with the integral estimate already established above.  The underlying
integer determinant argument and Smith normal form are standard tools;
no historical-priority claim is made for them.

\begin{lemma}[An integer-output determinant bound]
\label{unified-signed-determinant-lemma}
Let $\Lambda\subseteq\mathbb Z^d$ be a lattice, let
$a,b:\Lambda\to\mathbb Z$ be homomorphisms, and let $\theta$ be real.
Suppose that, for constants $A,E>0$ and every $c\in\Lambda$,
\[
 |a(c)|\le A\|c\|_\infty,
 \qquad
 |\theta a(c)-b(c)|\le E\|c\|_\infty.
\]
Put $\mathcal L=\{(a(c),b(c)):c\in\Lambda\}\subseteq\mathbb Z^2$.
If $\mathcal L$ has rank two, let
$j=[\mathbb Z^2:\mathcal L]$.  Then every pair with independent
outputs satisfies
\begin{equation}\label{unified-signed-height-product}
 \|c_1\|_\infty\|c_2\|_\infty\ge\frac{j}{2AE}.
\end{equation}
In particular, if $2AEK^2<1$, all outputs with
$\|c\|_\infty\le K$ span a space of dimension at most one,
whether or not the full image has rank two.

When the full image has rank two, $j$ is the greatest common divisor
of the determinants of all pairs of image vectors.  Thus it concerns
the actual image lattice, which need not be saturated in $\mathbb Z^2$.
\end{lemma}

\begin{proof}
Set $e_i=\theta a(c_i)-b(c_i)$.  Cancellation of $\theta$ gives
\[
 a(c_1)b(c_2)-a(c_2)b(c_1)
   =a(c_2)e_1-a(c_1)e_2.
\]
Its absolute value is at most
$2AE\|c_1\|_\infty\|c_2\|_\infty$.
A nonzero determinant of two vectors of $\mathcal L$ is a nonzero
integer multiple of the covolume $j$, proving
\eqref{unified-signed-height-product}.

For completeness, a rank-two sublattice of $\mathbb Z^2$ has a basis
$(a_0,b_0),(0,b_1)$ with $a_0,b_1>0$: choose $a_0$ as the positive
generator of its first-coordinate image, choose a vector above it,
and choose $b_1$ as the positive generator of its vertical kernel.
Subtracting a multiple of the first vector proves generation by these
two vectors.  The index and the absolute basis determinant are both
$a_0b_1$.  Every pair determinant is a multiple of this number, and
the basis pair attains it, proving the gcd assertion.

Finally, if $2AEK^2<1$, every determinant of two outputs in the
specified height ball is an integer of absolute value less than one,
and therefore is zero.  If one such output is nonzero, every other
one is collinear with it; if none is nonzero, the assertion is
immediate.
\end{proof}

\begin{lemma}[The exact image index and the fibres of the output map]
\label{unified-signed-image-index}
Let $C\in\mathbb Z^{r\times d}$ have row rank $r$, let
$B\in\mathbb Z^{2\times d}$, and suppose
$\operatorname{rank}\binom CB=r+2$.
For a full-row-rank integer matrix $T$ with $s$ rows, write
$\Delta_s(T)$ for the positive gcd of its $s$-by-$s$ minors;
use $\Delta_0(C)=1$ when $r=0$.  Then
\begin{equation}\label{unified-signed-index-formula}
 \left[\mathbb Z^2:B(\ker_{\mathbb Z}C)\right]
 =\frac{\Delta_{r+2}\!\left(\binom CB\right)}{\Delta_r(C)}.
\end{equation}
For any attained output $y=Bc_0$ with $Cc_0=0$, its entire fibre is
\[
 \{c\in\mathbb Z^d:Cc=0,\ Bc=y\}
 =c_0+\ker_{\mathbb Z}\binom CB.
\]
Both integer kernels in this statement are saturated in
$\mathbb Z^d$.  This does not assert that the image
$B(\ker_{\mathbb Z}C)$ is saturated in $\mathbb Z^2$.
\end{lemma}

\begin{proof}
Put $\mathcal H=\binom CB\mathbb Z^d$,
$\mathcal C=C\mathbb Z^d$, and
$\mathcal L=B(\ker_{\mathbb Z}C)$.
Projection to the first $r$ coordinates gives the exact sequence
of finite abelian groups
\[
 0\longrightarrow\mathbb Z^2/\mathcal L
 \longrightarrow\mathbb Z^{r+2}/\mathcal H
 \longrightarrow\mathbb Z^r/\mathcal C
 \longrightarrow0.
\]
Indeed, a class whose first coordinate lies in $\mathcal C$ can
be represented by $(0,y)$ after subtracting a vector of
$\mathcal H$, and $(0,y)$ lies in $\mathcal H$ exactly when
$y\in\mathcal L$.  Taking group orders gives the ratio of the
two image indices.

For a full-row-rank integer matrix, its image index is the gcd of
its maximal minors.  To see this, apply integer unimodular row and
column operations to obtain Smith normal form.  Such operations
preserve the image index and the gcd of the maximal minors; on
the diagonal form both are the product of the nonzero diagonal
entries.  This proves \eqref{unified-signed-index-formula}.
The fibre identity follows by subtracting $c_0$.
An integer kernel is saturated because $kTc=0$, for a nonzero
integer $k$, implies $Tc=0$ in the torsion-free group
$\mathbb Z^s$.
\end{proof}

\begin{theorem}[Independent signed outputs have an exponential height cost]
\label{unified-signed-actual-barrier}
Use the proper kernels, their rational bounds $U_n$ in
\eqref{stage3-proper-U}, and the thin blocks of
Theorem~\ref{stage3-counting-main}.  Thus
\[
 \mathcal K_N=\ker_{\mathbb Z}V,
 \quad D_N=4\operatorname{lcm}(1,\ldots,3M),
 \quad \delta_*=\log(729/16)-3.
\]
Set the explicit positive rational quantities
\[
 \mathcal S_N=\sum_{n=N}^M16^{n-N}
     =\frac{16^{M-N+1}-1}{15},
 \qquad
 \mathcal T_N=\sum_{n=N}^M U_n,
 \qquad G_N=D_N16^N.
\]
For $c\in\mathcal K_N$, write the integer form as
\[
 F_c(\gamma)=Q_c\gamma-P_c,
 \qquad
 a_N(c)=\frac{Q_c}{G_N}
       =\sum_{n=N}^M c_n16^{n-N}\in\mathbb Z.
\]
Let $\sigma_2(N)$ be the least possible maximum input height of
two vectors whose original coefficient pairs $(Q_c,P_c)$ are
independent, and set it to $\infty$ if no such pair exists.
Then, for every $N$ for which the block is defined,
\begin{equation}\label{unified-signed-actual-finite}
 \sigma_2(N)\ge
 \bigl(2D_N\mathcal S_N\mathcal T_N\bigr)^{-1/2}.
\end{equation}
If the normalized image
\[
 \mathcal L_N
   =\{(a_N(c),P_c):c\in\mathcal K_N\}
\]
has rank two and index $j_N$ in $\mathbb Z^2$, the sharper bound is
\begin{equation}\label{unified-signed-actual-index}
 \sigma_2(N)\ge
 \sqrt{\frac{j_N}{2D_N\mathcal S_N\mathcal T_N}}.
\end{equation}
In particular, with the extended-real convention for $\log\infty$,
\begin{equation}\label{unified-signed-actual-rate}
 \liminf_{N\to\infty}\frac{\log\sigma_2(N)}N
 \ge\frac{\delta_*}{2}
 =\frac{\log(729/16)-3}{2}>0.
\end{equation}

More strongly, if $K_N\ge1$ and
\[
 \limsup_{N\to\infty}\frac{\log K_N}{N}
 <\frac{\delta_*}{2},
\]
then, for all sufficiently large $N$, the formal outputs of
\emph{all} $c\in\mathcal K_N$ with $\|c\|_\infty\le K_N$
span a space of dimension at most one.  In particular this holds
for every subexponential height schedule $K_N=\exp(o(N))$,
including the heights produced by the counting argument.
\end{theorem}

\begin{proof}
The exact formulas preceding
\eqref{stage3-counting-upper-bound} give
\[
 P_c=D_N\sum_{n=N}^M c_np_n\in\mathbb Z,
 \qquad F_c(\gamma)=D_N\sum_{n=N}^M c_nI_n.
\]
Consequently
\[
 |a_N(c)|\le\mathcal S_N\|c\|_\infty,
 \qquad
 |G_N\gamma a_N(c)-P_c|
 \le D_N\mathcal T_N\|c\|_\infty.
\]
The two original pairs are independent exactly when the normalized
pairs are independent, since their first coordinates differ by
the same nonzero factor $G_N$.
Apply Lemma~\ref{unified-signed-determinant-lemma} with
\[
 \theta=G_N\gamma,\quad A=\mathcal S_N,
 \quad E=D_N\mathcal T_N.
\]
Every nonzero normalized determinant is an integer, and in the
rank-two case it is a nonzero multiple of $j_N$.
This proves both finite bounds; if there is no independent pair,
\eqref{unified-signed-actual-finite} holds by the convention
$\sigma_2(N)=\infty$.

The thin-block choice gives $M-N=o(N)$, so
\[
 \log\mathcal S_N=o(N).
\]
Since $U_n$ decreases,
$U_N\le\mathcal T_N\le(M-N+1)U_N$.  The proved asymptotic for
$U_N$ therefore gives
\[
 \log\mathcal T_N=-\log(729/16)N+o(N).
\]
The already used least-common-multiple asymptotic gives
$\log D_N=3N+o(N)$.  Hence
\begin{equation}\label{unified-signed-determinant-scale}
 2D_N\mathcal S_N\mathcal T_N
 =\exp\{-\delta_*N+o(N)\}.
\end{equation}
Taking logarithms of \eqref{unified-signed-actual-finite}
proves \eqref{unified-signed-actual-rate}.
Under the stated hypothesis on $K_N$, the product of
\eqref{unified-signed-determinant-scale} with $K_N^2$ tends to zero.
The final conclusion follows from the last assertion of
Lemma~\ref{unified-signed-determinant-lemma}.
\end{proof}

\paragraph{Exact arithmetic and scope.}
The index $j_N$, whenever the normalized image has rank two, can be
computed without any evaluation of $\gamma$.
Choose a maximal independent set of rows of the integer log matrix
as $C$, and let the two rows of $B$ be
$(16^{n-N})_{n=N}^M$ and $(D_Np_n)_{n=N}^M$.
Then Lemma~\ref{unified-signed-image-index} gives its exact value.
All vectors with a fixed formal output differ by the saturated
integer kernel of the stacked matrix; changing a vector inside
that fibre changes neither its formal linear form nor its
primitive arithmetic error.

The new conclusion is a lower height bound for the actual
proper-kernel matrices, not just for an unrelated model matrix.
It does not prove that every short output is zero, and it does
not exclude independent outputs at larger heights.
In particular, the previously stated sufficient target
$\limsup\log\sigma_2(N)/N<\delta_*$ is not contradicted:
an exponential-rate window between $\delta_*/2$ and $\delta_*$
remains.  Neither existence of suitable pairs in that window nor
irrationality of $\gamma$ is established here.

\section{The remaining estimates and how the results fit together}
\label{unified-final-scope}

The positive and signed constructions address different requirements.
The following table records proved statements and the precise estimates
that are not supplied by this manuscript. A hypothesis in the last column
is not used as if it had been established.

\small
\begin{longtable}{@{}>{\raggedright\arraybackslash}p{0.22\textwidth}
 >{\raggedright\arraybackslash}p{0.35\textwidth}
 >{\raggedright\arraybackslash}p{0.35\textwidth}@{}}
\toprule Construction & Proved properties & Remaining requirement\\\midrule
\endfirsthead
\toprule Construction & Proved properties & Remaining requirement\\\midrule
\endhead
Canonical interval selection &
Integer coefficients, exact log cancellation, positive forms, terminating
selection, divergent denominators, exact limiting dichotomy &
An independent proof that
\(\widehat q_m(r_{m+1}-r_m)\to0\).\\
Primitive optimization &
An exact finite candidate set; a terminating relative optimizer; a
three-fraction exact certificate when its hypotheses hold &
A proved zero limit of the interval minima for \(\xi=\gamma\).\\
Positive target realization &
Every strict rational lower target has a global positive certificate;
complete reduction gives the target's exact primitive pair &
Selection of an infinite target sequence with primitive error tending
to zero.\\
Localized squares with the repaired cutoff &
Positive integral error of order \(M^{-1}\), exact rational bounds,
and removal of the designated prime when \(M\) is prime &
A global denominator bound \(q_M=o(M)\) along a specified infinite
sequence. The \(3\)- and \(5\)-power towers fail.\\
Signed thin-block combinations &
All logarithms cancel; integerization is included in the exponentially
small upper bound; formal-zero combinations also exist &
Nonzero integrated output with the required height, or sufficiently
small independent output pairs.\\
\bottomrule
\end{longtable}
\normalsize

The new results sharpen this map. The positive dyadic degree law is
optimal on its specified fixed pole grid, even with unlimited coefficient
height. It settles an analytic complexity question but leaves the
target denominator uncontrolled. The interval theorem proves that
least-denominator selection, true primitive optimization, and relative
optimization share a limiting rationality dichotomy. Its Fibonacci
counterexample shows that the finite errors of these procedures can
nevertheless differ by arbitrarily large factors. Finally, the actual
signed matrix requires exponential height for two independent outputs.
In the notation of the signed construction,
\[
 \liminf_{N\to\infty}\frac{\log\sigma_2(N)}N
 \ge\frac{\delta_*}{2},
\]
whereas a sufficient estimate for its proposed two-square route would be
\[
 \limsup_{\substack{N\to\infty\\N\in\mathcal N}}
       \frac{\log\sigma_2(N)}N<\delta_*
\]
on an infinite set \(\mathcal N\). The lower bound does not imply that
such a set exists, and it does not contradict the stated sufficient
upper target.

There is therefore no transfer of nonvanishing from one family to the
smallness estimate of another family. A completed irrationality argument
would have to establish the missing property within one specified
infinite construction. The present results provide exact constructions,
optimization tools, arithmetic exclusions, and complexity laws with
which that further problem can be studied.

\appendix
\section{Source integration, corrected normalization, and reproducibility}
\label{unified-reproducibility}

\subsection{What was integrated}

All ten selected result groups are proved in this manuscript. The
canonical family, target interpolation, rationality dichotomy, and
Laguerre arithmetic come from source A; global positive realization,
effective square tests, and support compression come from source B;
nested brackets, finite optimization, interpolation valuations,
localization, cutoff repair, and signed cancellation come from source C.
The additional synthesis results are the sharp degree law, the general
interval dichotomy, the Fibonacci separation example, the exact repaired
prime-power costs, and the actual-matrix independence barrier.

Repeated definitions and one duplicate proof of the half-lcm clearing
factor have been consolidated. The monomial moment is
\[
 \int_0^\infty t^s(t-1)\log t\,e^{-t}\,dt
 =s!\bigl(1+sH_s-s\gamma\bigr).
\]
One redundant displayed formula in source C contained an extra \(s\)
inside these braces. That typographical error is corrected here.
The quadratic matrices, the source B moments, and the localization
arguments use the normalization displayed above. For example
\(A(L_1)=5\) and \(B(L_1)=2\), so \(I(L_1)=5-2\gamma\).
The integration does not assume that the superseded expression is valid.

Classical inputs include Binet's integral, digamma duplication,
Laguerre and Meixner--Pollaczek identities, the prime number theorem,
and Meyer's rational isotropy theorem. Each is identified at its point
of use. The new monic-polynomial estimate and the image-index formula
include their own proofs from factorization, integration by parts,
and integer lattice operations.

\subsection{Exact checks and their logical role}

The proof of the sharp degree estimate is symbolic. As a separate
finite check, the leading coefficients, orthogonality, and norms in
its shifted Jacobi calculation were verified in rational arithmetic for
\(\alpha=0,1\) and degrees \(0,\ldots,12\). The signed determinant
argument and the image-index identity were independently checked
against the actual proper-kernel coefficient map.

The short Python program below uses only the standard library and
exact fractions. It verifies the seven displayed finite error bounds,
membership of their targets in the claimed canonical intervals, and
the three-fraction certificate for the exact optimum at \(m=10\).
It also checks the two strict rational bounds surrounding that optimum.
Its successful output is
\begin{center}
\small\texttt{PASS: rational targets, seven error bounds, and exact m=10 optimum}
\end{center}
The reported finite candidate counts come from the finite
ceiling-Euclidean construction; they are not needed by this verifier
or by any asymptotic theorem.

These checks supplement the all-index proofs. In particular they do
not prove convergence of the gamma-specific primitive errors.
The TeX source contains the entire manuscript and bibliography:
there are no external \texttt{\string\input}, \texttt{\string\include},
figure, or BibTeX-file dependencies. A standard TeX installation
with the listed packages can compile it using \texttt{pdflatex} twice
or \texttt{latexmk -pdf}.

\subsection{A self-contained rational certificate verifier}
The following code contains no decimal approximation of \(\gamma\).
Finite decimal strings in the comparison table are parsed as exact
rational upper bounds.

\begingroup\footnotesize
\begin{verbatim}

from fractions import Fraction as F
from math import gcd

def bernoulli(n):
    a = [F(0)] * (n + 1)
    for m in range(n + 1):
        a[m] = F(1, m + 1)
        for j in range(m, 0, -1):
            a[j-1] = j * (a[j-1] - a[j])
    return a[0]

def mu(j):
    return abs(bernoulli(2*j + 2)) / (2*j + 2)

# Store only the harmonic numbers needed by dyadic targets.
harm, h = {}, F(0)
for n in range(1, 2**14 + 1):
    h += F(1, n)
    if n & (n-1) == 0:
        harm[n] = h

def target(m):
    p = [1]
    for j in range(1, m+1):
        t = [0] * (len(p)+1)
        for i, v in enumerate(p):
            t[i] -= v
            t[i+1] += 4**j * v
        p = t
    a = sum(p)
    s = F(sum(i*v for i, v in enumerate(p)), a)
    c = [F(0)] * (len(p)+1)
    for i, v in enumerate(p):
        c[i] += (1+s) * F(v, a)
        c[i+1] -= s * F(v, a)
    assert sum(c) == 1
    assert sum(i*v for i, v in enumerate(c)) == 0
    return sum(v*(harm[2**i]-F(1, 2**(i+1)))
               for i, v in enumerate(c))

r = [target(m) for m in range(14)]
assert all(r[i] < r[i+1] for i in range(13))
rho = lambda m: F(9*(m+1)**2, 16*4**(m+1))
upper = lambda m: r[m] + (r[m+1]-r[m])/(1-rho(m))

# A rational Binet bracket, with an explicit positive remainder.
ell_lo = 2*sum(F(1, (2*j+1)*3**(2*j+1))
               for j in range(110))
ell_hi = ell_lo + F(2, 221*3**221)*F(9, 8)
s = sum((-1)**j * mu(j)/64**(2*j+2) for j in range(20))
rem = mu(20)/64**42
g_lo = harm[64]-F(1, 128)-6*ell_hi+s
g_hi = harm[64]-F(1, 128)-6*ell_lo+s+rem
assert g_lo < g_hi

rows = [
    (0, 4, 7, '4.05096544e-2'),
    (1, 15, 26, '7.60728744e-3'),
    (2, 157, 272, '2.66085322e-3'),
    (4, 33841, 58628, '1.84706856e-6'),
    (6, 9062634, 15700603, '1.53583719e-9'),
    (8, 8788062966751, 15224921118955, '9.04971062e-14'),
    (10, 2323006970406828281, 4024504377931443899,
     '3.81603071e-19'),
]
for m, p, q, bound in rows:
    assert gcd(p, q) == 1
    assert r[m] <= F(p, q) <= r[m+1]
    assert 0 < q*g_lo-p < q*g_hi-p < F(bound)

p, q = 2323006970406828281, 4024504377931443899
pc, qc = 1055624077073011532, 1828820909171081307
assert q*pc-p*qc == 1
mediant = F(p+pc, q+qc)
assert r[11] < mediant < r[12]
assert upper(12) < F(pc, qc)
assert F(38160307067862345768, 10**38) < q*g_lo-p
assert q*g_hi-p < F(38160307067862345769, 10**38)
print('PASS: rational targets, seven error bounds, and exact m=10 optimum')

\end{verbatim}
\endgroup

\begin{thebibliography}{99}

\bibitem{source-a}
Keio--Shaitetsu research notes, \emph{Constructive Gamma Kernels and
Exact Arithmetic Costs}, 31 October 2026, internal source manuscript.

\bibitem{source-b}
Keio--Shaitetsu research notes, \emph{Eighteen Auxiliary Theorems for
Gamma Approximation}, 31 October 2026, internal source manuscript.

\bibitem{source-c}
Keio--Shaitetsu research notes, \emph{Primitive Costs for Euler's Constant},
31 October 2026, internal source manuscript.

\bibitem{dlmf-binet}
NIST Digital Library of Mathematical Functions, \S5.9,
especially equation 5.9.15 (Binet's digamma integral).
\url{https://dlmf.nist.gov/5.9.E15}.

\bibitem{dlmf-duplication}
NIST Digital Library of Mathematical Functions, \S5.5,
equation 5.5.8 (digamma duplication).
\url{https://dlmf.nist.gov/5.5.E8}.

\bibitem{dlmf-laguerre}
NIST Digital Library of Mathematical Functions, \S18.12,
equation 18.12.13 (Laguerre generating function), and \S18.3
(orthogonality and normalization).
\url{https://dlmf.nist.gov/18.12.E13};
\url{https://dlmf.nist.gov/18.3}.

\bibitem{dlmf-primes}
NIST Digital Library of Mathematical Functions, \S27.2,
equation 27.2.3 (the prime number theorem).
The corresponding $\vartheta(x)\sim x$ follows by partial summation.
\url{https://dlmf.nist.gov/27.2.E3}.

\bibitem{dlmf-laguerre-asymptotic}
NIST Digital Library of Mathematical Functions,
Section 18.15(iv), Laguerre asymptotics, equations 18.15.14--15.
\url{https://dlmf.nist.gov/18.15#iv}.

\bibitem{dlmf-mp}
NIST Digital Library of Mathematical Functions,
Section 18.19, Meixner--Pollaczek definitions and orthogonality,
equations 18.19.6--9.
\url{https://dlmf.nist.gov/18.19}.

\bibitem{dlmf-mp-generating}
NIST Digital Library of Mathematical Functions,
Section 18.23, Meixner--Pollaczek generating function,
equation 18.23.7.
\url{https://dlmf.nist.gov/18.23.E7}.

\end{thebibliography}
\end{document}
